\documentclass[10pt]{article}
% \documentclass[10pt, draft]{article} % draft mode (no figures)

\usepackage[margin=1in]{geometry}
\usepackage{amsmath,amssymb,amsthm,bm,mathtools}
\usepackage{enumitem}
\usepackage{graphicx}
\usepackage{subcaption}
\usepackage{booktabs}
\usepackage{tabularx}
\usepackage{makecell}
\usepackage{float}
\usepackage{xcolor}
\usepackage{hyperref}
\usepackage[style=numeric,sorting=nyt]{biblatex}

\addbibresource{references.bib}
\addbibresource{references_additions.bib}

\hypersetup{hidelinks}

% Let biblatex break long DOIs and URLs in the bibliography.
\setcounter{biburlnumpenalty}{7000}
\setcounter{biburlucpenalty}{7000}
\setcounter{biburllcpenalty}{7000}

\newtheorem{theorem}{Theorem}
\newtheorem{proposition}{Proposition}
\newtheorem{lemma}{Lemma}
\newtheorem{corollary}{Corollary}
\theoremstyle{definition}
\newtheorem{definition}{Definition}
\newtheorem{assumption}{Assumption}
\newtheorem{remark}{Remark}

\newcommand{\R}{\mathbb{R}}
\newcommand{\E}{\mathbb{E}}
\newcommand{\Cov}{\operatorname{Cov}}
\newcommand{\Corr}{\operatorname{Corr}}
\newcommand{\Var}{\operatorname{Var}}
\newcommand{\diag}{\operatorname{diag}}
\newcommand{\1}{\mathbf{1}}
\newcommand{\D}{\mathcal{D}}
\newcommand{\Hcal}{\mathcal{H}}
\newcommand{\Vcal}{\mathcal{V}}
\newcommand{\argminop}{\operatorname*{arg\,min}}
\newcommand{\argmaxop}{\operatorname*{arg\,max}}


% Allows the source to compile outside the project directory.
\newcommand{\safeincludegraphics}[2][]{%
  \IfFileExists{#2}{\includegraphics[#1]{#2}}{%
  \fbox{\parbox[c][0.19\textheight][c]{0.92\linewidth}{\centering
  Figure file not found in this folder.\\[2mm]{\scriptsize\ttfamily\detokenize{#2}}}}}}


\title{Targeting Price-Related Regressivity in Machine-Learning Mass Appraisal:
Covariance-Guided Objectives, Post-Hoc Comparison, and Temporal Limits}

\author{
Nicol\'as Acevedo\thanks{Operations Research Center, MIT, Cambridge, MA 02139. \texttt{nacevedo@mit.edu}}
\qquad
Haihao Lu\thanks{Sloan School of Management, MIT, Cambridge, MA 02139. \texttt{haihao@mit.edu}}
\\[1ex]
Michael Wagner\thanks{Cook County Assessor's Office, Chicago, IL 60602. \texttt{Michael.Wagner2@cookcountyil.gov}}
\qquad
Timothy Sparer\thanks{Cook County Assessor's Office, Chicago, IL 60602. \texttt{Timothy.Sparer@cookcountyil.gov}}
}
\date{}

\begin{document}
\maketitle


\begin{abstract}

Machine-learning mass-appraisal models can be predictively strong yet leave a systematic
price-related trend in valuation ratios: the pattern assessors call \emph{regressivity}. We
study a narrow intervention in a research translation of the Cook County Assessor's Office
residential LightGBM workflow. The target is one global, first-order quantity: the covariance
between log valuation errors and log sale price among observed sales, which is
proportional to the classical log-ratio/log-price slope. A Direct penalty acts on that squared covariance; a Jensen-derived
Surrogate is observation-additive and so fits the standard boosting custom-objective interface.
In a fixed prediction space with an intercept, Direct is exactly a one-parameter centered spread
map and the Surrogate a generally multi-directional weighted projection. We trace prespecified penalty grids with
seven chronological rolling-origin folds, a later held-out block, and a separate 2025
forward period, holding the base learner fixed. A centered-spread post-hoc rescaling
solved on development predictions reproduces much of Direct's first-order movement over the
range Direct attains, so training-time regularization is not necessary to move this trend; yet
at matched first-order correction the families are not interchangeable in residual structure,
and no family is uniformly better across the reported diagnostics. Across those folds, the held-out block, and 2025 the correction-strength
relationship does not transport, so no penalty strength is selected. The contribution is one
auditable first-order mechanism for observed sales with its temporal limits, not a fairness
guarantee or a claim about tax burdens.

\end{abstract}


\section{Introduction}
\label{sec:introduction}

Property taxes are one of the largest sources of local government revenue in the United States. They help finance public schools, roads and streets, and other essential local services \cite{IAAO2020PropertyTaxPolicy}. Using 2016 U.S. Census Bureau data, \textcite{Berry2021} reports that property taxes accounted for roughly 72\% of local tax revenue and 47\% of local own-source general revenue, and totaled about \$500 billion annually.

{Property taxation in the United States is generally based on property values}: local assessors estimate property values, which are then used, together with jurisdiction-specific assessment and taxation rules, to determine each taxpayer's share of the property-tax burden \cite{IAAO2020PropertyTaxPolicy,IAAO2025Mass}. Consequently, systematic overassessment of some properties and underassessment of others can shift tax liability from one group of property owners to another, creating or exacerbating inequities \cite{Berry2021,IAAO2020PropertyTaxPolicy}.

The {valuation} of groups of properties for tax purposes is known as \emph{mass appraisal}. The International Association of Assessing Officers (IAAO) defines mass appraisal as ``the process of valuing a group of properties as of a given date using common data, standardized methods, and statistical testing'' \cite{IAAO2025Mass}. Within this process, assessors use statistical models to estimate property values across large property groups. {When implemented in computer-based valuation systems, such models are commonly referred to as automated valuation models (AVMs)} \cite{IAAOAITaskForce2022,BidansetRakow2022}. Recent mass-appraisal research increasingly examines nonlinear machine-learning (ML) methods, including gradient-boosted trees, alongside traditional regression approaches, given their high predictive power \cite{Rootetal2023,Sevgen2024100111,Zilli2024}.

However, when evaluating the performance of a mass-appraisal valuation model, predictive accuracy alone is insufficient. The IAAO recommends the use of \emph{ratio studies}, which compare {appraised values} with observed sale prices.
{A central quantity in these studies is the ratio of a valuation estimate to a market-value proxy such as sale price. Because the CCAO model studied here predicts \emph{fair market value} before the statutory conversion to assessed value and the subsequent steps that determine taxable value and tax liability, we evaluate its model-estimated market value relative to the observed sale price and refer to this quantity as the \emph{valuation-to-sale ratio}, or \emph{valuation ratio} for brevity \cite{CCAOModelResAVM2026,CCAOAssessmentNotice2026,CCAOPropertyTaxSystem2026}.} Ratio studies use such ratios to evaluate valuation level and uniformity, including whether similar properties are valued consistently (\emph{horizontal uniformity}) and whether valuation ratios vary systematically across property-value levels (\emph{vertical equity}) \cite{IAAO2013Ratio,IAAO2025Mass,IAAO2026ExposureRatio}.


{A particularly relevant form of vertical inequity is assessment \emph{regressivity}: valuation ratios tend to be higher for properties with lower underlying market values and lower for properties with higher market values. Because market value is latent, empirical vertical-equity diagnostics necessarily rely on observed sales or other market-value proxies \cite{IAAO2026ExposureRatio}. This pattern can place a disproportionate share of the assessment, and potentially tax, burden on owners of lower-valued properties \cite{McMillen2020,Berry2021,AvenancioLeonHoward2022}. Moreover, regressivity may remain hidden in aggregate measures of predictive accuracy because these measures do not identify whether valuation errors vary systematically across the value distribution. A valuation model can therefore achieve strong predictive accuracy while still exhibiting a systematic value-related ratio pattern.}


{This motivates asking whether regressivity that remains after fitting a strong predictive model can be reduced through an explicit training or calibration criterion rather than treated only as an ex-post diagnostic. In mass appraisal, a useful correction must operate within existing sales-based modeling, validation, and review processes: current prices are observed only for properties that transact, while valuation systems must produce values within fixed assessment schedules and remain compatible with professional and administrative review \cite{IAAO2020SalesVerification,IAAO2018AVM,IAAO2025Mass}. Accordingly, the correction should reduce regressivity while preserving predictive performance and requiring minimal changes to existing valuation workflows.}

This paper arises from a collaboration between the Cook County Assessor's Office (CCAO) and academic researchers, motivated by an operational goal: reducing regressivity in the CCAO's residential valuation workflow. The CCAO currently uses a county-wide LightGBM (Light Gradient-Boosting Machine) model \cite{lightgbm} as its primary residential valuation model, {following earlier township-level linear-regression and gradient-boosting workflows} \cite{CCAOModelResAVM2026}. The current model is trained on screened residential sales and estimates a \emph{fair market value} using property characteristics, location, environmental variables, and market trends. It operates within a broader assessment pipeline that separates data ingestion, model training, assessment, evaluation, interpretation, and finalization \cite{CCAOModelResAVM2026}.


{The model's fair-market-value estimate is only one input into the final property-tax bill. In Cook County, that value is converted to assessed value according to property classification, equalized using the state equalization factor, reduced by eligible exemptions, and ultimately combined with local tax rates determined from taxing-district levies and taxable property values \cite{CCAOAssessmentNotice2026,CCAOPropertyTaxSystem2026}. We therefore focus on the valuation stage rather than the downstream steps that convert estimated market values into final tax liabilities.} Everything reported below is a Python research translation of the office's published residential workflow, evaluated on its verified sale sample against observed sale prices as the transaction-based market-value proxy. It is not an official production run, an assessment decision, or a deployment recommendation.


{Within the valuation stage, the central problem addressed in this paper is the interaction between predictive performance and regressivity. In the translated CCAO workflow the current model class is highly accurate and still price-regressive: on the held-out evaluation it attains an $R^2_P$ of 0.894 on the price scale while its exposure-draft vertical-equity index is -26.5\% and its coefficient of dispersion is 21.6\%, above the range the adopted ratio-study guidance gives for single-family residential property (Section~\ref{subsec:baseline_tension}). Predictive accuracy alone therefore does not remove the ratio trend. This makes the design problem a within-model one: we seek to retain the predictive behaviour of the current model class while reducing the regressive value-related pattern, rather than to choose between model classes. From an operational perspective, the correction should fit within CCAO's existing modeling and validation workflow without requiring a fundamental redesign of the surrounding valuation pipeline \cite{CCAOModelResAVM2026}.}


{We address this challenge by modifying the valuation model's training objective, leaving the data, feature set, prediction target, and surrounding validation and evaluation workflow unchanged. The quantity we train on is not a new statistic. It is the classical first-order relation the assessment literature has long used to diagnose price-related inequity---the slope of log valuation ratios on log sale prices, whose direct antecedent in this setting is \textcite{Cheng1974}---which on the log-price scale is proportional to the covariance between log valuation errors and log sale prices and is therefore something a learner can be penalized for. We study two mechanisms built on it. The \emph{Direct} penalty adds that squared covariance to the objective. The \emph{Surrogate} penalty is a Jensen upper bound on it that turns out to be observation-additive, and is exactly a weighted squared-error loss whose weights grow with squared log-price distance from the training center. The two are not implementation variants of one objective: Direct applies a single sample-coupled, rank-one correction, whereas the Surrogate reshapes the data-fit loss itself and can combine several learnable modes, so their paths need not coincide. Neither the classical log-ratio/log-price target nor covariance regularization nor dependence-based fair regression originates here; what this paper supplies is the translation of that classical target into a trainable objective, the exact fixed-prediction-space distinction between the two mechanisms, and the evaluation of what they attain. For a fixed penalty formulation the correction introduces a single additional tuning parameter, the penalty strength.
The framework is intentionally targeted: it does not, by itself, determine overall valuation level, eliminate nonlinear or local inequity, or replace the assessor-facing vertical-equity diagnostics used for evaluation.}


{The empirical strategy traces paths rather than searching for a preferred point on one. Both objectives are run over prespecified penalty grids with the base learner held fixed, so what varies along a path is the objective and not the model's capacity. Development evaluation uses rolling-origin chronological folds; a later held-out block is the primary out-of-time evaluation of the development fit, and a separate 2025 forward sample is evaluated after refitting on all eligible 2016--2024 sales. Beside the retrained configurations we evaluate the simple post-hoc alternatives that act on already-fitted predictions, and we compare the families at \emph{matched first-order correction} rather than at matched penalty, which is what makes it possible to ask whether moving the targeted slope by the same amount leaves the same residual structure behind. Price-scale comparisons are read under a stated retransformation convention, because one ordering in the matched-correction comparison reverses under the alternative one. We report where the paths behave well and where they fail, including how far development-period path structure does not carry over to later market periods, and we select no penalty strength anywhere.}


{The paper makes three contributions, all of them application-driven. \emph{Objective translation and mechanism.} It carries the classical assessor-side first-order price-related pattern
into a trainable objective inside a research translation of CCAO's residential LightGBM workflow, and settles the relationship between the two resulting penalties in the fixed prediction space: the Direct correction is a one-parameter centered spread map and the Surrogate a generally multi-directional weighted projection. \emph{Mechanism comparison rather than winner selection.} It compares the training-time corrections against a serious centered-spread post-hoc comparator, and compares the three families at matched first-order correction, which shows that moving the first-order slope does not determine the complete residual geometry. \emph{Temporal failure-mode audit.} It traces the chronological, held-out, and forward paths of both objectives and documents the limits of raw-$\rho$ portability, rather than selecting or recommending a deployment penalty.}

The remainder of this section reviews related work on machine learning in mass appraisal, vertical equity, and dependence-based fair regression. Section~\ref{sec:setting} then defines the CCAO setting and evaluation framework, Section~\ref{sec:method} develops the covariance-guided correction, and the remaining sections describe the empirical design, report the regularization-path evidence, and discuss practical implications, limitations, and extensions.


\subsection{Related Work}
\label{subsec:related_work}

\paragraph{Machine Learning in Mass Appraisal}
Recent reviews document the {growing research on} automated valuation models and machine learning (ML) in mass appraisal and {real estate} valuation. {Multiple regression remains common, alongside tree ensembles such as random forests and gradient boosting, as well as neural-network-based methods} \cite{WangLi2019,Rootetal2023,Zilli2024,purdueAssessingTechniques}. These reviews organize the work they survey primarily by predictive criteria, such as out-of-sample accuracy and performance relative to alternative algorithms. Vertical equity, although routinely evaluated in assessor practice, is not among the criteria around which they organize model training or model selection \cite{WangLi2019,Rootetal2023,Zilli2024}. Tree ensembles are among the methods these reviews both survey and record as predictively competitive \cite{Rootetal2023,Sevgen2024100111,Zilli2024}. A practical training-time equity intervention should therefore be compatible with the model classes these reviews document rather than require a specialized valuation architecture.

 {Operational adoption by assessment offices nevertheless remains partial and uneven.} \textcite{BidansetRakow2022} identify {communication challenges, limited technical capacity, and budget constraints among perceived barriers to wider adoption}. {Moreover, predictive modeling is only one component of the assessment workflow; the use of ML does not replace sales verification, out-of-time evaluation, ratio studies, model interpretation, or administrative review} \cite{IAAO2020SalesVerification,IAAOAITaskForce2022}. {The} \textcite{IAAOAITaskForce2022} emphasizes explainability, defensibility, administrative feasibility, and integration into existing workflows as central conditions for operational use of advanced ML models. {These considerations favor equity-aware methods that can be incorporated into existing model-training, validation, and review pipelines with minimal disruption.}

\paragraph{Vertical equity and regressivity.}
Equity has long been a central objective in mass appraisal practice and standards. The assessor literature distinguishes \emph{horizontal equity}, which concerns the uniform treatment of similar properties, from \emph{vertical equity}, which concerns systematic differences in assessment outcomes across property-value levels \cite{IAAO2013Ratio,cornia2006property,McMillenSingh2023}. A common form of vertical inequity is \emph{regressivity}: assessment-to-sale ratios tend to decline with property value, so lower-valued properties tend to be relatively overassessed and higher-valued properties tend to be relatively underassessed. {Ratio studies evaluate assessment uniformity using complementary diagnostics. For example, the Coefficient of Dispersion (COD) summarizes dispersion in assessment ratios, whereas the Price-Related Differential (PRD) and Price-Related Bias (PRB) are commonly used to diagnose price-related vertical inequity} \cite{IAAO2013Ratio,IAAO2023VerticalReview}. The statistical measurement of price-related vertical inequity has a long history in this literature. \textcite{PaglinFogarty1972} reframed property-tax equity around the distribution of assessment-to-value ratios; \textcite{Cheng1974} asked how assessment performance itself should be measured and, in doing so, supplied the direct antecedent of the diagnostic this paper trains on: he writes a log assessment-to-sale ratio as a function that includes the log settlement price, entering at coefficient $\theta-1$, so a log-ratio/log-price slope of exactly this form is already present in the classical assessment literature. Cheng also discusses how appraisal lag and heterogeneous value growth can affect the observed ratio--price relationship. \textcite{Edelstein1979} appraised residential property-tax regressivity directly; and \textcite{SundermanEtAl1990} developed statistical tests for vertical inequity. That work supplies both the ratio-based framing used here and the recurring lesson that alternative specifications of the same test need not agree.


{No single diagnostic fully characterizes vertical equity. Methodological studies show that vertical-equity measures can differ materially in their statistical behavior, and that some regression-based diagnostics can be biased when assessed values are themselves generated by regression-based valuation models} \cite{IAAO2023VerticalReview,McMillenSingh2023}. {Accordingly, the assessor literature recommends examining multiple complementary diagnostics rather than relying on a single global measure.} The classical ratio--price slope is therefore one member of a complementary set that also contains the price-related differential, price-related bias, the modified Kakwani index and the proposed exposure-draft vertical-equity index; this paper trains on the slope and continues to report the others, rather than treating any one of them as a complete account of vertical equity.  {Empirical studies further document substantial regressivity and consequential distributional disparities across jurisdictions} \cite{Berry2021,AvenancioLeonHoward2022}. Cook County is also the setting for direct evidence on algorithmic assessment reform: \textcite{Liu2026} exploits the staggered triennial reassessment cycle around the office's adoption of a county-wide machine-learning valuation model and reports a substantial flattening of the log assessed-to-sale-price gradient afterwards. That is external evidence that reforming an assessment algorithm can materially change measured price regressivity; it evaluates none of the objectives studied here and supplies no benchmark for them. {Recent large-scale evidence also shows that predictive accuracy and assessment equity need not be intrinsically opposed: improvements in the information available to valuation models can, in many settings, improve both simultaneously \cite{SmithEtAl2026}.}

{A smaller literature moves beyond aggregate diagnosis toward model specification, localized validation, and explicit correction intended to improve assessment equity.} Proposed approaches include spatial and geographically weighted methods, stratified procedures, ratio-study-based model validation, and fairness-aware valuation models \cite{Quintos2014,BidansetEtAl2019,BidansetEtAl2025,HermansEtAl2022,HermansEtAl2025,CandoganHanLu2023}. These studies approach assessment equity along several distinct dimensions: the modeling architecture, the information set, the spatial or stratification structure, the validation design, or an explicit equity or fairness formulation written into the estimation problem. The route taken here is narrower and different rather than better: one first-order training target inside an existing learner, with no change to architecture, features, or sample construction. No experiment in this paper combines it with any of those alternatives, so no claim of compatibility with them is made.
{A closely related approach is \textcite{CandoganHanLu2023}, who develop a \emph{K-segment} property-valuation model that explicitly targets regressivity and studies the associated accuracy--fairness relationship. Our approach addresses the same assessor-side objective through a different mechanism: rather than introducing a segmented valuation architecture, we use a first-order log-error and log-sale-price association as a training target within an existing regression learner. The method therefore targets a deliberately narrow first-order mechanism during training while retaining standard ratio studies, proxy-aware vertical-equity diagnostics, and subgroup review for out-of-sample evaluation.}

\paragraph{Fair regression and dependence-based constrained learning.}


Although much of the fairness-in-ML literature initially focused on classification, \emph{fair regression} studies how application-specific fairness requirements can be incorporated when predictions are continuous, as in property valuation. Existing approaches consider criteria such as statistical parity, bounded group loss, mean parity, or distributional parity and implement them through constrained or regularized empirical risk minimization, post-processing, or recalibration \cite{AgarwalDudikWu2019,FitzsimonsEtAl2019,ChzhenEtAl2020Wasserstein,ChzhenEtAl2020Recalibration,WeiEtAl2023}. The broader methodological principle is that predictive risk can be optimized subject to, or jointly with, an application-specific fairness criterion. Varying the strength of this requirement produces a family of solutions that makes the relationship between predictive performance and the targeted fairness criterion explicit, without requiring that improvements in one necessarily worsen the other.

{A particularly relevant strand expresses fairness requirements as restrictions on statistical dependence between predictions and a sensitive variable.} {Prior work has quantified or controlled such dependence using correlation and coefficients of determination \cite{KomiyamaEtAl2018}, kernel-based measures such as the Hilbert--Schmidt Independence Criterion (HSIC) \cite{PerezSuayEtAl2017,LiEtAl2022Fairness}, and maximal-correlation criteria \cite{MaryEtAl2019,lee2022maximal}. Related regularization approaches also control the contribution of sensitive attributes to regression predictions \cite{ScutariPaneroProissl2022}.}
{These approaches are particularly relevant when the variable of concern is continuous and can often be incorporated directly into regression objectives. Our application, however, differs from conventional group-fairness settings. Sale price is the observed transaction-based proxy for market value and the prediction target, not a protected attribute, and our objective is not demographic parity. We adapt the broader dependence-control principle to a narrower structural target: the signed first-order association between log valuation errors and log sale prices. The assessor literature supplies the domain-specific objective---reducing value-related regressivity---while the fair-regression literature supplies a general mechanism for incorporating statistical associations into model training.}

A global slope correction is also not the only way to address regressivity: segmented valuation, richer spatial and temporal information, and post-hoc affine recalibration \cite{BeheshtiEtAl2019} are close alternatives, and the last of these is close enough that it is carried through this paper as an explicit comparator rather than left as a possibility. The narrower methodological contribution is to use that classical first-order relationship as a trainable target in the CCAO workflow, derive and distinguish the Direct and observation-additive Surrogate objectives, and evaluate their temporal paths and failure modes. Because a centered post-hoc spread map is the exact fixed-space form of the Direct correction, the comparison against it is reported here rather than deferred, and no claim that in-processing is preferable is made on the strength of it.

\paragraph{Residual--outcome decorrelation outside mass appraisal.}
The mechanism controlled here, namely dependence between a model's residuals and the
very outcome that model predicts, has already been formulated, and controlled, outside
mass appraisal. Some of the closest methodological prior art comes from outside the
valuation literature, particularly brain-age prediction, where the same structure appears
as the ``brain-age delta'' problem: the residual of predicted age on true age is
systematically associated with true age.
\textcite{SmithEtAl2019} set out that estimation problem together with its
regression-based correction, and \textcite{BeheshtiEtAl2019} give a post-hoc affine
bias-adjustment scheme for it. \textcite{TrederEtAl2021} impose a correlation constraint
directly on the training objective and trace a tunable family of solutions as the
constraint strength varies. \textcite{WangEtAl2023} instead reshape the loss so that it
corrects the target-dependent component of predictive bias, and \textcite{RenEtAl2026}
place an explicit delta--outcome correlation loss inside a nonlinear learner.
\textcite{LeeChen2026} supply a statistical framing for this class of procedures, through
outcome-calibrated regression and inference based on predicted outcomes. Related
fair-regression work also expresses dependence restrictions through measures such as the
coefficient of determination \cite{KomiyamaEtAl2018}.

\paragraph{The novelty boundary.}
It is worth stating in one place what this paper does and does not claim, since
each of the four lineages above already contains part of the construction. We claim no
novelty for the log-ratio/log-price vertical-equity diagnostic, which is classical in
assessment practice and appears in the form used here in \textcite{Cheng1974}; for
covariance or correlation regularization as a training criterion
\cite{KomiyamaEtAl2018,TrederEtAl2021}; for fairness-aware empirical-risk minimization in
regression generally; for residual--outcome decorrelation \cite{SmithEtAl2019}; for a
tunable decorrelation path indexed by a penalty strength \cite{TrederEtAl2021}; for
target-dependent corrective losses \cite{WangEtAl2023}; for nonlinear in-processing
correction \cite{RenEtAl2026}; or for affine post-hoc orthogonalization of residuals
against the outcome \cite{BeheshtiEtAl2019}, whose statistical framing is likewise
established \cite{LeeChen2026}. What we do claim is narrower and
specific to this application: the CCAO-centered translation of the classical first-order
assessment relation into a LightGBM-compatible training objective, expressed in the
office's own diagnostics; the exact fixed-prediction-space distinction between the Direct
mechanism and the Jensen-derived, observation-additive Surrogate, together with an accurate
account of the curvature structure the custom-objective interface admits for each; the comparison against
post-hoc correction at matched first-order targets rather than at matched penalty; and the
chronological, held-out and forward audit of the resulting paths, reported with its failure
modes rather than summarized as a pass. No penalty strength is selected anywhere in this
paper.


\section{Preliminaries: The CCAO Baseline and Empirical Motivation}
\label{sec:setting}


This section fixes the CCAO application setting, prediction notation, and evaluation framework used throughout the paper. We first describe the residential valuation workflow and the baseline model. We then distinguish predictive performance, assessor-facing measures of valuation level and ratio dispersion, complementary vertical-equity diagnostics, and mechanism and residual-structure diagnostics for first-order, nonlinear conditional-mean, and broader dependence patterns. Finally, we report the baseline level that motivates the proposed correction.


\subsection{The CCAO Residential Valuation Workflow}
\label{subsec:ccao_workflow}

The CCAO values a large and heterogeneous residential inventory on a triennial reassessment cycle \cite{CCAOAssessmentCalendar2026}. Its residential modeling evolved from township-level linear-regression models, with linear-regression and gradient-boosting alternatives used in 2019--2020, to a county-wide LightGBM model beginning in 2021 \cite{CCAOModelResAVM2026}. LightGBM is an implementation of gradient-boosted decision trees \cite{lightgbm}. That transition has itself been studied as an algorithmic assessment reform: using the staggered triennial reassessment cycle and the log assessed-to-sale-price ratio, \textcite{Liu2026} reports a substantial flattening of the price gradient afterwards.

The CCAO's current published residential automated valuation model (AVM) uses verified sales and a broad set of property, spatial, neighborhood, market, and contextual predictors to estimate residential market values \cite{CCAOModelResAVM2026}. The prediction model is embedded in a broader operational workflow comprising data preparation, model training, valuation, evaluation, interpretation, and final review. Candidate valuations are evaluated using both predictive-performance measures and assessor-specific ratio-study statistics across multiple geographic areas and property classes.


The empirical analysis in this paper is a within-model one. The earlier township-level linear-regression workflow is historical context for the diagnostics used here and is not re-estimated: the reported baseline is unpenalized LightGBM, representing the current model class. The prediction target, underlying feature information, development sample, temporal design, and evaluation procedures are held fixed across every configuration compared.
Every comparison reported in this paper is therefore a within-model contrast under one fixed empirical design; it is neither a comparison of model classes nor a reconstruction of official historical assessments. The exact sample construction, transaction filters, and temporal validation procedures are described in Section~\ref{subsec:ccao_design}.


\subsection{Prediction Setting and the Baseline Model}
\label{subsec:baseline_models}

Let $\mathcal{I}$ denote the set of residential sales retained after the CCAO sales-validation and transaction-eligibility procedures described in Section~\ref{subsec:ccao_design}. For each $i\in\mathcal{I}$, we observe
\[
(x_i,P_i,t_i),
\]
where $x_i$ contains property and parcel characteristics; tract-level socioeconomic and demographic variables from the U.S. Census Bureau's American Community Survey (ACS); and geographic, neighborhood, environmental, accessibility, market, and temporal features.


{The variable $P_i>0$ is the verified sale price and is treated as the observed transaction-based proxy for the property's market value, while $t_i$ denotes the sale date used to construct the chronological samples \cite{CCAOModelResAVM2026,IAAO2026ExposureRatio}. Any time-derived predictors used by the models are included in $x_i$. Market value itself is latent: a verified sale provides an observed market-value indicator rather than a directly observed market value. The prediction problem is to estimate market value from the information available in $x_i$.}

Since residential sale prices are positive and highly right-skewed, the baseline model is trained on the log-price scale. Define
\begin{equation}
\label{eq:prediction_scale}
y_i=\log P_i,
\qquad
\widehat{y}_i=f(x_i),
\qquad
e_i(f)=\widehat{y}_i-y_i,
\end{equation}
{where $f\in\mathcal{H}$ is a candidate prediction function and $e_i(f)$ is the \emph{log valuation error}, defined with prediction minus observed log price so that it coincides with the log valuation ratio below.}
The corresponding prediction on the original price scale is
\begin{equation}
\label{eq:price_prediction}
\widehat{P}_i=\exp\bigl(f(x_i)\bigr).
\end{equation}


{The log-price scale is the estimation scale, whereas $\widehat{P}_i$ is the directly retransformed valuation used for price-scale predictive evaluation and ratio-study diagnostics. Throughout the analysis, $\widehat{P}_i$ is defined operationally by~\eqref{eq:price_prediction}; direct exponentiation is not presented as a bias correction for the conditional mean on the original price scale. We define the corresponding \emph{valuation-to-sale ratio} as}
\begin{equation}
\label{eq:valuation_ratio}
r_i(f)
=
\frac{\widehat{P}_i}{P_i}
=
\exp\bigl(e_i(f)\bigr).
\end{equation}


{For brevity, we subsequently refer to $r_i(f)$ as the \emph{valuation ratio}. Its numerator is the model-estimated market value, not Cook County's statutory assessed value after downstream adjustments. When a fitted model is fixed, we suppress the dependence on $f$ and write simply $r_i$ and $e_i$. The log valuation error and valuation ratio are linked exactly by}
\begin{equation}
\label{eq:log_ratio}
e_i(f)
=
\log \widehat{P}_i-\log P_i
=
\log r_i(f).
\end{equation}
Thus, $e_i(f)>0$ (equivalently, $r_i(f)>1$) corresponds to a prediction above the observed sale price, whereas $e_i(f)<0$ (equivalently, $r_i(f)<1$) corresponds to a prediction below it.
{Estimation on the log-price scale therefore converts multiplicative price errors into additive log valuation errors, and $e_i(f)$ is exactly the log valuation ratio.}

The predictive baseline is a single model class, gradient-boosted decision trees, representing CCAO's current model class. As described above, the earlier township-level linear-regression workflow is historical context for the diagnostics used here and is not estimated in this paper. Across every configuration compared, we hold fixed the prediction target, development observations, underlying feature information, temporal design, and evaluation protocol; configurations differ only in the training objective.

For a training sample of size $n$, the ordinary baseline is fitted to reduce mean squared log-price error, i.e.,
\begin{equation}
\label{eq:baseline_learning}
L(f)
=
\frac{1}{n}\sum_{i=1}^n e_i(f)^2.
\end{equation}
This loss is the predictive component of the training objective modified in Section~\ref{sec:method}.


{LightGBM, the modeling foundation of CCAO's current published residential AVM, constructs an additive ensemble of regression trees \cite{CCAOModelResAVM2026,lightgbm}. A gradient-boosted tree predictor can be written as}
\begin{equation}
\label{eq:boosted_tree_model}
f_{\mathrm{gbm}}^{(M)}(x)
=
f^{(0)}(x)
+
\sum_{m=1}^{M}\nu b_m(x),
\end{equation}
where $f^{(0)}$ is an initial prediction, $b_m$ is the regression tree added at boosting iteration $m$, $M$ is the number of boosting iterations, and $\nu>0$ is the learning rate. The tree ensemble can represent nonlinearities, threshold effects, and interactions without requiring them to be specified in advance.

LightGBM fits the ensemble in~\eqref{eq:boosted_tree_model} sequentially. The gradient and Hessian calculations required for the proposed custom objectives are introduced with the penalties in Section~\ref{sec:method}.

{The ordinary baseline objective minimizes squared log-price error and contains no term that directly controls systematic value-related patterns in log valuation errors or valuation ratios. We next distinguish predictive performance from the assessor-facing criteria used to evaluate valuation level, ratio dispersion, and regressivity.}


\subsection{
\texorpdfstring{
Predictive, Assessor-Facing, and Residual-Structure Diagnostics}{Predictive, Assessor-Facing, and Residual-Structure Diagnostics}}
\label{subsec:metrics}


{The two scales play distinct roles in evaluation. Predictive performance is assessed both on the log-price scale used for estimation and on the original price scale on which valuations are reported. Assessor-facing ratio-study diagnostics are computed on the original price scale. The resulting measures answer different questions and are therefore interpreted jointly rather than collapsed into a single performance criterion.}

\paragraph{Predictive performance.}
{For an evaluation sample $S\subseteq\mathcal{I}$, let $n_S=|S|$ and $\bar P_S=\frac{1}{n_S}\sum_{i\in S}P_i$. We use four complementary headline accuracy-related measures: price-scale mean absolute error (MAE), price-scale mean absolute percentage error (MAPE), price-scale $R^2$, and log-scale root mean squared error (RMSE):}
\begin{align}
\label{eq:predictive_metrics}
\operatorname{MAE}_{P}(S)
&=
\frac{1}{n_S}\sum_{i\in S}\left|\widehat P_i-P_i\right|,
\\
{\operatorname{MAPE}_{P}(S)}
&{=100\times
\frac{1}{n_S}
\sum_{i\in S}
\left|
\frac{\widehat P_i-P_i}{P_i}
\right|
=100\times
\frac{1}{n_S}\sum_{i\in S}|r_i-1|},
\\
R^2_{P}(S)
&=
1-
\frac{\sum_{i\in S}(\widehat P_i-P_i)^2}
{\sum_{i\in S}(P_i-\bar P_S)^2},
\\
\operatorname{RMSE}_{\log P}(S)
&=
\left(
\frac{1}{n_S}\sum_{i\in S}e_i(f)^2
\right)^{1/2}.
\end{align}


{The price-scale MAE gives the average absolute prediction error in dollars, while MAPE gives the average absolute error relative to the observed sale price. In particular, MAPE measures absolute deviation from the unit-ratio benchmark $r_i=1$ and is therefore distinct from COD below, which measures dispersion around the sample median ratio. The statistic $R^2_P$ summarizes predictive fit on the original price scale. Finally, $\operatorname{RMSE}_{\log P}$ directly evaluates the squared-error scale used for estimation and, by~\eqref{eq:log_ratio}, is equivalently the root mean squared log valuation ratio.}


\paragraph{Valuation level.}
Recall from~\eqref{eq:valuation_ratio} that $r_i=\widehat P_i/P_i$. A ratio above one indicates that the predicted value exceeds the observed sale price, whereas a ratio below one indicates that it falls below the sale price. Define the median, arithmetic mean, and price-weighted mean ratios as
\begin{equation}
\label{eq:level_ratios}
m_S
=
\operatorname{median}_{i\in S}\{r_i\},
\qquad
\bar r_S
=
\frac{1}{n_S}\sum_{i\in S}r_i,
\qquad
\bar r_{W,S}
=
\frac{\sum_{i\in S}\widehat P_i}{\sum_{i\in S}P_i}
=
\frac{\sum_{i\in S}r_iP_i}{\sum_{i\in S}P_i}.
\end{equation}


{These quantities summarize aggregate valuation level. Because both numerator and denominator of $r_i$ are expressed on the market-value scale in this analysis, the neutral reference value for all three statistics is one. The median is relatively insensitive to extreme ratios, the arithmetic mean assigns equal weight to each sale, and the weighted mean places greater weight on higher-value properties. Agreement among these measures indicates broad aggregate calibration but does not establish either horizontal uniformity or vertical equity.}


\paragraph{Ratio dispersion (Horizontal uniformity).}
Ratio dispersion is an important assessor-facing dimension of horizontal uniformity. We summarize it using the coefficient of dispersion (COD) and coefficient of variation (COV). The COD is
\begin{equation}
\label{eq:cod}
\operatorname{COD}(S)
=
100\times
\frac{1}{n_S}
\sum_{i\in S}
\frac{|r_i-m_S|}{m_S},
\end{equation}
and the COV is
\begin{equation}
\label{eq:cov}
\operatorname{COV}(S)
=
100\times
\frac{s_{r,S}}{\bar r_S},
\qquad
s_{r,S}
=
\left(
\frac{1}{n_S-1}
\sum_{i\in S}(r_i-\bar r_S)^2
\right)^{1/2}.
\end{equation}


{The COD measures the average absolute deviation from the median ratio, whereas the COV measures standard deviation relative to the mean. Lower values indicate less aggregate ratio dispersion, but the reference ranges depend on property class and prevailing market conditions, and unusually low dispersion can itself warrant diagnostic review \cite{IAAO2013Ratio,IAAO2026ExposureRatio}.}
To avoid ambiguity with the covariance penalty introduced later, $\operatorname{COV}$ always denotes the coefficient of variation, whereas statistical covariance is written as $\operatorname{Cov}(\cdot,\cdot)$.

\paragraph{Regressivity (Vertical equity).}
{Vertical equity concerns whether valuation ratios vary systematically across underlying market value. A \emph{regressive} pattern arises when lower-value properties tend to have higher valuation ratios than higher-value properties; a \emph{progressive} pattern arises in the opposite direction. Because true market value is latent, empirical vertical-equity procedures necessarily use observed sales or other market-value proxies, and the diagnostics below differ partly in how they address this measurement problem \cite{IAAO2026ExposureRatio}. No single statistic fully characterizes vertical equity: different diagnostics summarize different aspects of the relationship between valuations, ratios, and market-value proxies and can therefore disagree \cite{IAAO2023VerticalReview,IAAO2026ExposureRatio}. We consequently use a complementary suite of ratio-based, regression-based, distributional, and grouped diagnostics rather than treating any single statistic as definitive. The May 2026 IAAO exposure draft places primary emphasis on the grouped VEI and identifies PRB, MKI, and PRD as complementary measures. Because the document remains an exposure draft, we use its proposed procedures and ranges as descriptive guidance rather than as final adopted compliance standards.}

The \emph{price-related differential} (PRD) is
\begin{equation}
\label{eq:prd}
\operatorname{PRD}(S)
=
\frac{\bar r_S}{\bar r_{W,S}}.
\end{equation}
Note that, since the weighted mean ratio assigns greater weight to higher-value properties, a PRD above one tends to arise when higher-value properties have lower valuation ratios than lower-value properties. The 2013 IAAO standard provides a reference range of $[0.98,1.03]$ for the PRD: values above this range tend to indicate regressivity, whereas values below it tend to indicate progressivity \cite{IAAO2013Ratio}.

{PRD is a simple and widely used ratio-based diagnostic, but it can be particularly sensitive to influential high-value observations and extreme ratios. Therefore, we interpret it jointly with the other vertical-equity measures rather than as a standalone test \cite{IAAO2026ExposureRatio}.}


{The following regression- and group-based vertical-equity diagnostics require a proxy for the property's value level. Ranking properties only by sale price can mechanically bias ratio-based vertical-equity tests toward finding regressivity because sale price also appears in the denominator of $r_i$, while using valuation alone can induce bias in the opposite direction. Following IAAO guidance, we therefore use the equal-weight market-value proxy}
\begin{equation}
\label{eq:market_value_proxy}
V_i^{\mathrm{proxy}}
=
\frac{1}{2}
\left(
P_i+\frac{\widehat P_i}{m_S}
\right),
\end{equation}
{which combines the verified sale price with the median-normalized predicted value to reduce this source of proxy-induced bias \cite{IAAO2013Ratio,IAAO2026ExposureRatio}.}

The 2013 IAAO Standard explicitly defines the equal-weight proxy as $0.50(\mathrm{AV}/\mathrm{Median})+0.50\,\mathrm{SP}$. As of August 2026, the May 2026 ratio-studies document remains an exposure draft and the 2013 Standard remains the recommended IAAO guidance. The exposure draft's VEI appendix describes an equal-weight proxy but displays the normalized valuation term without an explicit $0.50$ multiplier. For reproducibility, all PRB and VEI results in this paper use the single equal-weight convention in Eq.~\eqref{eq:market_value_proxy}, consistent with the 2013 Standard and the exposure draft's stated equal-weight interpretation.
The executed canonical metric code uses this equal-weight proxy exactly. PRB in \texttt{utils/motivation\_utils.py} computes $0.5\,\mathrm{SP}+0.5\,(\mathrm{AV}/\mathrm{Med})$, and VEI computes $0.5\,y_{\mathrm{true}}+0.5\,(y_{\mathrm{pred}}/m_S)$. Both match Eq.~\eqref{eq:market_value_proxy}. SHA256 of that file at population time is \texttt{4219a9d68e5266e300b319660f93d37ca184189ee340edb0785f40c9f3c39beb}.

{We first report the \textit{price-related bias} (PRB).}
Let $\widetilde V_S^{\mathrm{proxy}}$ denote the median of $V_i^{\mathrm{proxy}}$ over $S$. The PRB is the slope in the simple linear regression
\begin{equation}
\label{eq:prb}
\frac{r_i-m_S}{m_S}
=
a_S
+
\operatorname{PRB}(S)\cdot
\log_2\left(
\frac{V_i^{\mathrm{proxy}}}
{\widetilde V_S^{\mathrm{proxy}}}
\right)
+
\varepsilon_i.
\end{equation}

Negative PRB values indicate that valuation ratios tend to decline as the market-value proxy increases and are therefore associated with regressivity. Values between $-0.05$ and $0.05$ are conventionally preferred \cite{IAAO2013Ratio}.

{As a complementary distributional diagnostic, we report the \emph{Modified Kakwani Index} (MKI). Let $\pi(1),\allowbreak\ldots,\pi(n_S)$ index the observations in $S$ in nondecreasing order of sale price, and let
\[
q_j=\frac{j-\tfrac12}{n_S},
\qquad
\overline{\widehat P}_S
=
\frac{1}{n_S}\sum_{i\in S}\widehat P_i.
\]
Define the \emph{concentration index} of predicted values when observations are ordered by sale price and the corresponding \emph{Gini coefficient} of sale prices as
\begin{equation}
\label{eq:mki_components}
C_{\widehat P\mid P}(S)
=
\frac{2}{n_S\overline{\widehat P}_S}
\sum_{j=1}^{n_S}
\widehat P_{\pi(j)}q_j
-1,
\qquad
G_P(S)
=
\frac{2}{n_S\bar P_S}
\sum_{j=1}^{n_S}
P_{\pi(j)}q_j
-1,
\end{equation}
respectively. Then, the MKI is
\begin{equation}
\label{eq:mki}
\operatorname{MKI}(S)
=
\frac{C_{\widehat P\mid P}(S)}{G_P(S)}.
\end{equation}
When there is sufficient price variation, an MKI value of one indicates vertical neutrality, whereas values below one indicate a regressive direction, and values above one indicate a progressive direction \cite{IAAO2026ExposureRatio}. Because MKI orders observations using sale price, it can itself be affected by errors in the market-value proxy, which is another reason to interpret it as one component of the broader diagnostic suite rather than as a definitive test.}


Finally, we report the \emph{Vertical Equity Indicator} (VEI), the quantile-grouped vertical-equity measure emphasized in the May 2026 IAAO exposure draft \cite{IAAO2026ExposureRatio}.
The draft varies the number of percentile groups with sample size: for the large evaluation samples used in the main countywide analysis ($n_S\geq 501$), observations are divided into deciles according to $V_i^{\mathrm{proxy}}$. Let $S_{q_1}$ and $S_{q_{10}}$ denote the lowest and highest deciles, and let $m_{S_{q_1}}$ and $m_{S_{q_{10}}}$ denote their median valuation ratios. The VEI is
\begin{equation}
\label{eq:vei}
\operatorname{VEI}(S)
=
100\times
\frac{m_{S_{q_{10}}}-m_{S_{q_1}}}{m_S}.
\end{equation}

A negative VEI indicates that higher-value properties have lower median valuation ratios than lower-value properties. The exposure draft proposes a reference band of $[-10\%,10\%]$, together with confidence-interval-based testing for point estimates outside that band.

{Again, since the May 2026 document remains an exposure draft, we use VEI and the proposed band as descriptive guidance rather than as a final adopted compliance standard. Moreover, note that the first-versus-last quantile comparison in~\eqref{eq:vei} does not necessarily reveal non-monotone patterns within the median ratios distribution. Hence, we also report the median valuation ratio and its 90\% confidence interval for every VEI percentile group, together with the overall median ratio.}


\begin{table}[!ht]
\centering
\footnotesize
\setlength{\tabcolsep}{4pt}
\renewcommand{\arraystretch}{1.10}

% Indentation for individual measures within each family
\newcommand{\metriclabel}[1]{\hspace*{0.9em}#1}

\caption{Predictive and assessor-facing evaluation measures.}
\label{tab:assessment_metrics_summary}

\begin{tabularx}{\textwidth}{
@{}
>{\raggedright\arraybackslash}p{3.8cm}
>{\centering\arraybackslash}p{1.8cm}
>{\centering\arraybackslash}p{2.9cm}
>{\centering\arraybackslash}p{2.3cm}
>{\raggedright\arraybackslash}X
@{}
}
\toprule
\textbf{Measure}
& \textbf{\shortstack{Ideal value}}
&
\textbf{\shortstack{Reference range}}
& \textbf{Guidance status}
& \textbf{Interpretation} \\
\midrule

\textbf{Predictive performance}
& & & & \\
\cmidrule(r){1-1}

\metriclabel{$R^2_P$}
& $1$
& --
& n/a
& Original price-scale goodness of fit. \\

\metriclabel{$\operatorname{MAE}_P$}
& $0$
& --
& n/a
& Mean absolute prediction error in dollars. \\

\metriclabel{$\operatorname{MAPE}_P$}
& $0$
& --
& n/a
& Mean absolute prediction error relative to sale price. \\

\metriclabel{$\operatorname{RMSE}_{\log P}$}
& $0$
& --
& n/a
& RMSE on the log-price estimation scale. \\


\midrule
\addlinespace[5pt]
\textbf{{Valuation level}
}
& & & & \\
\cmidrule(r){1-1}

\metriclabel{Median ratio $m_S$}
& $1$
& $[0.90,1.10]$
& IAAO 2013 (adopted)
& Robust aggregate valuation level. \\

\metriclabel{Mean ratio $\bar r_S$}
& $1$
& $[0.90,1.10]$
& IAAO 2013 (adopted)
& Unweighted aggregate valuation level. \\

\metriclabel{Weighted mean $\bar r_{W,S}$}
& $1$
& $[0.90,1.10]$
& IAAO 2013 (adopted)
& Value-weighted valuation level. \\


\midrule
\addlinespace[6pt]
\textbf{{Horizontal uniformity}
}
& & & & \\
\cmidrule(r){1-1}

\metriclabel{COD}
& --
& $[5.0,15.0]$
& IAAO 2013 (adopted)
& Absolute dispersion relative to the median ratio. \\

\metriclabel{COV}
& --
& $[6.25\%,18.75\%]$
& derived approximation
& Relative dispersion around the mean ratio. \\


\midrule
\addlinespace[5pt]
\textbf{Vertical equity}
& & & & \\
\cmidrule(r){1-1}

\metriclabel{PRD}
& $1$
& $[0.98,1.03]$
& IAAO 2013 (adopted)
& Mean-to-weighted-mean ratio comparison. \\

\metriclabel{PRB}
& $0$
& $[-0.05,0.05]$
& IAAO 2013 (adopted)
& Regression-based value-related ratio trend. \\

\metriclabel{MKI}
& $1$
& $[0.95,1.05]$
& May 2026 (exposure draft)
& Distributional concentration relative to sale prices. \\

\metriclabel{VEI}
& $0$
& $[-10\%,10\%]$
& May 2026 (exposure draft)
& High-versus-low value-group median-ratio difference. \\

\bottomrule
\end{tabularx}

\vspace{1mm}
\begin{minipage}{\textwidth}
\scriptsize
\emph{Notes.}
``--'' denotes that no universal ideal or guidance range is imposed. The ranges below combine adopted IAAO guidance, explicitly identified derived approximations, and supplementary exposure-draft or implementation guidance as detailed next.
The median-ratio, mean-ratio, weighted-mean-ratio, PRD, PRB, and COD ranges
follow IAAO ratio-study and mass-appraisal guidance
\cite{IAAO2013Ratio,IAAO2018AVM,IAAO2025Mass}.
For COV, the approximate range $[6.25\%,18.75\%]$ is derived from the
IAAO relationship $\operatorname{COV}\approx1.25\,\operatorname{COD}$ under
approximately normal ratio distributions and the residential-improved COD
reference range $[5,15]$.
The MKI range $[0.95,1.05]$ is used by CCAO's ratio-study implementation
and is also identified in recent IAAO exposure-draft guidance as an
optimal range for lower-variability strata.
The May 2026 IAAO exposure draft proposes $[-10\%,10\%]$ for VEI.
The exposure draft remains under review; the 2013 Standard remains the
recommended IAAO guidance until a final revision is approved.
The guidance-status column records that standing explicitly, because
the two are not interchangeable: the level-ratio, COD, PRD and PRB ranges come
from the \emph{adopted} 2013 Standard, whereas the MKI and VEI bands come from
the May 2026 \emph{exposure draft} and are proposed rather than adopted
guidance. The COV range is neither: it is a derived approximation obtained from
the COD range through the stated relationship, and is labelled as such rather
than presented as a published band. Predictive measures carry no guidance status
because no ratio-study standard defines a range for them.
\end{minipage}

\end{table}


{Table~\ref{tab:assessment_metrics_summary} summarizes the principal predictive and assessor-facing criteria used throughout the paper. Predictive metrics quantify different aspects of valuation error, valuation-level statistics summarize aggregate calibration, COD and COV summarize ratio dispersion, and VEI, PRB, MKI, and PRD provide complementary grouped, regression-based, distributional, and ratio-based views of vertical equity. Strong predictive performance or low aggregate dispersion does not by itself imply neutral vertical-equity behavior.}
The mechanism and residual-structure diagnostics below are deliberately excluded from the reference-range table because they are not assessor-standard compliance measures.


\paragraph{
Mechanism and residual-structure diagnostics.}
{The assessor-facing measures above evaluate valuation outcomes. The following three
diagnostics instead examine the residual--price structure most directly relevant to
the proposed training mechanism. They form an increasingly general sequence:
$\beta_{\log}$ measures the targeted first-order trend, $\Delta_{\mathrm{NL}}$
measures non-affine conditional-mean structure beyond that trend, and distance
correlation measures broader statistical dependence. None is treated as an
assessor-standard compliance statistic.

\emph{First-order mechanism.}
We first report the slope from regressing log valuation ratios on log sale prices:
\begin{equation}
\label{eq:log_ratio_slope}
\beta_{\log}(S)
=
\frac{
\widehat{\operatorname{Cov}}_S(e,y)
}{
\widehat{\operatorname{Var}}_S(y)
}
=
\frac{
\widehat{\operatorname{Cov}}_S(\log r,\log P)
}{
\widehat{\operatorname{Var}}_S(\log P)
}.
\end{equation}
A value of zero indicates no first-order log-ratio--log-sale-price trend. Negative
values indicate the regressive direction and positive values the progressive
direction. For a fixed evaluation sample, $\widehat{\operatorname{Var}}_S(y)$ is
common across models, so $\beta_{\log}(S)$ is the signed first-order slope
corresponding directly to the covariance targeted during training.
This slope is not introduced here: it has a classical
property-assessment antecedent in \textcite{Cheng1974}, who already writes a log
assessment-to-sale ratio as a function that includes the log settlement price,
entering at coefficient $\theta-1$. What follows from this
paper's own log-error notation is the exact covariance identity used for training
in~\eqref{eq:training_covariance}, not the ratio--price relationship it summarizes.
Because sale
price is an imperfect proxy for latent market value and also appears in the
denominator of $r_i$, we use $\beta_{\log}$ as a mechanism diagnostic rather than
as a standalone definition of vertical equity.

\emph{Nonlinear conditional-mean structure.}
Let
\[
Z=\frac{y-\E[y]}{\sqrt{\Var(y)}}.
\]
Following the nonparametric coefficient-of-determination literature, define Pearson's
correlation ratio
\[
\eta^2(e\mid Z)
=
\frac{\Var\!\left(\E[e\mid Z]\right)}{\Var(e)},
\]
which measures the fraction of log-error variance explained by the unrestricted
conditional mean of $e$ given $Z$ \cite{DoksumSamarov1995}. The squared Pearson
correlation $\Corr(e,Z)^2$ is the corresponding fraction explained by the best
affine relationship. We therefore define the \emph{correlation-ratio nonlinearity
gap}
\begin{equation}
\label{eq:nonlinearity_gap}
\Delta_{\mathrm{NL}}
:=
\eta^2(e\mid Z)-\Corr(e,Z)^2.
\end{equation}
The descriptive name and the symbol $\Delta_{\mathrm{NL}}$ are paper-specific;
the underlying correlation-ratio construction is established in the statistical
literature. The gap is nonnegative at the population level and equals zero when the
conditional mean of the log valuation error is affine in log sale price. Positive
values therefore measure conditional-mean explanatory power beyond the best affine
trend. This distinction is especially relevant along a regularization path:
$\lvert\beta_{\log}\rvert$ may continue to fall even if stronger regularization
induces an S-shaped, U-shaped, or otherwise non-affine residual--price pattern.
Appendix~\ref{app:metrics} gives the equivalent projection representation, relates
the construction to parametric-versus-nonparametric regression-fit comparisons
\cite{HardleMammen1993}, and specifies the fixed cross-fitted estimator used in the
held-out and 2025 evaluations.
The same fixed estimator was later applied, fold by fold, to the already-frozen chronological-validation predictions as a retrospective diagnostic; those reconstructed CV values do not define the $\rho$ grid or any model-selection rule.

\emph{Broader dependence.}
Finally, we report distance correlation,
\begin{equation}
\label{eq:dcor_diagnostic}
\operatorname{dCor}(e,y)
=
\operatorname{dCor}(\log r,\log P),
\end{equation}
as an omnibus diagnostic of dependence between log valuation errors and log sale
prices \cite{SzekelyRizzoBakirov2007}. At the population level and under the usual
moment conditions, distance correlation is zero if and only if the two variables
are independent. It can therefore respond to nonlinear mean relationships,
price-related dispersion, and other forms of dependence that are intentionally
outside the narrower scope of $\beta_{\log}$ and $\Delta_{\mathrm{NL}}$. Distance
correlation does not identify a regressive or progressive direction, and
distributional independence is stronger than vertical neutrality as conventionally
evaluated in ratio studies. We use it as a broader residual-structure diagnostic,
not as an assessor-facing measure of vertical equity.
}


{Three benchmark principles govern the empirical use of these measures. First, point estimates are treated as descriptive model diagnostics rather than as formal determinations of compliance with IAAO standards; uncertainty is reported where it is material to substantive comparisons, including the confidence intervals accompanying the VEI percentile-group profile \cite{IAAO2013Ratio,IAAO2026ExposureRatio}. Second, countywide statistics can conceal geographic and property-stratum heterogeneity; such subgroup analyses remain necessary complements for later operational validation when sample size and value variation support meaningful comparison. Third, during model development, evaluation statistics are computed on chronological validation folds before fold-level results are aggregated to summarize each regularization path. The full suite is used for evaluation rather than imposed jointly as a collection of model constraints.}


\subsection{Baseline Level and Empirical Motivation}
\label{subsec:baseline_tension}

This subsection establishes the \emph{level} of the assessor-facing diagnostics attained by unpenalized LightGBM, CCAO's current model class. Township-level linear regression historically preceded that workflow \cite{CCAOModelResAVM2026}, and the classical vertical-equity testing literature \cite{PaglinFogarty1972,Cheng1974,Edelstein1979,SundermanEtAl1990} developed on models of that kind, which is why the diagnostics below are the ones an assessor would recognize. The earlier model class is not re-estimated in this paper, and no quantitative linear-versus-nonlinear comparison is reported. Every configuration reported below is estimated and evaluated under the same prediction target,
underlying feature information, sample-construction rules, temporal-validation
framework, and evaluation procedure. What the baseline table reports is therefore a
level for the current model class rather than a reconstruction of official
assessments produced under different historical workflows.
Section~\ref{subsec:ccao_design} provides the complete empirical design.

{The primary evaluation below uses the held-out evaluation sample defined in Section~\ref{subsec:ccao_design}. We also report a separate 2025 forward evaluation after the baseline specification is refit using all eligible 2016--2024 sales; the 2025 sample is not used for model selection.}

% Tier B2.2: tab:ccao_baseline_results is re-sourced from the frozen full zero-penalty
% control, analysis/p0_major_revision_validation/tables/zero_control_full.csv, cell_id=A.
% The former provenance line pointed at output/paper_v6_preselection, which is not frozen
% evidence and is not read by this pass.

For readability and to keep the empirical hierarchy explicit, the main text retains the primary baseline table for prediction and assessor-facing vertical equity, while the complementary valuation-level, horizontal-uniformity, and mechanism/residual-structure baseline diagnostics are reported in Appendix~\ref{app:ccao_path_details}.

{\begin{table}[!ht]
\centering
\scriptsize
\setlength{\tabcolsep}{2.4pt}
\renewcommand{\arraystretch}{1.06}
\newcommand{\baselineprimarymetric}[1]{\hspace*{0.6em}#1}
\caption{Primary baseline level: prediction and vertical equity for \emph{Ordinary LightGBM (standard raw-label native)}, the assessor-facing workflow benchmark, on the held-out and 2025 forward evaluations. This is a level statement about one model class, not a comparison between model classes and not a penalty effect.}
\label{tab:ccao_baseline_results}
\begin{tabularx}{\textwidth}{@{} >{\raggedright\arraybackslash}p{6.4cm} >{\centering\arraybackslash}X >{\centering\arraybackslash}X @{}}
\toprule
\textbf{Measure} & \textbf{Held-out evaluation} & \textbf{2025 forward evaluation} \\
\midrule
\multicolumn{3}{@{}l}{\textbf{Prediction}} \\
\cmidrule(r){1-1}
\baselineprimarymetric{$R^2_P$} & 0.894 & 0.904 \\
\baselineprimarymetric{$\operatorname{MAE}_P$} & \$75,655 & \$78,484 \\
\baselineprimarymetric{$\operatorname{MAPE}_P$} & 21.2\% & 20.8\% \\
\baselineprimarymetric{$\operatorname{RMSE}_{\log P}$} & 0.289 & 0.278 \\
\addlinespace[5pt]
\multicolumn{3}{@{}l}{\textbf{Vertical equity}} \\
\cmidrule(r){1-1}
\baselineprimarymetric{PRD} & 1.069 & 1.079 \\
\baselineprimarymetric{PRB} & -0.091 & -0.106 \\
\baselineprimarymetric{MKI} & 0.923 & 0.907 \\
\baselineprimarymetric{VEI} & -26.5\% & -28.6\% \\
\bottomrule
\end{tabularx}
\vspace{1mm}
\begin{minipage}{\textwidth}
\scriptsize
\emph{Notes.}
Every entry is cell A of the frozen full zero-penalty control, the assessor-facing workflow benchmark named \emph{Ordinary LightGBM (standard raw-label native)} in Section~\ref{subsec:comparators}. The linear-regression columns this table previously carried have been removed rather than requalified: no frozen artifact reproduces any of those values, and one of them printed the same $R^2_P$ to three decimals on two different evaluation samples. A qualitative historical note on the earlier township-level workflow survives in the text above; no linear-regression number is reported anywhere in this paper.
There is no second column to compare against, so no entry is boldfaced and no entry denotes statistical significance, standards compliance, model selection, or a preferred penalty strength.
Preferred values and applicable reference ranges, and the guidance status of each measure, are defined in Table~\ref{tab:assessment_metrics_summary}: PRD and PRB belong to the adopted Standard on Ratio Studies, whereas MKI and VEI are Exposure-Draft diagnostics and therefore proposed rather than adopted.
MAPE and VEI are reported in percent; MAE is reported in dollars.
For the held-out evaluation the model is fit on the 344,607-sale development pool; for the 2025 forward evaluation it is refit on all 382,897 eligible 2016--2024 sales.
\end{minipage}
\end{table}

}

The level is itself the empirical motivation. On the primary held-out evaluation the
workflow benchmark attains an $R^2_P$ of 0.894 on the price scale while its price-related bias is
-0.091, its exposure-draft vertical-equity index is -26.5\%, and its log-ratio/log-price slope is
-0.150. Its coefficient of dispersion is 21.6\% --- above the $[5,15]$ range the adopted Standard on
Ratio Studies gives for single-family residential property (Table~\ref{tab:assessment_metrics_summary}). On the 2025 forward sample the entire 95\% confidence
interval for the price-related bias runs from -0.111 to -0.101 and therefore lies below $-0.10$. High predictive
accuracy and a substantial price-related ratio trend therefore coexist in the current model class,
which is all the motivation the rest of the paper needs; it does not depend on a comparison with any
other model class.


The corresponding complementary baseline diagnostics are retained in Appendix Table~\ref{tab:ccao_baseline_complementary}.


\begin{figure}[!htbp]
\centering
\safeincludegraphics[width=0.8\textwidth]{img/generated_v12_994/baseline_models_motivation_2024_2025.pdf}
\caption{Descriptive valuation-ratio profiles against sale price in the primary held-out evaluation and the 2025 forward evaluation. Curves show median ratios in 30 equal-count sale-price bins; shaded regions show the corresponding interquartile ranges. The panel retains the historical township-level linear-regression trace as visual context only: no linear-regression quantitative result is reported in this paper, none should be read off this figure, and the assessor-facing level for the workflow benchmark is the statement accompanying Table~\ref{tab:ccao_baseline_results}.}
\label{fig:baseline_motivation}
\end{figure}


The corresponding VEI percentile-group profile is reported in Appendix~\ref{app:ccao_path_details} (Figure~\ref{fig:vei_group_profile_placeholder}); the intervals drawn there are descriptive bootstrap intervals on the plotted group medians and carry no exposure-draft compliance verdict.
Table~\ref{tab:ccao_baseline_results} and Figure~\ref{fig:baseline_motivation} present the primary baseline level in the main text. Appendix Table~\ref{tab:ccao_baseline_complementary} reports the complementary valuation-level, horizontal-uniformity, and mechanism/residual-structure diagnostics for the same cell.


Nothing in this level statement should be read as evidence that machine-learning or nonlinear valuation models are intrinsically more regressive than linear ones. No linear-regression result is reported here, the frozen evidence for this application reproduces none, and the finding is application-specific in any case. What the level does establish is that high predictive accuracy and a substantial price-related ratio trend coexist in the current model class. The design question is therefore a within-model one --- whether that predictive behaviour can be retained while the value-related pattern is reduced through the training objective --- rather than a choice between model classes.


{The assessor-facing diagnostics above determine how the resulting valuations are evaluated; they need not be optimized directly during training. The first-order mechanism slope in Eq.~\eqref{eq:log_ratio_slope} isolates the signed log-ratio--log-sale-price trend:
\[
\beta_{\log}(S)
=
\frac{\widehat{\operatorname{Cov}}_S(e,y)}
{\widehat{\operatorname{Var}}_S(y)}.
\]
Within a fixed training sample, $\widehat{\operatorname{Var}}(y)$ does not depend on the fitted model. Penalizing squared covariance therefore controls the same first-order slope up to a fixed scaling. Using $e_i(f)=\log r_i(f)$ and $y_i=\log P_i$,
\begin{equation}
\label{eq:training_covariance_bridge}
\widehat{\operatorname{Cov}}\bigl(e(f),y\bigr)
=
\widehat{\operatorname{Cov}}\bigl(\log r(f),\log P\bigr).
\end{equation}
We use covariance as the trainable mechanism, while the assessor-facing measures and the residual-structure diagnostics remain separate out-of-sample checks on what the correction does and does not remove. Section~\ref{sec:method} develops the correction.}


\section{Covariance-Guided Regressivity Correction}
\label{sec:method}


{This section translates the value-related regressivity identified in Section~\ref{sec:setting} into a training-time correction penalty. We first define a global first-order covariance-based regressivity target on the log-price scale, then introduce a direct squared-covariance penalty, and finally derive a sample-additive upper-bound surrogate that is compatible with standard boosting software. The correction penalty leaves the prediction target, feature information, and base model class unchanged. That is, for a fixed penalty formulation, the only additional tuning parameter is the penalty strength, which we denote $\rho$.}

{Let $J_0(f)$ denote the baseline training objective, whose data-fit component is the squared log-price loss $L(f)$ defined in Section~\ref{subsec:baseline_models}. All quantities below are computed from the ``current'' training sample. In temporal validation, the centering quantities and penalties are recomputed using the training block of each fold only.}


\subsection{{Log-Scale Regressivity Target}}
\label{subsec:training_regressivity}

Let
\begin{equation}
\label{eq:centered_log_price}
\bar y=\frac{1}{n}\sum_{i=1}^n y_i,
\qquad
c_i=y_i-\bar y,
\qquad
\bar e(f)=\frac{1}{n}\sum_{i=1}^n e_i(f).
\end{equation}
The empirical covariance between log-residuals and log-prices is
\begin{equation}
\label{eq:training_covariance}
C(f)
=
\widehat{\operatorname{Cov}}\bigl(e(f),y\bigr)
=
\frac{1}{n}\sum_{i=1}^n
\bigl(e_i(f)-\bar e(f)\bigr)(y_i-\bar y)
=
\frac{1}{n}\sum_{i=1}^n e_i(f)c_i.
\end{equation}
The last equality follows from $\sum_{i=1}^n c_i=0$, so centering the residuals does not change the covariance.

This covariance has a direct slope interpretation. Let
\[
V_y=\frac{1}{n}\sum_{i=1}^n c_i^2.
\]
Provided $V_y>0$, the least-squares slope from regressing log valuation ratios on log prices is
\begin{equation}
\label{eq:covariance_slope}
\operatorname{slope}\bigl(e(f)\sim y\bigr)
=
\frac{C(f)}{V_y}.
\end{equation}
Because $V_y$ is fixed by the training sample, reducing $|C(f)|$ is equivalent to reducing the magnitude of this global linear trend. In particular, $C(f)<0$ indicates that log valuation ratios tend to decline with log sale price, which is the first-order pattern associated with regressivity.


A qualification is essential because sale price is both the prediction target and the observed value benchmark. Let $m(X)=\E[Y\mid X]$ be the population squared-error predictor and let $e^\star=m(X)-Y$. Then $\E[e^\star]=0$ and
\begin{equation}
\label{eq:bayes_residual_covariance}
\Cov(e^\star,Y)=-\E\!\left[\Var(Y\mid X)\right]\leq0.
\end{equation}
Thus, even the Bayes predictor has negative residual--outcome covariance whenever sale prices contain conditionally unpredictable variation. Appendix~\ref{app:theory} gives the finite-sample analogue: in a fixed linear prediction space containing an intercept, ordinary least squares satisfies $C_0=-\lVert f_0-y\rVert_n^2\leq0$. Driving training covariance toward zero can therefore spread predictions toward realized transaction prices rather than recover latent market values. Negative covariance is not, by itself, evidence of latent vertical inequity, and zero covariance is not a fairness condition. We use $C(f)$ only as a controllable first-order mechanism and judge the resulting valuations with chronological out-of-sample, proxy-aware ratio and residual-structure diagnostics.

{This slope is a training-scale mechanism measure, not the assessor-facing PRB statistic defined in Section~\ref{subsec:metrics}. We use covariance rather than correlation because $V_y$ is fixed within the training sample, whereas correlation additionally normalizes by the model-dependent residual variance. Covariance therefore provides a direct training target for the global value-related slope.}

{Sale price enters this construction as the observed market-value benchmark for verified training sales; it is not introduced as a protected attribute. The broader connection to dependence-based fair regression is discussed in Section~\ref{subsec:related_work}.}


\subsection{{Direct Squared-Covariance Penalty}}
\label{subsec:direct_penalty}

{We first penalize the chosen first-order target directly. The \emph{covariance-penalized learner} is}


{
\begin{equation}
\label{eq:direct_penalty}
f_\rho^{\mathrm{cov}}
\in
\argminop_{f\in\Hcal}
\left\{
J_0(f)+\frac{\rho}{2}C(f)^2
\right\},
\qquad
\rho\geq 0.
\end{equation}
}

When $\rho=0$, formulation~\eqref{eq:direct_penalty} reduces to the baseline learner. As $\rho$ increases, the optimization places greater weight on reducing the global residual--price trend. Squaring the covariance treats negative and positive trends symmetrically. This is important because the objective is not to replace regressivity with progressivity, but to move the first-order relationship toward zero.


{For interpretation, it is enough to consider the contribution of the covariance penalty to the gradient of observation $i$:}
{
\begin{equation}
\label{eq:direct_penalty_gradient}
\frac{\partial}{\partial \widehat y_i}
\left[
\frac{\rho}{2}C(f)^2
\right]
=
\frac{\rho}{n}C(f)c_i.
\end{equation}
}


{Suppose, for example, that the fitted model is regressive due to $C(f)<0$. For a sale below the geometric-mean price, $c_i<0$, and thus the covariance component of the gradient is positive, i.e., a descent step lowers its predicted value. For a sale above the geometric-mean price, $c_i>0$, and thus the covariance component is negative, i.e., a descent step raises its predicted value. The penalty therefore acts in the direction required to flatten the global log-ratio--log-price trend. Note that the direction and interpretation reverse automatically when the fitted trend is progressive.}


Note that the direct penalty targets how valuation ratios vary with sale price, but not their
overall valuation level. Multiplying every predicted value by the same positive constant shifts all log valuation ratios by the same amount but leaves the covariance $C(f)$ unchanged. Thus, even a model with $C(f)=0$ can have a median or mean valuation ratio different from one. Hence, valuation level must be evaluated separately.


{For second-order models (e.g., LightGBM), the corresponding Hessian of the squared log-price loss plus direct penalty is}
{
\begin{equation}
\label{eq:direct_gradient_hessian}
\nabla_{\widehat y}^2
\left(
L(f)+\frac{\rho}{2}C(f)^2
\right)
=
\frac{2}{n}I+\frac{\rho}{n^2}cc^\top.
\end{equation}
}
The rank-one Hessian term $\frac{\rho}{n^{2}}cc^\top$ introduces cross-observation curvature. Standard LightGBM custom objectives provide one gradient and one Hessian value per observation, so this dense term cannot be represented exactly through the usual interface. An exact second-order implementation therefore requires additional machinery: using only the diagonal is possible, but discards the cross-observation curvature. The next subsection proposes an alternative penalty term to handle this caveat. Further, the model-specific derivations are given in Appendix~\ref{app:implementation}.


\subsection{{Sample-Additive Covariance Upper-Bound Surrogate}}
\label{subsec:surrogate}


{For standard second-order boosting software, we seek a penalty that decomposes exactly across observations. Given that~\eqref{eq:training_covariance} already gives $C(f)=\frac{1}{n}\sum_{i=1}^n e_i(f)c_i$, we can make use of Jensen's inequality that yields a sample-additive upper bound to the squared-covariance term directly.}


{
\begin{equation}
\label{eq:surrogate_upper_bound}
C(f)^2
=
\left(
\frac{1}{n}\sum_{i=1}^n e_i(f)c_i
\right)^2
\leq
\frac{1}{n}\sum_{i=1}^n e_i(f)^2c_i^2.
\end{equation}
}
{This inequality motivates the fully observation-additive surrogate penalty}
\begin{equation}
\label{eq:surrogate}
\Psi^{\mathrm{surr}}(f)
=
\frac{1}{n}\sum_{i=1}^n e_i(f)^2c_i^2
=
\frac{1}{n}\sum_{i=1}^n
e_i(f)^2(y_i-\bar y)^2.
\end{equation}
The resulting \emph{surrogate-penalized learner} is therefore
\begin{equation}
\label{eq:surrogate_learning}
f_\rho^{\mathrm{surr}}
\in
\argminop_{f\in\Hcal}
\left\{
{J_0(f)}
+
\rho\Psi^{\mathrm{surr}}(f)
\right\}.
\end{equation}
{
Thus, $\Psi^{\mathrm{surr}}(f)$ is a genuine upper bound on the squared empirical covariance in~\eqref{eq:training_covariance}. Its exact relationship with the direct-covariance penalty can be seen by making the observation
\begin{equation}
\label{eq:surrogate_decomposition}
\Psi^{\mathrm{surr}}(f)
=
C(f)^2
+
\frac{1}{n}
\sum_{i=1}^n
\left(e_i(f)c_i-C(f)\right)^2.
\end{equation}
Hence, $\Psi^{\mathrm{surr}}(f)$ exceeds $C(f)^2$ by the nonnegative dispersion term in~\eqref{eq:surrogate_decomposition}: opposite-signed observation-level contributions can cancel in $C(f)$ but not in the surrogate. This upper-bound relation should not be interpreted as control of a broader class of nonlinear residual--price relationships; neither the bound nor its weighted-loss representation directly penalizes nonlinear conditional-mean lack-of-fit.
}


Combining the baseline loss with the surrogate gives
\begin{equation}
\label{eq:surrogate_weighted_loss}
L(f)+\rho\Psi^{\mathrm{surr}}(f)
=
\frac{1}{n}\sum_{i=1}^n
\left[1+\rho(y_i-\bar y)^2\right]e_i(f)^2.
\end{equation}
Thus, the surrogate is equivalent to a weighted squared-error objective with fixed observation weights
\begin{equation}
\label{eq:surrogate_weights}
w_i(\rho)=1+\rho(y_i-\bar y)^2.
\end{equation}
Properties near the center of the log-price distribution receive weights close to the baseline weight, whereas properties in the lower and upper tails receive larger weights.


{Since $\exp(\bar y)$ is the geometric mean sale price, the weight function is symmetric in log-distance from this center: prices that differ from the geometric mean by reciprocal multiplicative factors receive the same additional weight. This symmetry concerns the weighting function only: the empirical lower and upper tails, and the resulting prediction changes, need not be symmetric. Moreover, because the additional weight grows quadratically with log-price distance, observations far from the center can receive disproportionate influence. The log transformation compresses this leverage relative to an analogous raw-price construction, but does not bound it.}


{This representation clarifies the mechanism without imposing a directional assumption on individual observations. The surrogate raises the cost of residuals farther from the center of the log-price distribution. Unlike the direct penalty, its observation-level correction is not determined by the sign of the aggregate covariance: it depends on the residual $e_i(f)$ and the squared log-price deviation $c_i^2$. Its resulting effect on the signed covariance and on assessor-facing vertical-equity diagnostics must therefore be evaluated out of sample.}


For squared loss, the observation-wise gradient and Hessian are
\begin{align}
\label{eq:surrogate_derivatives}
g_i
&=
\frac{2}{n}e_i(f)
\left[1+\rho(y_i-\bar y)^2\right],
\\
h_i
&=
\frac{2}{n}
\left[1+\rho(y_i-\bar y)^2\right].
\end{align}
Both quantities depend only on observation $i$ once $\bar y$ has been computed. They can therefore be supplied through a standard second-order boosting custom-objective interface. Equation~\eqref{eq:surrogate_weighted_loss} also shows that the same mathematical objective can, in principle, be expressed as weighted squared loss in regression software that supports fixed observation weights. The CCAO paths reported in this paper use the custom-objective implementation below; an alternative native-weighted implementation should not be assumed to reproduce the identical finite boosting path without an explicit implementation-equivalence check.


{For the LightGBM implementation, we multiply the complete objective by the common positive factor $n/2$. This leaves the minimizer and the relative weighting indexed by $\rho$ unchanged while putting the $\rho=0$ derivatives on LightGBM's native squared-error scale. The supplied surrogate derivatives are therefore
\[
g_i=e_i(f)\left[1+\rho c_i^2\right],
\qquad
h_i=1+\rho c_i^2.
\]
The corresponding direct-objective implementation is given in Appendix~\ref{app:implementation}.}

The Jensen step is what makes the objective observation-additive, and it is not
free. Because the Surrogate multiplies each observation's squared error by
$1+\rho c_i^{2}$, the implied weights concentrate on sales far from the centre of the
log-price distribution as $\rho$ grows, and the frozen weight audit quantifies how far
that goes. At $\rho=100$ the effective sample size
$\bigl(\sum_i w_i\bigr)^{2}/\sum_i w_i^{2}$ on the development block falls to 86,609 of
its 344,607 rows, or 25.1\% of $n$; across the nine fitting blocks it ranges from 25.0\%
to 28.6\% of $n$, the median weight on the development block is 19.1 rather than near
one, and the most extreme one percent of rows carry 10.7\% of the total weight. Weight
concentration of this magnitude is an implementation-side diagnostic, and a plausible
explanation for the Surrogate's behaviour at strong penalties. It is an association, not
an established mechanism: nothing here demonstrates that the concentration causes the
observed path behaviour, and no such claim is made.


{Unlike the direct covariance penalty, the surrogate is not invariant to a common shift in all log predictions because it penalizes the magnitude of $e_i(f)$ itself. It can therefore affect
overall valuation level in addition to the value-related trend.
}


\subsection{{Interpretation, Validation, and Scope}}
\label{subsec:correction_scope}


The Direct objective's curvature was audited directly, and the audit does not
license the convenient reading. Gate E.4 of the frozen validation returns the verdict
INDETERMINATE, and neither of its two acceptance criteria is met: the maximum curvature
ratio at the display anchors is 40.04 against an acceptance threshold of 1.2 for
$\rho\ge1$, and the maximum penalty contribution $M$ over all fitting blocks and anchors
is 0.0118 against an acceptance threshold of 0.01. We therefore do not claim that the
Direct objective behaves as though it were gradient-only, and we do not claim the
opposite either. The frozen evidence supports neither reading, and reporting it as
indeterminate is the accurate statement.

What can be stated exactly is the structure. The Hessian of the Direct
objective is $\frac{2}{n}I+\frac{\rho}{n^{2}}cc^{\top}$, from
\eqref{eq:direct_gradient_hessian}: a diagonal part plus a \emph{dense rank-one} term
that couples every pair of observations. LightGBM's custom-objective interface accepts
one gradient and one scalar Hessian per observation, so the cross-observation rank-one
term is not representable through it and is omitted rather than approximated. What the
implementation supplies in its place is the diagonal of the exactly scaled Hessian, which
the frozen objective-scaling audit records identically for all nine fitting blocks:
$1+\frac{\rho}{2n}c_i^{2}$ for Direct and $1+\rho c_i^{2}$ for the Surrogate, with the
common factor $n/2$ applied analytically so that the supplied derivatives reduce to
$g_i=e_i$ and $h_i=1$ at $\rho=0$. The Direct optimizer therefore combines the exact
penalty gradient with an approximate, diagonal curvature. How much that approximation
matters is precisely what E.4 was designed to settle, and it did not settle it.

{
The two formulations intervene differently. The Direct penalty responds to the aggregate signed first-order association between log valuation ratios and log sale prices. The Surrogate instead changes the data-fit objective itself by upweighting squared errors farther from the center of the log-price distribution. Although this observation-additive objective upper-bounds squared covariance, it is not merely an implementation approximation to the Direct penalty: its additional dispersion term in~\eqref{eq:surrogate_decomposition} gives it a distinct mechanism and can produce different path behavior. Because the two penalties have different scales, the numerical values of $\rho$ are not directly comparable across formulations. Their appropriate scale can also vary across training samples and jurisdictions, so the two regularization paths are traced separately through chronological validation.}


{Neither formulation directly optimizes PRD, PRB, VEI, COD, or valuation level. Nor does zero training covariance imply the absence of nonlinear, geographic, property-class, temporal, or other local value-related patterns. The training covariance is therefore treated as a mechanism target rather than as a sufficient measure of vertical equity. Predictive performance, valuation level, PRD, PRB, VEI, ratio curves, the correlation-ratio nonlinearity gap $\Delta_{\mathrm{NL}}$ in~\eqref{eq:nonlinearity_gap}, and the broader distance-correlation diagnostic remain separate validation and out-of-sample evaluation criteria; subgroup diagnostics remain necessary complements for later operational validation.}

{Sale price is required to construct the correction only for verified sales used during model fitting. Once a model has been trained, the resulting valuation function continues to generate predictions from $x$ alone, as in the baseline AVM; the correction does not require a current sale price for an unsold property.}

For comparisons across training samples, the raw numerical value of $\rho$ is not intrinsically portable. Under a rescaling $y\mapsto ay$, squared error scales as $a^2$ while $C(f)^2$ scales as $a^4$; the scale-invariant Direct coordinate is therefore
\begin{equation}
\label{eq:normalized_rho}
\widetilde\rho=\rho\,\widehat{\Var}_{\mathrm{train}}(y).
\end{equation}
Raw $\rho$ remains the coordinate of the CCAO implementation reported below. Any cross-jurisdiction comparison must recompute the training-block variance without leakage and report $\widetilde\rho$; even then, normalization is a comparability device, not a universal operating rule.

Section~\ref{sec:empirical_design} specifies the chronological validation and regularization-path procedure used in the empirical application. Appendix~\ref{app:theory} gives parallel fixed-prediction-space benchmarks for the two penalties. For the Direct penalty, increasing $\rho$ shrinks the targeted covariance smoothly and, when the space contains an intercept, the exact path is a one-parameter centered spread rescaling of the unpenalized fit. For the Surrogate, the exact path is instead a log-price-distance-weighted projection of the outcome onto the same prediction space; locally, it moves in the learnable component of the baseline residuals weighted by squared log-price distance from the sample center. In both cases, the increase in baseline in-sample squared error is second order near $\rho=0$. These results clarify the rank-one geometry of the Direct correction and the generally multi-directional geometry of the Surrogate, while remaining interpretive benchmarks rather than guarantees for fully retrained boosted trees, whose splits, leaves, and stopping behavior can change with the objective.


\section{Application and Evaluation Design}
\label{sec:empirical_design}

{This section specifies the empirical design used to evaluate the proposed correction. We first describe the CCAO application, data, and implementation; then define the chronological validation and regularization-path protocol; and finally summarize the model comparators. All model comparisons reported in the next section use the evaluation criteria defined in Section~\ref{subsec:metrics}.}


\subsection{{The CCAO Residential Pipeline: Application, Data, and Implementation}}
\label{subsec:ccao_design}


{For the primary application, we implement in Python the components of CCAO's residential valuation workflow required for model development and evaluation. The implementation follows the data, feature, model-fitting, and ratio-study structure of CCAO's published residential pipeline, but it is a research translation rather than a direct production run. The reported results should therefore be interpreted as evidence within this translated workflow, not as official CCAO assessment outputs.}


{Within the LightGBM comparison, the intervention is deliberately restricted to the training objective. We hold fixed the sale filters, log-price target, underlying predictor information, temporal design, transformation back to the price scale, and evaluation procedures. The
fixed Section-2 LightGBM hyperparameter configuration is held fixed across the standard and penalized fits, so $\rho$ is the only dimension varied within each penalized family.
Holding the surrounding model specification fixed makes the training-objective change as controlled as possible. The native-versus-custom $\rho=0$ comparison is reported in Appendix~\ref{app:implementation}; because the completed parity audit identifies this difference as an execution-path distinction, attribution of positive-penalty changes is made within each custom-objective family.}


The primary CCAO application uses verified Cook County residential sales from the final residential-modeling extract used to construct the evaluation samples. The residential scope follows the property types represented in the translated CCAO residential pipeline, including single-family homes, townhomes, small multifamily properties, and condominiums; following the office pipeline, we exclude multicard PINs and sales marked as outliers.
The extract is identified by content rather than by a version string. Every
sample reported here is constructed from
\texttt{data/CCAO/2025/training\_data.parquet}, of 215,400,916 bytes, with sha256
\begin{center}
\texttt{\small b1fc00b514041af5aa7135d85ed59a028afe40569bba1014c4dcf3ad2f3a7b51}
\end{center}
as recorded in the frozen preflight provenance. Two limitations belong
with that identification rather than after it. The file is hash-pinned but not
redistributed: it is not part of the replication package, and no claim of public
availability is made for it. And no build or retrieval date was recorded when the extract
was taken, so none is asserted here -- the hash fixes \emph{which} extract was used, not
when it was produced.


Eligible 2016--2024 sales are ordered chronologically. The oldest 90\% (344,607 sales, through 2023-11-09) form the development/CV pool, and the newest 10\% (38,290 sales, 2023-11-09 through 2024-12-31) form the primary held-out evaluation. The separate 2025 block contains 26,641 sales and is evaluated after each specification is refit on all 382,897 eligible 2016--2024 sales.
The original out-of-time paths had already been inspected during paper development. For the later lower-$\rho$ resolution extension described in Section~\ref{subsec:path_design}, however, the augmented chronological-CV result was frozen before the newly added lower-tail points were evaluated on the held-out and 2025 samples.
Because earlier project development inspected 2024 outputs, we describe the primary sample as held out rather than as an untouched or preregistered confirmatory test.

Because the 90/10 split is defined on chronologically ordered observations rather than on distinct calendar dates, sales occurring on the cutoff date can appear on both sides of the split; no sale observation appears in both samples.

\begin{table}[!ht]
\centering
\caption{Final CCAO temporal samples.}
\label{tab:ccao_samples}
\begin{tabular}{llr}
\toprule
Sample & Period & Observations \\
\midrule
Development / validation & 2016-01-01--2023-11-09 & 344,607 \\
\makecell[l]{
\\Primary held-out evaluation} & 2023-11-09--2024-12-31 & 38,290 \\
\makecell[l]{
\\Full 2016--2024 refit sample} & 2016-01-01--2024-12-31 & 382,897 \\
Secondary forward evaluation & 2025-01-01--2025-12-29 & 26,641 \\
\bottomrule
\end{tabular}
\end{table}
{\footnotesize Observed eligible sales in the final 2016--2024 modeling universe begin on 2016-01-01.}

The model uses 95 predictors, including property characteristics, location and proximity variables, Census measures, time variables, parcel-shape measures, neighborhood identifiers, and other assessor variables. Twenty-three are treated as categorical.
{All predictors are constructed upstream by the CCAO pipeline. Standard and penalized LightGBM fits use the same 95 predictors, with the same 23 variables treated as categorical; the penalty introduces no additional predictors or feature engineering.}

{Table~\ref{tab:feature_groups} summarizes the predictor information used by the models. The predictor set is fixed before model fitting and is identical across the standard and penalized fits. It is defined by the restricted extract identified above and is not redistributed as a separate artifact; Appendix~\ref{app:ccao_results} states what is and is not released.}

{
\begin{table}[!ht]
\centering
\caption{CCAO predictor groups used in the empirical models.}
\label{tab:feature_groups}
\begin{tabular}{lrl}
\toprule
Feature group & Count & Examples \\
\midrule
Assessor/property variables & 30 & Class, building area, land area, age, rooms, neighborhood \\
Location variables & 9 & Coordinates, census tract, school district, municipality \\
Proximity variables & 26 & Transit, parks, roads, hospitals, water, airports \\
ACS 5-year measures & 16 & Income, education, poverty, employment, tenure \\
Time variables & 8 & Sale year, month, quarter, day, post-COVID indicator \\
Parcel-shape variables & 6 & Edge/angle dispersion, bounding ratios, vertex count \\
\midrule
Total & 95 & 23 categorical predictors in LightGBM \\
\bottomrule
\end{tabular}
\end{table}
}


\subsection{Temporal Validation and Regularization-Path Design}
\label{subsec:path_design}

Within the 344,607-sale development/CV pool, we use seven expanding-window
rolling-origin folds. The cumulative windows advance in 15-month steps anchored at
2016-01-01, and within each cumulative window the newest 10\% of observations form
validation. For fold $v$, let $\mathcal T_v$ and $\mathcal V_v$ denote its training
and validation samples. The training window expands across folds, and all
training-dependent quantities used by the proposed objectives, including $\bar y$
and the centered log-price deviations, are recomputed from $\mathcal T_v$ only.
The exact fold dates and sample sizes are reported in Appendix~\ref{app:ccao_path_details}, Table~\ref{tab:fold_structure}; because the newest-10\% validation rule is identical in every fold, it is stated once rather than repeated as a table column.

\paragraph{Development coordinates.}
Three development coordinates appear below, and they are not interchangeable.
The primary coordinate, D1, preserves the CCAO-inspired overlapping rolling-origin
workflow: it is the equal-weight mean of the seven chronological fold metrics. D3 is the
one-sale-one-vote sensitivity, in which every distinct sale carries total weight exactly
one through $w_{ik}=1/m_i$ across the $m_i$ folds in which sale $i$ appears; it is defined
on the development pool alone and covers 130,165 distinct sales. D2 is the
duplicate-weighted pooled out-of-fold coordinate, retained as a historical and
reproducibility sensitivity and never elevated over D1 or D3.

The three differ because of a property of the fold construction that has to be
stated before any result is read. The fold-6 and fold-7 origins are only about four months
apart, while the rule reserving the most recent tenth of each cumulative window for
validation spans close to a year of the development pool, so those two validation blocks
overlap. The concatenated pooled sample therefore contains 151,153 appearances of 130,165
distinct sales: 20,988 sales appear twice and none appears more often, which the frozen
accounting records as a duplicated share of 13.9\%.

D1's status under that overlap is exactly five things, and the last two are the
ones easily lost. D1 remains the primary coordinate; its mathematical definition is
unchanged; every fold-level prediction is still genuinely out of training sample, because
no overlapping row lies in the training block of the model that predicted it; but D1 is
\emph{not} an independent-observation aggregate, since some sales contribute to more than
one fold-specific value; and fold-level statistics are therefore not independent. We do
not describe D1 as unaffected by the overlap. It follows that the fold-level standard
deviations reported alongside fold means are descriptive chronological-window variation
and not standard errors: they support no IID reading of the seven folds and no across-fold
significance statement. D2 is the only one of the three coordinates that double-counts by
construction, and where the nonlinearity gap is computed on that pooled sample a composite
\texttt{fold|row\_id} identifier is supplied to the otherwise unchanged frozen estimator,
so the quantity stays a function of observation identity alone.

The experiment traces the full regularization paths rather than choosing a penalty
strength at this stage.
The substantive comparison contains standard LightGBM and the positive-penalty Direct and Surrogate LightGBM families.
Each family was first evaluated on 50 geometrically spaced positive values from $0.1$ to $100$. A one-time lower-tail extension then added 32 positive values between approximately $0.00110$ and $0.08685$, aligned to the same logarithmic spacing as the original grid, to increase resolution near zero. The final augmented path therefore contains 82 positive values per family, with the originally fitted $0.1$--$100$ points retained unchanged.
Because the direct and surrogate objectives have different numerical scales, their $\rho$
values are not interpreted as comparable magnitudes. The LightGBM hyperparameters were tuned on the same seven chronological folds before any penalty path was estimated. The tuning criterion was the equal-weight mean of fold validation price RMSE, $\operatorname{RMSE}_P(S)=\sqrt{n_S^{-1}\sum_{i\in S}(\widehat P_i-P_i)^2}$. The adopted configuration uses 994 trees and was then frozen for ordinary LightGBM, Direct, and Surrogate fits. The held-out and 2025 samples were not used for this hyperparameter choice. With that vector held fixed, $\rho$ is the only dimension varied within each penalized family. The complete parameter vector is archived in the frozen preflight provenance together with its configuration hash, which is re-verified there against the executed configuration file and against both experiment specifications; that hash is quoted in Appendix~\ref{app:implementation}, so the main text does not list every parameter.


For each configuration and fold, we retain the complete set of predictive,
valuation-level, ratio-dispersion, vertical-equity, and mechanism/residual-structure diagnostics.
We report equal-weight fold means and standard deviations to describe the
chronological behavior of the regularization paths.
No deployment operating point is selected in the present analysis. The chronological-CV path is used only to characterize the path.

For compact presentation, we use five nominal decade-spaced positive display anchors, $\rho\in\{0.01,\allowbreak 0.1,\allowbreak 1,\allowbreak 10,\allowbreak 100\}$, and represent each nominal value by the nearest already-tested point on the frozen 82-point positive grid. The corresponding tested values are $0.0104811$, $0.1$, $0.954095$, $10.4811$, and $100$. This display convention is adopted only to make the tables and ratio-shape figures easy to read; it is not part of the frozen transition rule and does not select a penalty strength. Ordinary LightGBM remains the unpenalized baseline.

All 82 positive values per family remain in the final path figures; no tested value is dropped from the reported paths.

Ratio-shape figures additionally show the zero-penalty path origin for visual continuity. Because the Direct and Surrogate objectives have different numerical scales, equal nominal display anchors are common presentation coordinates only and must not be interpreted as equal-strength interventions across families.


The same final augmented regularization paths are evaluated on the two out-of-time samples. For the held-out evaluation, every valid configuration is refit on the full
344,607-sale development pool and evaluated on the 38,290-sale held-out sample. For
the 2025 forward evaluation, the same configuration is refit on all 382,897 eligible
2016--2024 sales and evaluated on the 26,641-sale 2025 sample.
These held-out and forward paths are descriptive. For the v2 lower-tail extension, the augmented grid and its CV-derived transition result were frozen before the newly added lower-$\rho$ points were evaluated out of time; those new held-out and forward results were not used to revise the grid, the transition rule, the objective definitions, or the frozen LightGBM configuration. The chronological CV, held-out, and 2025 paths are therefore complementary temporal views of the final augmented experiment, not sequential stages that select a penalty strength.

All scalar evaluation statistics are computed on the full relevant validation or
evaluation sample after predictions are transformed back to the price scale. For
the ratio-shape figures in the Results, each evaluation sample is partitioned into
30 equal-count sale-price bins and we plot the median valuation ratio in each bin
against sale price on a logarithmic horizontal axis. The binning is used only for
visualization; all scalar metrics are computed on the full evaluation sample.

\begin{table}[t]
\centering
\caption{
CCAO application and regularization-path design.}
\label{tab:ccao_design}
\begin{tabularx}{\textwidth}{@{}p{3.5cm}X@{}}
\toprule
Item & Design \\
\midrule
Population & Verified Cook County residential sales \\
Development / CV & Oldest 90\% of eligible 2016--2024 sales: 344,607, through 2023-11-09 \\
Primary held-out evaluation & Newest 10\%: 38,290 sales, 2023-11-09--2024-12-31 \\
Full 2016--2024 refit sample & 382,897 eligible 2016--2024 sales \\
2025 forward evaluation & 26,641 sales, 2025-01-01--2025-12-29 \\
Predictors & 95 total; 23 categorical \\
Validation & Seven expanding-window rolling-origin folds; newest 10\% validates \\
Base nonlinear learner & LightGBM on log sale price; 994 trees; hyperparameters frozen before the penalty sweep \\
Corrections evaluated & Direct squared covariance and sample-additive covariance upper bound \\
Penalty path & 82 positive values per family: original 50-point $[0.1,100]$ path plus a 32-point lower-tail extension $[0.00110,0.08685]$ \\
Current analysis & Full CV, held-out, and forward paths; descriptive path landmarks only; no screening result and no deployment operating-point selection \\
\bottomrule
\end{tabularx}
\end{table}


\subsection{Models and Comparisons}
\label{subsec:comparators}

The empirical comparison in this paper covers the baseline model class, the custom-objective implementation control, and the Direct and Surrogate training-time regularization paths, together with a global post-hoc centered rescaling of the same unregularized predictors. That comparator is reported in Sections~\ref{subsec:posthoc_comparator} and~\ref{subsec:matched_beta_results}, both descriptively and at matched development first-order correction. \emph{Structured} post-hoc calibration by geography, property class or time is a different construction and remains outside the empirical scope.

\paragraph{Reference cells, and what each one licenses.}
Three reference cells appear below, and the distinction between two of them is
load-bearing for every comparison in the Results. \emph{Ordinary LightGBM (standard
raw-label native)} is the assessor-facing workflow benchmark: the standard pipeline as an
office would run it. \emph{Centered-label native L2 (initialization-aligned)} is an
implementation-decomposition control only, and no full regularization path is traced for
it. \emph{Custom-objective $\rho=0$ origin} is the primary within-path penalty-isolating
reference: the same custom-objective code path as the penalized fits, evaluated at zero
penalty.

The attribution rule that goes with them is not a nicety. Comparisons against
the workflow benchmark are descriptively valid, and they are the practitioner's question
-- can a different objective beat the standard pipeline? -- but a change relative to it
must not be attributed to $\rho$ alone, because that cell differs from the penalized fits
in execution path as well as in penalty. The custom-objective $\rho=0$ origin to
positive-$\rho$ contrast, taken within the same family, is therefore the primary
within-path penalty-isolating comparison: it isolates the change associated with
activating the penalty within the same frozen implementation path. Where both comparisons
are reported they are labelled as the different questions they are, and no
workflow-benchmark difference is presented as a pure penalty effect.
Appendix~\ref{app:implementation} reports the $\rho=0$ implementation audit that makes
the distinction necessary rather than merely cautious.

\paragraph{The workflow benchmark.}
Standard LightGBM, trained with the ordinary squared log-price objective, is the workflow benchmark for every comparison reported here; the penalized and post-hoc configurations use the same outcome, development observations, and underlying predictor information under the common temporal design. The earlier township-level linear-regression workflow is historical context only: it is not re-estimated, and no quantitative linear-versus-boosted comparison is reported.

\paragraph{Direct and surrogate penalties.}
The direct LightGBM comparator uses the diagonal-Hessian approximation described in Appendix~\ref{app:implementation}: its covariance gradient is exact, while the cross-observation second-order curvature is omitted so that the objective can be used through the standard LightGBM interface. The surrogate uses the exact observation-additive objective from Section~\ref{subsec:surrogate}.

\textbf{Post-hoc centered rescaling.} A global one-parameter post-hoc centered rescaling of an already fitted predictor is \emph{inside} the empirical scope of this paper. The fixed-space Direct result in Appendix~\ref{app:theory} makes that comparison scientifically natural, and Sections~\ref{subsec:posthoc_comparator} and~\ref{subsec:matched_beta_results} carry it out: first descriptively, then at matched development first-order correction, with the rescaling of the custom-objective $\rho=0$ origin as the primary post-hoc family and the rescaling of the workflow benchmark as a secondary practical variant. The paper's contribution remains the design, implementation, mechanism analysis and evaluation of training-time covariance-guided corrections; the comparator is what keeps that contribution from resting on an untested assumption that retraining is necessary.


In many settings, richer spatial, temporal, neighborhood, and property information can improve both predictive accuracy and assessor-facing equity measures \cite{SmithEtAl2026}, but that is a modeling extension rather than a comparator in the headline experiment. We therefore discuss richer information as a separate direction in Section~\ref{sec:discussion} rather than treating it as a completed empirical control here.


\subsection{Scope of External Evidence}
\label{subsec:external_scope}

The primary evidentiary claim is CCAO-specific. A newer protocol-frozen, normalized multi-jurisdiction benchmark exists in the project repository, but it uses commercially harmonized data and a common research pipeline rather than the local data, legal rules, feature engineering, or review procedures of the participating assessor offices. Its screening labels and aggregate success summaries require a separate inferential and implementation audit before manuscript integration. We therefore do not use that benchmark to strengthen the paper's main claim or to imply transfer into another office's production workflow. The legacy six-county experiment is retained in the source history below but removed from the active manuscript because it is superseded and was developed iteratively.


\section{Results}
\label{sec:results}


The primary CCAO Results report the complete Direct and Surrogate regularization paths. The fixed cross-fitted correlation-ratio nonlinearity gap $\Delta_{\mathrm{NL}}$ is shown along those paths: the held-out and 2025 curves are the original out-of-time diagnostics, while the chronological-CV curves are reconstructed from the already-frozen validation predictions using the same fixed estimator. The artifact that reproduces the positive-$\rho$ path anchors carries no $\Delta_{\mathrm{NL}}$ column, so no positive-$\rho$ path-anchor value of it is reported for these paths and no screening result is derived from it; positive-$\rho$ values of the gap are reported instead at the matched first-order and crossover coordinates, which rest on their own frozen sources. No deployment operating point is selected.
We also report the chronological-fold behaviour of those paths, together with the temporal failure-mode audit of Section~\ref{subsec:path_stability_results}. These characterize the paths and define no selected or recommended range.

One implementation qualification applies to every headline path comparison, and Section~\ref{subsec:zero_reference_result} states it as a rule rather than a caveat. In the frozen CCAO artifacts the two custom-objective fits at $\rho=0$ agree with each other but are not numerically identical to ordinary LightGBM, despite algebraically matching squared-error derivatives. Positive-$\rho$ changes relative to the within-family $\rho=0$ origin describe the effect of regularization along that custom-objective training path; differences relative to ordinary LightGBM also carry the native-versus-custom execution-path difference. Appendix~\ref{subsec:zero_control_results} reports the audit.

\subsection{Accuracy Alone Does Not Remove the CCAO Ratio Trend}
\label{subsec:ccao_baseline_results}

Section~\ref{subsec:baseline_tension} reports the level of the assessor-facing
diagnostics for \emph{Ordinary LightGBM (standard raw-label native)}, the workflow benchmark, on the
held-out and 2025 forward samples, re-sourced from the frozen full zero-penalty control; that level
involves no penalized-model choice and no comparison with another model class. The remainder of the
primary CCAO Results treats the \emph{regularization path itself} as the empirical
object of interest. No $\rho$, penalty family, or operating point is selected in the
present analysis.


\subsection{The Zero-Penalty Reference Cells and What They Permit}
\label{subsec:zero_reference_result}

Every path statement below is read against one of two reference cells, and the
difference between them governs how the rest of the Results is organized. \emph{Ordinary LightGBM
(standard raw-label native)} is the assessor-facing workflow benchmark: the standard pipeline as an
office would run it. \emph{Centered-label native L2 (initialization-aligned)} is an
implementation-decomposition control only, and no regularization path is traced for it.
\emph{Custom-objective $\rho=0$ origin} is the primary within-path penalty-isolating reference: the
same custom-objective code path as the penalized fits, evaluated at zero penalty.

Three findings from the frozen zero-penalty audit fix how those cells may be used. First,
the two custom-objective controls are the same object: at $\rho=0$ the Direct and Surrogate
predictions coincide exactly, with a maximum absolute log difference of $0$ in every audited
capacity-by-sample cell. Second, the custom-objective origin is \emph{not} numerically identical to
ordinary LightGBM. Evaluated on the same held-out and 2025 samples, the two differ by a mean
absolute log deviation of 3.24e-02, with a maximum of 3.09e-01, on the held-out sample, and by
3.06e-02, with a maximum of
2.83e-01, on the 2025 forward sample. Those differences are large enough to move reported metrics
and are not a rounding artifact. Third, the audit locates the difference in the execution path
rather than in the objective: at the first boosting iteration the supplied zero-penalty gradients
and Hessians match native squared-error derivatives exactly, so nothing in the covariance objective
is implicated. Appendix~\ref{subsec:zero_control_results} reports the audit table and the
implementation decomposition behind it.

The consequence is a rule rather than a caveat. A difference between ordinary LightGBM
and a penalized fit is a workflow comparison: it carries the execution-path difference as well as
the penalty, and it must not be attributed to $\rho$ alone. The comparison that isolates the penalty
is the within-family one, from the custom-objective $\rho=0$ origin to positive $\rho$ along the same
frozen implementation path. Both are reported throughout, marked with different devices, and
labelled as the different questions they are. Because the two zero-penalty custom-objective controls
coincide, that origin is shared by the Direct and Surrogate paths, and it is therefore also the
common starting point against which the two penalized families are compared with each other.


\subsection{\texorpdfstring{
What the Final Augmented Regularization Paths Show}{What the Final Augmented Regularization Paths Show}}
\label{subsec:path_results}

The Direct and Surrogate objectives are evaluated over their complete final augmented grids. Since the two penalties have different numerical scales, their
$\rho$ values are not compared as equal-strength interventions. The principal
question at this stage is descriptive: how do prediction, valuation level,
horizontal uniformity, assessor-facing vertical-equity diagnostics, and the targeted
dependence measures evolve as regularization increases?

The compact main-text view is split in two, because the provenance of its two
halves differs. Table~\ref{tab:path_anchor_frozen} reports prediction together with PRD and MKI at
the two positive penalties whose fitted artifacts the frozen reproduction audit re-derives at
reproduction tier R1 --- tested $\rho\approx 0.954095$ and $\rho=100$, where cached and refit metrics
agree at the displayed precision --- alongside both reference cells.
Table~\ref{tab:path_anchor_standards} reports the two standards-facing vertical-equity measures ---
PRB, from the adopted 2013 Standard on Ratio Studies, and VEI, from the May-2026 Exposure Draft and so
proposed rather than adopted --- at all five nominal decade-spaced display anchors, because those two
re-derive exactly at every anchor from the frozen inferential tables. The three remaining display anchors therefore do not
appear in Table~\ref{tab:path_anchor_frozen} at all: their prediction, PRD and MKI values rest on
legacy fitted artifacts that are hash-pinned but present in no checkout, so no frozen artifact
reproduces them and they are not printed here. Appendix
Table~\ref{tab:path_anchor_complementary} reports the complementary valuation-level,
horizontal-uniformity and first-order-mechanism diagnostics at the same two frozen anchors. The full
82-point positive path for each family remains the scientific object of analysis; the anchor rows are
presentation-only snapshots and do not constitute model selection.


{\begin{table}[!htbp]
\centering
\scriptsize
\setlength{\tabcolsep}{3.4pt}
\renewcommand{\arraystretch}{1.08}
\caption{Regularization path at the two frozen positive-$\rho$ anchors: prediction, PRD
and MKI, with both reference cells shown as explicit rows. Boldface marks movement toward the
measure's own preferred direction relative to \emph{Ordinary LightGBM (standard raw-label native)},
the assessor-facing workflow benchmark; a dagger (\dag) marks movement toward that direction relative
to the \emph{Custom-objective $\rho=0$ origin}, the primary within-path penalty-isolating reference.
The two reference rows carry no markers. Neither device denotes statistical significance, standards
compliance, model selection, or a preferred penalty strength: no penalty strength is selected
anywhere in this paper.}
\label{tab:path_anchor_frozen}
\resizebox{\textwidth}{!}{%
\begin{tabular}{@{}llrrrrrr@{}}
\toprule
Method & $\rho$
& $R^2_P$
& MAE
& MAPE
& $\operatorname{RMSE}_{\log P}$
& PRD
& MKI \\
\midrule
\multicolumn{8}{l}{\textit{Panel A: Held-out evaluation}} \\
\addlinespace[2pt]
\multicolumn{8}{l}{\textbf{Reference cells}} \\
\cmidrule(lr){1-8}
Ordinary LightGBM (standard raw-label native) & -- & 0.894 & \$75,655 & 21.2\% & 0.289 & 1.069 & 0.923 \\
Custom-objective $\rho=0$ origin & $0$ & 0.893 & \$75,976 & 21.2\% & 0.289 & 1.069 & 0.923 \\
\addlinespace[3pt]
\multicolumn{8}{l}{\textbf{Direct penalty}} \\
\cmidrule(lr){1-8}
Direct & $\approx 0.954095$ & \textbf{0.899}\textsuperscript{\dag} & \textbf{\$74,485}\textsuperscript{\dag} & \textbf{21.1\%}\textsuperscript{\dag} & 0.290 & \textbf{1.060}\textsuperscript{\dag} & \textbf{0.940}\textsuperscript{\dag} \\
Direct & $100$ & 0.869 & \$84,918 & 23.9\% & 0.325 & \textbf{1.029}\textsuperscript{\dag} & \textbf{0.998}\textsuperscript{\dag} \\
\addlinespace[3pt]
\multicolumn{8}{l}{\textbf{Surrogate penalty}} \\
\cmidrule(lr){1-8}
Surrogate & $\approx 0.954095$ & \textbf{0.897}\textsuperscript{\dag} & \textbf{\$75,468}\textsuperscript{\dag} & \textbf{21.0\%}\textsuperscript{\dag} & 0.298 & \textbf{1.049}\textsuperscript{\dag} & \textbf{0.953}\textsuperscript{\dag} \\
Surrogate & $100$ & 0.889 & \$82,094 & 22.8\% & 0.357 & \textbf{1.010}\textsuperscript{\dag} & \textbf{1.013}\textsuperscript{\dag} \\
\midrule
\addlinespace[2pt]
\multicolumn{8}{l}{\textit{Panel B: 2025 forward evaluation}} \\
\addlinespace[2pt]
\multicolumn{8}{l}{\textbf{Reference cells}} \\
\cmidrule(lr){1-8}
Ordinary LightGBM (standard raw-label native) & -- & 0.904 & \$78,484 & 20.8\% & 0.278 & 1.079 & 0.907 \\
Custom-objective $\rho=0$ origin & $0$ & 0.904 & \$79,139 & 20.8\% & 0.279 & 1.078 & 0.909 \\
\addlinespace[3pt]
\multicolumn{8}{l}{\textbf{Direct penalty}} \\
\cmidrule(lr){1-8}
Direct & $\approx 0.954095$ & \textbf{0.910}\textsuperscript{\dag} & \textbf{\$77,139}\textsuperscript{\dag} & \textbf{20.7\%}\textsuperscript{\dag} & 0.279\textsuperscript{\dag} & \textbf{1.069}\textsuperscript{\dag} & \textbf{0.926}\textsuperscript{\dag} \\
Direct & $100$ & 0.873 & \$89,309 & 23.1\% & 0.307 & \textbf{1.036}\textsuperscript{\dag} & \textbf{0.983}\textsuperscript{\dag} \\
\addlinespace[3pt]
\multicolumn{8}{l}{\textbf{Surrogate penalty}} \\
\cmidrule(lr){1-8}
Surrogate & $\approx 0.954095$ & \textbf{0.908}\textsuperscript{\dag} & \textbf{\$78,196}\textsuperscript{\dag} & \textbf{20.4\%}\textsuperscript{\dag} & 0.284 & \textbf{1.060}\textsuperscript{\dag} & \textbf{0.932}\textsuperscript{\dag} \\
Surrogate & $100$ & 0.900 & \$85,257 & 22.0\% & 0.336 & \textbf{1.020}\textsuperscript{\dag} & \textbf{0.987}\textsuperscript{\dag} \\
\bottomrule
\end{tabular}}
\vspace{1mm}
\begin{minipage}{\textwidth}
\scriptsize
\emph{Notes.}
The two rows under \textit{Reference cells} are the two frozen reference roles, and they
answer different questions. \emph{Ordinary LightGBM (standard raw-label native)} is the
assessor-facing workflow benchmark: the standard pipeline as an office would run it. The
\emph{Custom-objective $\rho=0$ origin} is the primary within-path penalty-isolating reference: the
same custom-objective code path as the penalized fits, evaluated at zero penalty. Comparisons against
the workflow benchmark are descriptively valid, but a difference relative to it must not be attributed
to $\rho$ alone, because that cell differs from the penalized fits in execution path as well as in
penalty; only the $\rho=0$ origin to positive-$\rho$ contrast isolates the penalty. Reference-cell
values are taken from the frozen full zero-penalty control, and the two positive-$\rho$ rows from the
frozen reproduction audit. Direct and Surrogate penalty magnitudes are not numerically comparable
across families. The nominal display anchors $\{0.01,0.1,1,10,100\}$ map to tested values
$\{0.0104811,0.1,0.954095,10.4811,100\}$; only the last two of those reproduce here, and the other
three are omitted rather than shown blank.
Improvement is defined against each measure's own direction: higher $R^2_P$; lower MAE, MAPE and
$\operatorname{RMSE}_{\log P}$; PRD and MKI closer to one. The adopted 2013 guidance range for PRD is
$[0.98,1.03]$, whereas the MKI range $[0.95,1.05]$ is a May-2026 Exposure Draft diagnostic band rather
than an adopted compliance standard, so a marker on MKI is directional only. Comparisons are evaluated
on the unrounded canonical values, so two cells that appear tied at the displayed precision can still
be marked differently. MAPE is reported in percent and MAE in dollars.
\end{minipage}
\end{table}
}

At the nominal $\rho=1$ display anchor (tested $\rho\approx 0.954095$) both families improve
price-scale accuracy relative to the penalty-isolating origin and relative to the workflow benchmark.
Direct reaches held-out $R^2_P=0.899$ with MAE of \$74,485, against $0.893$ and \$75,976 at the
\emph{Custom-objective $\rho=0$ origin} and $0.894$ and \$75,655 at \emph{Ordinary LightGBM (standard
raw-label native)}; the corresponding 2025 forward figures are $0.910$ and \$77,139, against $0.904$
and \$79,139 at the origin and $0.904$ and \$78,484 at the workflow benchmark. PRD and MKI move toward
one at the same anchor. At $\rho=100$ the price-scale picture reverses while PRD and MKI keep moving
toward one: Direct's held-out $R^2_P$ falls to $0.869$ with MAE of \$84,918, and the Surrogate
$\operatorname{RMSE}_{\log P}$ degrades from $0.289$ at the origin to $0.357$, its largest displayed
log-scale cost.
Thus the empirical paths are not a simple monotone exchange of prediction accuracy for vertical equity. At stronger penalties, price-scale accuracy eventually deteriorates, while some vertical-equity diagnostics approach or cross their neutral values and others become non-monotone.


{\begin{table}[!htbp]
\centering
\scriptsize
\setlength{\tabcolsep}{5pt}
\renewcommand{\arraystretch}{1.08}
\caption{Standards-facing vertical equity at all five nominal display anchors, with the
two measures' guidance status kept separate: PRB belongs to the adopted 2013 Standard on Ratio
Studies, while VEI is a May-2026 Exposure Draft proposal and carries no adopted compliance status.
Boldface marks movement toward zero relative to \emph{Ordinary LightGBM (standard raw-label native)},
the assessor-facing workflow benchmark; a dagger (\dag) marks movement toward zero relative to the
\emph{Custom-objective $\rho=0$ origin}, the primary within-path penalty-isolating reference. The two
reference rows carry no markers. Neither device denotes statistical significance, standards
compliance, model selection, or a preferred penalty strength: no penalty strength is selected anywhere
in this paper.}
\label{tab:path_anchor_standards}
\resizebox{\textwidth}{!}{%
\begin{tabular}{@{}llrrrr@{}}
\toprule
& & \multicolumn{2}{c}{PRB (adopted 2013 Standard)}
  & \multicolumn{2}{c}{VEI (May-2026 Exposure Draft, proposed)} \\
\cmidrule(lr){3-4}\cmidrule(l){5-6}
Method & $\rho$ & Held-out & 2025 forward & Held-out & 2025 forward \\
\midrule
\multicolumn{6}{l}{\textbf{Reference cells}} \\
\cmidrule(lr){1-6}
Ordinary LightGBM (standard raw-label native) & -- & -0.091 & -0.106 & -26.5\% & -28.6\% \\
Custom-objective $\rho=0$ origin & $0$ & -0.088 & -0.103 & -26.2\% & -28.4\% \\
\addlinespace[3pt]
\multicolumn{6}{l}{\textbf{Direct penalty}} \\
\cmidrule(lr){1-6}
Direct & $\approx 0.0104811$ & \textbf{-0.089} & \textbf{-0.105} & -26.9\% & -29.8\% \\
Direct & $0.1$ & \textbf{-0.088}\textsuperscript{\dag} & \textbf{-0.102}\textsuperscript{\dag} & -27.0\% & \textbf{-27.7\%}\textsuperscript{\dag} \\
Direct & $\approx 0.954095$ & \textbf{-0.075}\textsuperscript{\dag} & \textbf{-0.090}\textsuperscript{\dag} & \textbf{-21.9\%}\textsuperscript{\dag} & \textbf{-23.9\%}\textsuperscript{\dag} \\
Direct & $\approx 10.4811$ & \textbf{-0.040}\textsuperscript{\dag} & \textbf{-0.049}\textsuperscript{\dag} & \textbf{-10.5\%}\textsuperscript{\dag} & \textbf{-11.9\%}\textsuperscript{\dag} \\
Direct & $100$ & \textbf{-0.015}\textsuperscript{\dag} & \textbf{-0.029}\textsuperscript{\dag} & \textbf{0.7\%}\textsuperscript{\dag} & \textbf{-4.5\%}\textsuperscript{\dag} \\
\addlinespace[3pt]
\multicolumn{6}{l}{\textbf{Surrogate penalty}} \\
\cmidrule(lr){1-6}
Surrogate & $\approx 0.0104811$ & \textbf{-0.088}\textsuperscript{\dag} & \textbf{-0.101}\textsuperscript{\dag} & \textbf{-25.8\%}\textsuperscript{\dag} & \textbf{-27.4\%}\textsuperscript{\dag} \\
Surrogate & $0.1$ & \textbf{-0.079}\textsuperscript{\dag} & \textbf{-0.095}\textsuperscript{\dag} & \textbf{-22.7\%}\textsuperscript{\dag} & \textbf{-24.7\%}\textsuperscript{\dag} \\
Surrogate & $\approx 0.954095$ & \textbf{-0.045}\textsuperscript{\dag} & \textbf{-0.062}\textsuperscript{\dag} & \textbf{-12.6\%}\textsuperscript{\dag} & \textbf{-15.3\%}\textsuperscript{\dag} \\
Surrogate & $\approx 10.4811$ & \textbf{0.013}\textsuperscript{\dag} & \textbf{0.000}\textsuperscript{\dag} & \textbf{0.2\%}\textsuperscript{\dag} & \textbf{-2.0\%}\textsuperscript{\dag} \\
Surrogate & $100$ & \textbf{0.042}\textsuperscript{\dag} & \textbf{0.030}\textsuperscript{\dag} & \textbf{4.0\%}\textsuperscript{\dag} & \textbf{1.7\%}\textsuperscript{\dag} \\
\bottomrule
\end{tabular}}
\vspace{1mm}
\begin{minipage}{\textwidth}
\scriptsize
\emph{Notes.}
PRB is the price-related bias measure of the adopted 2013 Standard on Ratio Studies, whose
acceptable range is $[-0.05,0.05]$. VEI is taken from the May-2026 Exposure Draft; its $\pm 10\%$ band
is proposed guidance rather than an adopted compliance standard, so no VEI entry in this table should
be read as a compliance determination. All forty PRB and VEI entries re-derive from the frozen
inferential tables, which is why this table reaches all five anchors while
Table~\ref{tab:path_anchor_frozen} reaches only two.
Two entries need their own qualification. The Direct held-out VEI at the nominal $\rho=10$ anchor is
-10.5\%: it is a movement toward zero relative to both reference cells, but it still lies just outside
the proposed $\pm 10\%$ Exposure-Draft band, and it is therefore not a compliance result. The
Exposure Draft does not treat a point estimate outside that band as a finding on its own: its
procedure escalates such a case to a significance test, which this paper does not report. The
Surrogate 2025 PRB at that same anchor displays as 0.000; the frozen estimate is 0.000135 with a 95\%
confidence interval running from $-0.004$ to $0.005$, so it is a rounded near-zero estimate and not
exact first-order neutrality. Under the adopted 2013 classification the entire 95\% confidence interval must
lie outside a band before that band is treated as exceeded, and an interval that merely crosses a
threshold does not qualify.
Boldface and the dagger are defined against distance from zero, evaluated on the unrounded canonical
values, so two cells that appear tied at the displayed precision can still be marked differently.
Direct and Surrogate penalty magnitudes are not numerically comparable across families. VEI is
reported in percent.
\end{minipage}
\end{table}
}


The anchor rows in Tables~\ref{tab:path_anchor_frozen}
and~\ref{tab:path_anchor_standards} remain descriptive snapshots of the full paths.


\subsection{How the Valuation-Ratio Shape Changes Along the Path}
\label{subsec:ratio_shape_results}


Scalar diagnostics do not show where along the value distribution the ratio pattern changes. We therefore complement the sweep table with a selection-free ratio-shape view using the same 30 equal-count-bin construction as the baseline motivation figure, the zero-penalty path origin, and the nominal decade-spaced positive display anchors $\{0.01,0.1,1,10,100\}$ mapped to their nearest tested grid points.

\begin{figure}[!htbp]
\centering
\safeincludegraphics[width=0.8\textwidth]{img/generated_v12_994/ratio_shape_evolution.pdf}
\caption{Valuation-ratio profiles against sale price at the zero-penalty path origin and nominal decade-spaced positive display anchors $\rho\in\{0.01,0.1,1,10,100\}$, each represented by the nearest tested grid point. The horizontal line at 1 is the principal neutrality reference; lines at 0.9 and 1.1 are aggregate appraisal-level reference guides and are not binwise acceptance criteria. The display anchors are presentation-only and do not select a penalty strength.}
\label{fig:ratio_shape_path_placeholder}
\end{figure}

Both formulations attenuate the broad decline in valuation ratios as regularization
increases, but they change the profile differently.
Across the decade-spaced display snapshots, the Direct profile changes more smoothly than the Surrogate profile. Under stronger Surrogate regularization, the profile
develops a pronounced non-monotone trough through the lower and middle sale-price
range before recovering at higher prices. This visual contrast is descriptive rather than a structural guarantee: Direct covariance controls the affine residual--price component but does not protect against nonlinear conditional-mean departures at stronger regularization. The ratio-shape contrast here is therefore read from the profiles themselves; the two families' non-affine residual structure is compared numerically at matched first-order correction in Section~\ref{subsec:matched_beta_results}, not along these paths.
Consequently, a near-neutral global slope or high-versus-low grouped statistic need not imply a flat valuation-ratio profile across the value distribution.

The fixed cross-fitted $\Delta_{\mathrm{NL}}$ nonlinearity gap is reported here only where a frozen artifact reproduces it. At the shared zero-penalty path origin it is $0.116$ on the held-out sample and $0.121$ in 2025. The frozen artifact that reproduces the positive-$\rho$ display anchors carries no $\Delta_{\mathrm{NL}}$ column, so no positive-$\rho$ path-anchor value of the gap is quoted from it and no claim is made about how the gap moves along these paths. (The matched first-order and crossover comparisons reported later are different coordinates resting on their own frozen sources, and supported positive-$\rho$ values of the gap do appear there.) The ratio-shape contrast described above is therefore read from the profiles themselves rather than from the scalar gap.

\subsection{
\texorpdfstring{
First-Order Correction and Residual Structure Along the Path}{First-Order Correction and Residual Structure Along the Path}
}
\label{subsec:mechanism_results}


{The Direct training target is covariance, and $\beta_{\log}$ is its signed first-order slope equivalent up to the fixed variance of log sale price within an evaluation sample. The correlation-ratio nonlinearity gap $\Delta_{\mathrm{NL}}$ then measures conditional-mean explanatory power beyond the best affine residual--price relationship, while distance correlation provides a broader omnibus check for remaining distributional dependence. These diagnostics answer nested but distinct questions---first-order trend, non-affine conditional mean, and broader dependence---and are therefore reported at the coordinates whose frozen artifacts carry them: the zero-penalty path origin, and the matched first-order, centered-spread and crossover comparisons of Sections~\ref{subsec:posthoc_comparator} and~\ref{subsec:matched_beta_results}.}

\begin{figure}[!htbp]
\centering
\safeincludegraphics[width=0.8\textwidth]{img/generated_v12_994/mechanism_vs_rho.pdf}
\caption{Mechanism and residual-structure paths versus $\rho$ for Direct and Surrogate on held-out (solid) and 2025 (dashed) evaluations. The $\beta_{\log}=0$ line is a first-order neutrality reference; $\Delta_{\mathrm{NL}}$ and dCor are not given assessor-style zero targets. The $\Delta_{\mathrm{NL}}$ and dCor panels are shown descriptively: the frozen artifact that reproduces the positive-$\rho$ path anchors carries neither column, so no positive-$\rho$ value and no path-turn location is quoted from them, and positive-$\rho$ residual structure is compared instead at the matched first-order coordinate of Section~\ref{subsec:matched_beta_results}. No penalty strength is selected in this paper.}
\label{fig:mechanism_path_placeholder}
\end{figure}

The first-order mechanism moves in the intended direction in both families. Along
the held-out Direct anchors, $\beta_{\log}$ changes from $-0.147$ at
$\rho=0$ to $-0.079$ at $\rho=100$; the corresponding 2025 values move from
$-0.161$ to $-0.093$. The Surrogate path moves closer to first-order neutrality,
reaching $-0.001$ on the held-out sample and $-0.015$ in 2025 at $\rho=100$.
Along these paths the broader distance-correlation diagnostic is quoted only at the zero-penalty origin, where it is $0.382$ on the held-out sample and $0.417$ in 2025. The frozen artifact that reproduces the positive-$\rho$ path anchors carries no distance-correlation column either, so no positive-$\rho$ path-anchor value and no claim about the shape of the dCor path are asserted for these paths. Supported positive-$\rho$ values of both residual-structure diagnostics do exist at the matched first-order and crossover coordinates, which rest on separate frozen artifacts.
Because distance correlation is an omnibus dependence measure, a movement in it would not by itself establish that nonlinear conditional-mean structure had changed; that distinction belongs to $\Delta_{\mathrm{NL}}$, which is unavailable along these paths for the same reason. Neither diagnostic is therefore used to characterize residual structure at positive $\rho$ from these path artifacts; the matched first-order comparison, which has its own frozen source, is where positive-$\rho$ residual structure is compared.
Figure~\ref{fig:mechanism_path_placeholder} traces the first-order coordinate on the complete grids: $\lvert\beta_{\log}\rvert$ declines from $\rho=0$ to $\rho=100$ on Direct and falls further on Surrogate, which reaches near first-order neutrality at the strongest tested penalty. First-order slope reduction and the correlation-ratio nonlinearity gap remain conceptually distinct objects, but along these paths $\Delta_{\mathrm{NL}}$ is quoted only at the zero-penalty origin, so no positive-$\rho$ comparison between the two is asserted for them.

\subsection{Accuracy--Equity Trajectories}
\label{subsec:tradeoff_results}

The central assessor-facing result is the path traced jointly by predictive
performance and vertical-equity diagnostics.

The main tradeoff figure reports all four assessor-facing vertical-equity diagnostics---PRD, PRB, MKI, and VEI---against $R^2_P$. Appendix~\ref{app:ccao_path_details} broadens the accuracy dimension to MAE, MAPE, and $\operatorname{RMSE}_{\log P}$ and separately reports mechanism-versus-accuracy trajectories.
\begin{figure}[!htbp]
\centering
\safeincludegraphics[width=0.98\textwidth]{img/generated_v12_994/accuracy_equity_trajectories_inprocessing_only.pdf}
\caption{Accuracy--equity trajectories with $R^2_P$ on the vertical axis and PRD, PRB, MKI, and VEI on the horizontal axis, for held-out and 2025 evaluations. Ordinary LightGBM is the workflow-benchmark context anchor; the historical linear-regression point is drawn for visual context only and supports no reported quantitative comparison. Dotted vertical lines mark metric neutrality where applicable, and shaded regions are assessor-facing guidance ranges rather than compliance claims. Arrows indicate increasing $\rho$ and do not mark a selected point.}
\label{fig:accuracy_equity_placeholder}
\end{figure}
The trajectories are not simple monotone accuracy--equity frontiers. At moderate regularization, parts of both paths move toward more neutral PRD and VEI while preserving or improving price-scale fit relative to the corresponding custom-objective $\rho=0$ control.
At stronger regularization, $R^2_P$ bends downward while several vertical-equity diagnostics approach or cross their neutral values; the resulting trajectories can themselves become non-monotone. This is why the complete path, rather than a single endpoint or a presumed accuracy cost, is
the relevant empirical object before selecting an operating point. Accordingly, the main accuracy--equity figure does not highlight any segment of either path as preferred, and temporal behaviour is assessed separately below.

\subsection{What Retraining Buys over Post-Hoc Centered Rescaling}
\label{subsec:posthoc_comparator}

Corollary~\ref{cor:path_scaling} shows that inside the exact fixed-prediction-space
benchmark the Direct path \emph{is} a one-parameter centered rescaling of the unpenalized predictor.
That turns a theoretical remark into an empirical question this paper can answer rather than defer:
along the actual CCAO paths, what does retraining buy relative to globally rescaling the same
unregularized predictor at the same first-order correction? This subsection reports the comparator
and the descriptive comparison the frozen record supports;
Section~\ref{subsec:matched_beta_results} then reports the formal comparison at matched
first-order correction.

\paragraph{The map, and the one sensitivity attached to it.}
The comparator rescales a cached prediction array about the mean log \emph{target} of
its own fitting block,
\[
  f_b(x) \;=\; \bar y_T + b\,\bigl(f_0(x)-\bar y_T\bigr),
\]
where $\bar y_T$ is the fold-$k$ training mean for fold $k$, the development-pool mean for the
held-out block, and the 2016--2024 production mean for the 2025 forward block. This is the
theorem-matched form, because Corollary~\ref{cor:path_scaling} writes the fixed-space Direct path in
exactly that shape, and it is the only map used for the full path. At $b=1$ the map returns the
cached array unchanged, so each post-hoc path starts at its own reference cell rather than at a
re-evaluation of it. No model is refit anywhere in this comparator: the inputs are the frozen
prediction arrays, hash-verified before use. The frozen grid holds 124 base values, giving a
path table of 2,480 rows.

A second centering --- about the fitting block's mean \emph{prediction}
$\bar f_{0,T}$ instead of $\bar y_T$ --- is a named sensitivity rather than an alternative primary
map, and the two are never merged. At the same $b$ they differ by $b\,(\bar f_{0,T}-\bar y_T)$, a
constant in $x$, so $\beta_{\log}$, the residual--price covariance, COD, COV and distance
correlation are exactly unchanged by the choice, and only level-sensitive measures can move at all.
Over the blocks on which the two post-hoc paths are traced the largest gap between the two centers
is 5.18e-06 --- across all three reference cells it is bounded by 8.36e-06 --- so at the largest
base value the induced log shift is at most 6.266e-06; the largest metric movement the alternative
centering produces anywhere is \$1.54{} of mean absolute error, and $\beta_{\log}$ moves
by 2.78e-17 --- machine precision, exactly as the constant-shift argument requires. The
sensitivity was therefore evaluated at four base values on each regime for both reference cells
instead of as a duplicate full path, and it is reported here with that bound rather than waved away
as numerically close.

\paragraph{Three roots, and they are three different objects.}
The rescaling coefficient can be defined on three different samples, and collapsing them
would misstate what the comparator does. The \emph{fitting-block} coefficient
$b^{\star}_{T}=\widehat{\Var}_T(y)/\widehat{\Cov}_T(f_0,y)$ of Appendix~\ref{app:theory} is an
in-sample theory diagnostic and sets no empirical endpoint: for the custom-objective $\rho=0$ origin
it runs from 1.065397 to 1.092817 over the fitting blocks, because in sample that learner's
$\beta_{\log}$ lies between -0.061383 and -0.084934 while out of sample it is about
$-0.14$ to $-0.16$, so an in-sample calibration would under-correct by roughly a factor of two in $b-1$. The
\emph{primary empirical} root is the development-coordinate one. On the D1 coordinate
$\beta_{\log}$ is exactly linear in $b$, with
$\beta_k(b)=b\,\widehat{\Cov}_{V_k}(f_0,y)/\widehat{\Var}_{V_k}(y)-1$ per fold and
$\beta_{\mathrm{D1}}(b)$ their equal-weight mean, so the root is
$b_{0}=\bigl[\tfrac{1}{7}\sum_{k}\widehat{\Cov}_{V_k}(f_0,y)/\widehat{\Var}_{V_k}(y)\bigr]^{-1}$. Note what
that expression is built from: fold-specific \emph{validation}-block covariance and variance, not a
training covariance/variance relation. It equals 1.160458 for the custom-objective $\rho=0$
origin and 1.160367 for the workflow benchmark, and it drives the development coordinate to
1.00e-16 against a target of $10^{-12}$. The \emph{third} root is the pooled out-of-fold one at
1.167027 and it is a sensitivity only: the pooled development sample gives the sales in the
overlapping validation blocks double observation weight by construction. The grid runs to
$b_{\max}=1+1.25\,(\max_{\text{roots}}-1)=$ 1.209082 --- a twenty-five percent overshoot in the
adjustment $b-1$ rather than in $b$ itself --- and every root and bound is computed from development
predictions alone. The familiar $1/R^{2}$ expression, as Appendix~\ref{app:theory} states, is the
linear/projection special case rather than the general coefficient for a boosted predictor.

\paragraph{The descriptive comparison, and it crosses over.}
Table~\ref{tab:posthoc_crossover} places Direct's most-corrected held-out
configuration and the Surrogate's near-neutral held-out configuration beside post-hoc rescalings
of the custom-objective $\rho=0$ origin at a comparable held-out first-order correction. The frozen record is explicit that this juxtaposition is keyed on
\emph{held-out} $\beta_{\log}$ and is descriptive: the comparison on the development coordinate,
which is the one that can be matched by design, is
Section~\ref{subsec:matched_beta_results}.

Read in one direction the comparison favours rescaling; read in the other it favours
retraining. At the strongest correction the Direct path attains, post-hoc rescaling is the more
accurate of the two on the price scale. At near-neutral correction, which only the Surrogate among
the retrained families approaches, the retrained fit is dramatically the more accurate: post-hoc
rescaling has by then given up more than a tenth of price-scale $R^2_P$. The frozen report's own
summary is that the accuracy comparison appears to cross over with correction strength, and that is
the honest reading. It is also why no superiority claim for either family appears anywhere in this
paper. What does hold in both regimes is a mechanism difference: at comparable first-order
correction the retrained configurations carry the lower correlation-ratio nonlinearity gap.

\paragraph{What a one-parameter rescaling costs, and what it cannot do.}
Driving the development coordinate to first-order neutrality post-hoc requires
$b=$ 1.160458, and Table~\ref{tab:centered_spread} records what that costs. On the held-out
block $R^2_P$ falls from 0.8928 to 0.8063 and mean absolute error rises from
\$75,976 to \$87,182; on the 2025 forward block $R^2_P$ falls from
0.9039 to 0.7991. The assessor-facing measures do not merely flatten, they
overshoot: held-out PRD moves from 1.0690 to 0.9871 and MKI from 0.9230 to
1.0726; VEI --- a May-2026 Exposure Draft diagnostic, proposed rather than adopted ---
from -26.17\% to $+$14.77\%, while COD worsens from 21.60 to
22.59. And the correction does not reach the nonlinear structure: the correlation-ratio
nonlinearity gap \emph{rises} from 0.1164 to 0.1343 while distance correlation
falls from 0.3822 to 0.2518 and then turns back up to 0.2681 at
$b_{\max}$. An affine map of a fixed predictor cannot change non-affine conditional-mean structure
except incidentally, and what it does incidentally here is make it larger. So the comparator
reproduces or improves much of the first-order accuracy and equity frontier that the Direct path
reaches over the range Direct attains, and it does not reproduce the residual-structure behaviour.

The comparator is built on the custom-objective $\rho=0$ origin because that cell is the
common origin of both penalized families, so rescaling it holds the execution path fixed, and the
contrast is then between penalized retraining and a one-parameter affine rescaling of the same
predictor rather than between two implementations. The same map applied to \emph{Ordinary LightGBM (standard raw-label native)} is
reported as a secondary, practical variant --- could a practitioner obtain a similar trade-off by
post-processing the standard pipeline? --- and Panel~C of Table~\ref{tab:centered_spread} shows that
it tracks the primary comparator closely, so the answer is descriptively yes. Under the frozen
convention that secondary variant determines nothing: not the common support, not the choice of
comparison targets, and not the primary mechanistic test. \emph{Centered-label native L2
(initialization-aligned)} has no post-hoc path at all, consistent with its role as an
implementation-decomposition control.


{\begin{table}[!htbp]
\centering
\small
\setlength{\tabcolsep}{5pt}
\caption{Descriptive juxtaposition of the two most-corrected retrained configurations
with post-hoc centered rescalings of the \emph{Custom-objective $\rho=0$ origin} at comparable
held-out first-order correction. Held-out evaluation. This comparison is keyed on the held-out
coordinate and is descriptive; the comparison matched on the development coordinate is
Table~\ref{tab:matched_beta} and Section~\ref{subsec:matched_beta_results}. No configuration is
selected or recommended here.}
\label{tab:posthoc_crossover}
\begin{tabular}{@{}rllrrrr@{}}
\toprule
$\beta_{\log}$ & Family & Configuration & $R^2_P$ & MAE & $\Delta_{\mathrm{NL}}$ & dCor \\
\midrule
-0.07941 & Direct & $\rho=100$ & 0.86868 & \$84,918 & 0.11535 & 0.26493 \\
-0.07905 & C-posthoc & $b=1.080148$ & 0.89242 & \$74,930 & 0.12946 & 0.28404 \\
\addlinespace[3pt]
+0.00084 & Surrogate & $\rho=75.4312$ & 0.88908 & \$81,930 & 0.12434 & 0.26804 \\
+0.00118 & C-posthoc & $b=1.174235$ & 0.77962 & \$90,429 & 0.13438 & 0.25408 \\
\bottomrule
\end{tabular}
\vspace{1mm}
\begin{minipage}{\textwidth}
\small
\emph{Notes.}
$\beta_{\log}$ is the held-out first-order slope, which is what the two rows in each
block are matched on approximately; the post-hoc base value is the frozen-grid value whose held-out
slope is nearest the retrained configuration's, not a development-coordinate solve. The
development-coordinate solves appear in Table~\ref{tab:matched_beta} instead. Post-hoc rows are
rescalings of a cached prediction array through
$f_b(x)=\bar y_T+b(f_0(x)-\bar y_T)$ with no refitting. $\Delta_{\mathrm{NL}}$ and dCor carry no
assessor-facing reference range and are residual-structure diagnostics, so a lower value is
directional only. Values are read from the frozen Stage-2 comparator record and the frozen
reproduction audit.
\end{minipage}
\end{table}
}


\subsection{Matched First-Order Correction: Direct, Surrogate and Post-Hoc Rescaling}
\label{subsec:matched_beta_results}

The comparison in Section~\ref{subsec:posthoc_comparator} is keyed on an outcome
sample. The comparison that can be designed rather than observed matches configurations on the
\emph{development} first-order coordinate and then reads their out-of-time behaviour. This
subsection reports it for the three primary families --- Direct, Surrogate and post-hoc rescaling of
the custom-objective $\rho=0$ origin --- with the corresponding rescaling of the workflow benchmark
carried as a secondary practical variant only.

\paragraph{The design, and the tolerance it was frozen with.}
Matching uses the primary development coordinate D1, the equal-weight mean of the seven
chronological fold values; the one-sale-one-vote D3 coordinate enters below as the sensitivity. The
three families share the development range from -0.138271 to -0.078651 whose upper endpoint is
set by Direct, the binding family. Six comparison targets are placed in that range by a
deterministic rule: probe equally spaced points, then take as the target the \emph{achieved}
development $\beta_{\log}$ of the fitted Direct configuration nearest the probe. Because the target
is an achieved value rather than the probe, Direct matches exactly by construction and no selection
discretion enters. Table~\ref{tab:matched_beta} lists the six targets and the configuration each
family used; every one of them is an actually fitted model or an exactly solved rescaling, and no
targeted new fit was required. The materiality and matching tolerance is
$\tau=$ 0.002. It was frozen in the matched-configuration record before any held-out or 2025
metric was read. That ordering is a strength of the design rather than a caveat, and it is what makes both the
attainment adjudication below and the materiality classification of the D3 sensitivity credible: the
worst matching gap over the twenty-four matched cells is 1.0569e-03 and lies inside $\tau$. The
$j=0$ row is an identity check rather than a result --- Direct at $\rho=0$, Surrogate at $\rho=0$
and the post-hoc map at $b=1$ are the same object --- and the largest relative disagreement anywhere
in that check, over all seventeen metrics and four evaluations, is 4.42e-12.

One qualification governs every reading of the outcome columns. The configurations are
matched on the \emph{development} coordinate, and nothing in the design equalises their held-out or
forward $\beta_{\log}$; matching a development statistic does not transport to an evaluation block.
The out-of-time differences below are therefore differences between configurations chosen to
approximately match the frozen development first-order coordinate, not effects measured at exactly
equal out-of-sample first-order correction.

\paragraph{At the strongest common target the retrained Direct path is the worst of the
three.}
Table~\ref{tab:matched_beta_headline} reports the held-out outcome at $j=5$, the
strongest correction all three families attain. Direct needs $\rho\approx87$ to reach it, and at
that penalty the fit has degraded: its $R^2_P$ is below both comparators, its mean absolute error is
several thousand dollars higher, its $\operatorname{RMSE}_{\log P}$ is the largest of the three, and
its COD is more than two points higher. That ordering is not the one the anchor tables of
Section~\ref{subsec:path_results} suggest, and it is available only because the comparison is made
at a matched first-order coordinate. Across the whole attainable range
(Table~\ref{tab:matched_beta}) the post-hoc comparator reproduces or improves much of the Direct
first-order frontier, so the Direct family's advantage over a one-parameter rescaling is not what
carries the paper. What does separate the families is residual structure: at every matched target
with a nonzero correction ($j=1,\dots,5$) the Surrogate carries a markedly lower correlation-ratio
nonlinearity gap than either comparator, and it is the only family whose gap falls below the origin
value. At $j=0$ the three families are the same object, so that row is an identity check rather
than a comparison.

\paragraph{Price-level and log-scale accuracy disagree, and both are reported.}
The extended targets in Table~\ref{tab:matched_beta_ext} make that disagreement
explicit. At the development target $-0.03$ the Surrogate holds held-out $R^2_P$ at
0.8917 while the post-hoc comparator has fallen to 0.8569 --- a large margin on
the price scale --- yet in $\operatorname{RMSE}_{\log P}$ the ranking reverses, 0.3443
against 0.2999. The same reversal is present at $-0.06$: 0.8934 against
0.8868 on the price scale, 0.3230 against 0.2944 in logs. It appears
again at the strongest common target, in a sharper form: there the price-scale difference between
the Surrogate and the post-hoc comparator is -0.00035 --- below its own frozen materiality
floor of 0.0005 and therefore a tie --- while the log-scale difference is 0.0174 and
sits well above the same floor while pointing the other way. The two metrics weight the price
distribution differently --- price-scale $R^2_P$ is dominated by squared errors on high-priced
sales, while $\operatorname{RMSE}_{\log P}$ weights all sales equally in logs --- so they can order
the same pair of configurations in opposite directions. The Surrogate's $1+\rho c_i^2$ weighting
up-weights sales far from the centre of the log-price distribution in \emph{either} direction, and
the frozen record does not decompose which part of the distribution drives the divergence, so the
pattern is reported as measured rather than attributed to one tail. Neither settles
``predictive performance'' on its own, and this paper does not let that phrase collapse onto one
metric.

\paragraph{Non-attainment over the evaluated path.}
Three of the extended targets lie outside the three-way common support, and the frozen
record carries an explicit attainment state for every family at each of them rather than an
interpolated number. Over the evaluated frozen path the Direct family did not attain
$\beta_{\log}\in\{-0.06,-0.03,0\}$, and the Surrogate family did not attain $0$ over its evaluated
path; the largest development corrections they reached were -0.078651 and
-0.018609 respectively. Both post-hoc families attain all three exactly, because a
one-parameter rescaling can be solved for any first-order target. Those four display states are
reported with their metric columns blank, never dropped and never filled. They are statements about
the penalty grid this experiment evaluated --- eighty-two positive values per family --- and not
about the objectives or about $\rho$ values the experiment never examined. The informative content
of the asymmetry is what the frozen record says it is: reaching a target is not by itself an
advantage, since driving the post-hoc coordinate to neutrality costs the accuracy recorded in
Section~\ref{subsec:posthoc_comparator}.

\paragraph{The one-sale-one-vote sensitivity moves a material subset of the
comparison.}
Because the fold-6 and fold-7 validation blocks overlap, the primary D1 coordinate
carries an observation weighting that was never chosen deliberately, and the frozen design requires
the row-balanced D3 coordinate as a sensitivity. It is not a small perturbation. The largest
disagreement between the two coordinates is 0.004867 in $\beta_{\log}$ units, larger than
$\tau$; the common support shifts to the range from -0.137744 to -0.076705; and re-running the
identical deterministic construction on the 130,165 distinct development sales reselects
13 of the 24 matched configurations --- eleven of them exact re-solves of the one-parameter
post-hoc map, which must move when the coordinate does, and two different fitted penalty values,
both at $j=2$. Of 315 headline sign comparisons, 259 agree in sign, 41 fall below the materiality
floor on both coordinates, 5 are structural zeros forced by the D1 match, and 10 are material sign
flips; raw sign agreement across all 315 is 0.9365. Rank order is not
preserved between the two coordinates. That reordering should not be read as the overlap scrambling
the paths: the $\beta_{\log}$ paths of both retrained families are non-monotone in $\rho$ under
\emph{both} coordinates, while both post-hoc families are monotone under both, so a shift of a few
thousandths reorders configurations that were already close to tied. The subsection's high-level
conclusions survive the change of coordinate: no $\Delta_{\mathrm{NL}}$ comparison flips sign under
D3 anywhere, and no flip touches the Direct degradation at $j=5$. What does not survive is the fine
ordering of near-ties. Ten pairwise comparisons cross a frozen floor and reverse sign, all of them
in small differences --- the largest is held-out mean absolute error at $j=2$, which moves from
$+$\$15 to $-$\$149 against an error of about \$76,000 --- and one of them is the $j=5$
price-scale tie reported above: on D3 the Surrogate-minus-post-hoc difference is $+$0.00055 rather
than -0.00035, above the floor and pointing the other way. That particular comparison is therefore
coordinate-sensitive and is not read as a result. The seven chronological folds are overlapping
windows rather than independent
replicates, so no statement here is a significance statement, and no fold spread is read as a
standard error.

Taken together, the three families are not interchangeable at matched first-order
correction, and no family is better on every diagnostic. Post-hoc rescaling reproduces much of the
Direct frontier and reaches first-order targets Direct did not reach on this path, at a large
accuracy cost when pushed that far. The Surrogate preserves price-scale accuracy and reduces
nonlinear residual structure where the comparator increases it, and is worse in logs. That is the
result: a diagnostic distinction between mechanisms, not a ranking, and no penalty strength is
selected anywhere. The price-scale half of that comparison is stated under this paper's
direct-exponentiation convention: Appendix~\ref{app:smearing} reports a frozen retransformation
sensitivity under which the Surrogate--post-hoc ordering on $R^2_P$ and MAE reverses, while every
vertical-equity, uniformity, and residual-structure diagnostic in this table is unchanged.

{\begin{table}[!htbp]
\centering
\small
\setlength{\tabcolsep}{5pt}
\caption{Held-out outcome at the strongest development first-order correction all three
primary families attain ($j=5$ of Table~\ref{tab:matched_beta}). Configurations are matched on the
\emph{development} coordinate D1; their held-out first-order slopes are not equalised by that match,
so these are differences between approximately matched configurations rather than effects at exactly
equal out-of-sample correction. No configuration is selected or recommended.}
\label{tab:matched_beta_headline}
\begin{tabular}{@{}llrrrrrr@{}}
\toprule
Family & Configuration & $R^2_P$ & MAE & $\operatorname{RMSE}_{\log P}$ & COD
 & $\Delta_{\mathrm{NL}}$ & dCor \\
\midrule
Direct & $\rho=86.8511$ & 0.8760 & \$82,764 & 0.3170 & 23.97 & 0.1198 & 0.2730 \\
Surrogate & $\rho=2.5595$ & 0.8963 & \$76,128 & 0.3094 & 21.40 & 0.0914 & 0.2714 \\
C-posthoc & $b=1.069186$ & 0.8966 & \$74,236 & 0.2920 & 21.61 & 0.1282 & 0.2944 \\
\bottomrule
\end{tabular}
\vspace{1mm}
\begin{minipage}{\textwidth}
\small
\emph{Notes.}
C-posthoc is the centered rescaling of the \emph{Custom-objective $\rho=0$ origin}, the
primary post-hoc family because that cell is the common origin of both penalized paths. Direct and
Surrogate penalty magnitudes are not numerically comparable across families. COD is reported in
percent and lies outside the adopted $[5,15]$ range for every row, as it does at the path origin, so
no compliance claim is implied. $\Delta_{\mathrm{NL}}$ and dCor are residual-structure diagnostics
with no assessor-facing reference range.
\end{minipage}
\end{table}
}


\subsection{Chronological Stability and Other Model-Quality Dimensions}
\label{subsec:path_stability_results}

A regularization path is more informative if its direction is stable across market
periods and if improvements in vertical-equity diagnostics do not conceal large
changes in valuation level, horizontal uniformity, or alternative prediction
metrics.
We therefore reserve the complete fold-level paths for prediction, valuation level and horizontal uniformity, vertical equity, and mechanism diagnostics for Appendix~\ref{app:ccao_path_details}, while reporting the main path findings in the figures above.

The seven chronological folds are overlapping expanding-window constructions rather
than independent replications. Their validation blocks are not mutually disjoint---the
sixth and seventh blocks share observations---so the equal-weight cross-validation mean
is an average over partly overlapping evaluation evidence. Every prediction is still
genuinely out of training sample, but across-fold dispersion is descriptive
chronological-window variation and not an IID standard error, and it is never reported
as one. On this evidence the folds support the direction of the targeted first-order
mechanism while they do not identify a stable multi-metric operating point: the
penalty scales at which individual diagnostics turn are not common across
metrics, and no materiality threshold is imposed on the fold-level path comparisons
themselves, so their differences are not labelled material or negligible.

Temporal portability is weaker than within-development reproducibility, and the two
out-of-time blocks do not play the same role. The held-out block is the later
development-period evaluation used in the primary analysis; the separate 2025 forward
sample is a stronger portability check on a subsequent market period. They do not
behave identically. Along the frozen Surrogate path the held-out first-order slope
attains a maximum marginally above zero, whereas the cross-validation mean and the 2025
forward path remain negative throughout, so the same correction does not leave the same
first-order signature on every evaluation period. A disagreement of this kind is
evidence about temporal sensitivity rather than a defect to be tuned away. Accordingly,
the appendix path figures trace the tested grids descriptively and do not mark a
selected, portable, or recommended operating interval.

Three alternative temporal designs were evaluated against the primary design without
refitting it. A strict-date design moves every boundary-date sale to the evaluation
side, so that training strictly precedes evaluation at each origin; an oracle
overlap-removal diagnostic deletes from each training block every parcel that reappears
in the paired evaluation block; and an unseen-parcel view restricts the frozen
predictions to evaluation rows whose parcel never appears in the corresponding training
block. The screening gate returned \texttt{ALERT} on a single criterion, the first-order
slope's sign and family ordering. Targeted local refinement of the affected regions of
the frozen grid returned \texttt{NOT\_CONFIRMED}, recording the alert as a near-tie and
near-zero artifact that did not survive refinement, with overall temporal status
\texttt{PASS\_PRIMARY\_STANDS}. The separation tolerance against which that comparison
judged materiality is $\tau=$ 0.002. It was not a constant chosen at decision time: it is
the first-order matching tolerance already frozen for the matched-correction analysis,
adopted here before this comparison read any outcome.

These designs sharpen rather than soften the reading. Under the strict-date design the
first-order path keeps its shape, sign, and family ordering wherever the two families
are separated by more than that tolerance, and the 2025 forward numbers are unchanged
because that boundary was already date-separated---a property of the design rather than
independent confirmation of portability. The oracle diagnostic could not be implemented
prospectively, since constructing its training blocks requires knowing which parcels
appear in a future evaluation block; it is therefore reported as a diagnostic and not as
a proposed sample-construction rule. It leaves the within-path conclusions in place, and
it makes the uncorrected origin measurably more regressive, which strengthens rather
than weakens the motivation for the correction. The unseen-parcel view refits nothing
and changes the denominators, so it is reported on its own base and is never differenced
against the frozen path; on it the Surrogate continues to attain substantially more
first-order correction than the Direct family at strong penalties. Direct and Surrogate
are accordingly not interchangeable, and their attainable spans and temporal behaviour
are reported separately rather than averaged together.

Taken together, the corrections move the targeted first-order mechanism on every
evaluation block, but they do not do so at a penalty scale that is stable across market
periods, and the nonlinear-residual and broader-dependence diagnostics are quoted along these paths only at the zero-penalty origin, so this section makes no claim about where along the path they turn. The penalty grid was frozen
in advance, no diagnostic reported here defined it or any selection rule, and no penalty
strength is selected. The contribution of this section is a controlled characterization
of the attainable paths together with a temporal failure-mode audit, rather than the
identification of a deployment penalty or a portable operating interval.


\subsection{External-Validity Boundary}
\label{subsec:external_results_boundary}

The CCAO paths establish behavior only for the stated research translation, sale sample, temporal design, and fixed LightGBM configuration. They do not establish that the same raw penalty or accuracy--equity path transfers to another jurisdiction. The repository's newer normalized external benchmark is therefore treated as a separate robustness project rather than as confirmatory evidence in this manuscript. In particular, official jurisdiction-level assessment-ratio profiles must not be conflated with valuation ratios generated by the common research AVM.


\section{Discussion and Conclusions}
\label{sec:discussion}

The CCAO application answers a bounded model-design question. Along the currently frozen custom-objective paths, covariance-guided regularization can materially attenuate the regressive first-order pattern among observed sales, and moderate configurations can do so before predictive performance deteriorates. Stronger correction can alter valuation level, ratio dispersion, and nonlinear ratio structure. These are fixed-base-learner path results: the 994-tree LightGBM configuration is held constant to isolate objective behavior, so the experiment does not estimate each objective family's fully retuned attainable frontier. In addition, a comparison against ordinary LightGBM carries a native-versus-custom execution-path difference as well as the penalty, which is why the within-family custom-objective $\rho=0$ origin, rather than the workflow benchmark, is the penalty-isolating reference; Appendix~\ref{subsec:zero_control_results} reports the completed parity audit that located that difference. The evidence therefore supports a practical control for one component of a sale-price-related ratio pattern, not a complete solution to vertical equity or a deployment model.


\paragraph{Predictive accuracy alone does not remove the price-related ratio trend.}
In this application the two coexist. The ordinary workflow benchmark, \emph{Ordinary
LightGBM (standard raw-label native)}, attains high price-scale predictive accuracy on both
out-of-time evaluations while its first-order log-ratio/log-price slope and its assessor-facing
vertical-equity diagnostics remain well away from their neutral values
(Section~\ref{subsec:baseline_tension}). Stronger predictive accuracy therefore does not by itself
remove the first-order assessment-ratio trend, which is what makes the mechanism worth targeting
directly. That statement is bounded by the application: it describes CCAO's residential sale
sample, its feature information, its temporal design, and one fixed base-learner configuration. It
is not a general result that machine-learning valuation is inherently regressive. No alternative
model class is re-estimated for comparison, no linear-regression quantitative result is reported
anywhere in this paper, and nothing in the evidence attributes the trend to the learner rather than
to the data, the prediction target, or the workflow.

\paragraph{The target is one global, first-order, price-related mechanism.}
Both objectives control a single scalar: the covariance $\Cov(e,y)$ between the log
valuation-ratio residual and log sale price in~\eqref{eq:training_covariance}, equivalently the
first-order log-ratio/log-price trend summarized by $\beta_{\log}$ in~\eqref{eq:log_ratio_slope}.
What that quantity is, and what it is not, governs every conclusion that follows. It is a mechanism
diagnostic; it is global over the evaluated sample; it is first order in log price; and it is
specific to a price-related association. It is not fairness in general, not individual fairness, not
local or subgroup equity, not causal fairness, not statistical independence of residuals and price,
and not fairness of the final tax burden. In many settings, richer information can improve both
predictive accuracy and assessor-facing equity measures \cite{SmithEtAl2026}; no experiment here
combines that route with these penalties, so nothing below speaks to how the two interact. What the
penalties target is a countywide first-order pattern among observed sales, and they do not determine
valuation level, guarantee assessor-facing reference ranges, remove nonlinear or local inequity,
establish residual--price independence, or identify effects on final tax liabilities.
Section~\ref{subsec:correction_scope} states the corresponding
denials for the objectives themselves, and every result below is read inside that scope.

\paragraph{A one-parameter post-hoc rescaling reproduces much of the first-order
movement.}
The closest post-processing alternative is not a straw comparator and is not treated as
one. It is the exact fixed-prediction-space form of the Direct path
(Corollary~\ref{cor:path_scaling}): a centered rescaling of a cached prediction array, solved rather
than fitted, refitting nothing. Over the first-order range the Direct family attains on the
evaluated grid, that comparator reproduces or improves much of the accuracy and first-order equity
frontier Direct reaches, and it reaches first-order targets the evaluated Direct grid did not reach
at all (Sections~\ref{subsec:posthoc_comparator} and~\ref{subsec:matched_beta_results}). Two
consequences follow, and this paper accepts both. First, the paper does not claim that training-time
regularization is necessary in order to move the first-order mechanism. Second, the Direct objective
nonetheless remains a trainable mechanism for moving along that dimension --- the penalty is what
moves it, read within family from the shared custom-objective $\rho=0$ origin rather than against
the workflow benchmark --- but the first-order correction itself is not unique to retraining.

\paragraph{Equalizing the first-order trend does not equalize the residual geometry.}
The three families are nevertheless not interchangeable, and the comparison that shows
it is the one made at matched first-order correction rather than at matched penalty. Matching uses
the primary development coordinate D1, with the row-balanced one-sale-one-vote coordinate D3 as the
declared sensitivity and the zero-correction row an identity check rather than a comparison. At
every matched target with a nonzero correction the Surrogate carries a markedly lower
correlation-ratio nonlinearity gap than either comparator and is the only family whose gap falls
below the origin value, whereas driving a fixed predictor to first-order neutrality by affine
rescaling \emph{raises} that gap. At the strongest correction all three families attain, the
retrained Direct configuration is the worst of the three on price-scale accuracy, on log-scale
accuracy, and on dispersion --- an ordering visible only because the comparison is made at a matched
first-order coordinate. The D3 sensitivity reorders a material subset of the near-ties, but it flips
no nonlinearity-gap comparison anywhere and does not touch the Direct degradation at the strongest
common target, and where a retrained family did not attain an extended target on its evaluated
path that display state is reported as a state rather than interpolated. The general statement is the one this evidence
supports: equalizing the first-order price-related trend does not equalize the full residual
geometry, and neither objective directly controls the non-affine conditional-mean structure that
$\Delta_{\mathrm{NL}}$ measures. These are diagnostic distinctions between mechanisms and not a
ranking of them, and they are quoted only at the coordinates where the frozen record reports them
--- the zero-penalty origin, the matched first-order and extended targets, the descriptive crossover
juxtaposition, and the centered-spread comparator path --- and never as a shape claim along the
whole positive-$\rho$ paths. Global first-order neutrality is therefore informative but not
sufficient evidence of a flat value-related profile, of broader residual--price independence, or of
locally uniform vertical equity.

The fixed-prediction-space analysis sharpens this interpretation rather than creating a separate theoretical claim. The Direct objective has rank-one geometry and, with an intercept in the fixed space, reduces exactly to a centered spread correction; this makes transparent both what the covariance penalty controls and how limited that control can be. The Surrogate, by contrast, is a log-price-distance-weighted projection whose correction can combine multiple learnable modes and therefore need not trace the same path as Direct. This comparison explains why the Surrogate should not be interpreted as a computational approximation to the Direct objective alone. At the same time, neither fixed-space result predicts the exact path of a retrained tree ensemble, and the spectral representation does not by itself explain the empirical high-$\rho$ ratio shape.

\paragraph{Part of the price-scale comparison depends on the retransformation
convention.}
Predictions are returned to the price scale by direct exponentiation, which is an
operational convention rather than a conditional-mean correction, and one half of the
matched-correction comparison turns on it. Under the frozen Duan sensitivity the estimated smearing
factor is not common across the arms compared at the selected matched first-order illustration, and
the Surrogate-versus-post-hoc ordering on price-scale $R^2_P$ and on mean absolute error reverses
(Appendix~\ref{app:smearing}). That ordering is accordingly a property of the retransformation
convention at that coordinate, and it is not evidence of general predictive superiority for either
arm. This should not be over-read into a claim that every metric ranking in the paper is unstable.
COD, COV, PRD, PRB, MKI, VEI, $\beta_{\log}$, $\Delta_{\mathrm{NL}}$ and distance correlation are
invariant to a common multiplicative factor as a matter of algebra, so no vertical-equity,
uniformity, or residual-structure conclusion here depends on the choice, while the valuation-level
ratios scale exactly by the factor. $\operatorname{RMSE}_{\log P}$ is defined on the log scale and
is not a price-scale rescaling quantity. Three groups of quantities should therefore be kept apart
whenever these comparisons are quoted: convention-sensitive price-scale quantities, log-scale
quantities, and exactly invariant ratio and dependence quantities.


\paragraph{The correction-strength relationship does not transport across market
periods.}
The regularization paths describe tradeoffs over the tested grids, and that is the whole
of what they describe. The chronological folds are overlapping expanding-window constructions rather
than independent replications --- the sixth and seventh validation blocks share observations, so
across-fold dispersion is chronological-window variation and is never read as an IID standard error
--- and on that evidence they support the direction of the targeted mechanism without identifying a
stable multi-metric operating point. The two out-of-time blocks keep distinct roles: the held-out
block is the later development-period evaluation used in the primary analysis, and the separate 2025
forward sample is a stronger portability check on a subsequent market period. They do not agree, and
that disagreement is itself the finding rather than a defect to be tuned away. Development-period
path landmarks did not transfer to those blocks, so this study identifies no candidate region, no
portable interval of $\rho$, no activity-onset boundary, and no upper guardrail: the evidence does
not support any such construct, and none is used here. No $\rho$ is selected, and no deployment
penalty is chosen anywhere in this paper. The correct reading is not that the paper failed to find
the right $\rho$. It is that the relationship between correction strength and path position varies
enough across market periods that the evidence does not support a single portable operating region
or a universal $\rho$. That is an empirical failure-mode result about the method rather than
unfinished parameter tuning, and the design is what makes it one: the penalty grid was frozen in
advance, no diagnostic reported here defined it or any selection rule, and the materiality tolerance
$\tau$ was the first-order matching tolerance already frozen for the matched-correction analysis
before the relevant outcome comparison was read.

\paragraph{What the study establishes, and where it stops.}
The value of the approach is operational. It converts a ratio-study concern into a model-training option that can be incorporated into a standard boosted-tree workflow: the Surrogate is exactly observation-additive, while the Direct implementation uses the exact covariance gradient with diagonal curvature through the standard LightGBM interface. Both can be evaluated through the same chronological-validation and ratio-study workflow used for the baseline model. Those features follow directly from the CCAO problem that motivated the collaboration.

Taken together, the contribution is a training mechanism for assessor workflows
rather than a deployment rule, and it is built on a classical assessor-side diagnostic rather
than on a new one. The paper establishes how to target that one specific
first-order mechanism through a standard boosted-tree interface; how the Direct and Surrogate
objectives move along the resulting tradeoff, read within family from the shared custom-objective
$\rho=0$ origin; how a simple one-parameter post-hoc rescaling behaves beside them; what a
comparison at matched first-order correction reveals about residual structure that a comparison at
matched penalty does not; and how far temporal instability limits portability. It does not establish
a universal fairness solution, a universal penalty, or a production-ready rule for selecting penalty
strength, and it makes no external-validity claim from the superseded multi-county analysis
(Section~\ref{subsec:external_results_boundary}). Covariance-guided regularization should therefore
be treated as one auditable modeling option inside a broader valuation-review process, and not as
evidence that training-time regularization is necessary or sufficient for vertical equity.

\subsection{\texorpdfstring{
Practical Implementation Considerations}{Practical Implementation Considerations}}
\label{sec:implementation_workflow}

For a later operational implementation, the correction can be incorporated into an existing AVM development workflow through the following sequence. The final step is a prospective decision this study does not make.

\begin{enumerate}[leftmargin=7mm]
    \item Build a defensible baseline model using verified sales, appropriate time adjustment, relevant property features, and local market information, fitted as the office would ordinarily run it. This is the workflow benchmark against which the rest is described.
    \item Reserve future observations for held-out and forward evaluation and create rolling or blocked validation folds that respect sale chronology.
    \item Diagnose the first-order mechanism on that benchmark before changing the objective. Report baseline prediction and valuation-level results. Add COD, PRD, PRB, MKI, VEI, value-group ratio profiles, $\beta_{\log}$, $\Delta_{\mathrm{NL}}$, and dCor.
    \item Fit the Direct and Surrogate covariance-guided objectives over prespecified logarithmic grids of $\rho$, including $\rho=0$ as an implementation control and as the within-family origin against which penalty effects are read.
    \item Summarize the complete regularization paths across chronological validation folds, including fold-level variability, and report the paths themselves rather than a preferred segment of them.
    \item Compare the training-time configurations against the simple post-hoc alternatives on the same predictions --- constant level calibration and a validation-fitted centered spread map --- using identical temporal splits and diagnostics.
    \item Compare families at matched first-order correction rather than at matched penalty, and read predictive accuracy on both the price and log scales alongside the assessor-facing vertical-equity measures and the broader residual-structure diagnostics. State the retransformation convention under which any price-scale comparison is made.
    \item Refit the same fixed paths on the full development samples and evaluate them descriptively on the held-out and forward periods without changing the grid or base learner.
    \item Repeat the diagnostics by geography, property class, value range, and time before any later deployment-oriented model-selection decision.
    \item If deployment eventually requires a single operating point, define that decision prospectively instead of reading it off the paths above: fix the predictive, valuation-level, and equity criteria, their acceptance thresholds, and the tie-breaking rule in advance, then validate that rule out of time on periods not used to formulate it. The retrospective path characterization reported here does not supply such a rule and is not a substitute for one.
\end{enumerate}

\begin{table}[t]
\centering
\caption{Different pipeline components address different problems.}
\label{tab:roles}
\begin{tabular}{p{4.2cm}p{10.0cm}}
\toprule
Component & Primary role \\
\midrule
Better property, spatial, and temporal data & Reduce omitted structure and improve the conditional valuation model; may improve accuracy and equity together \\
Constant level calibration & Move the overall ratio level toward its required target; does not change log-residual covariance \\
Covariance-guided regularization & Discourage a global first-order log-ratio/log-price trend during model training \\
Residual-structure and local diagnostics & Use $\beta_{\log}$, $\Delta_{\mathrm{NL}}$, distance correlation, and ratio-shape diagnostics to distinguish first-order, non-affine conditional-mean, and broader dependence patterns; use stratified diagnostics to detect geographic or property-class patterns hidden by countywide summaries. \\
Temporal validation and uncertainty analysis & Characterize how path behavior changes across market periods; by itself, chronological validation does not establish a temporally invariant penalty range \\
\bottomrule
\end{tabular}
\end{table}

The office should save the penalty definition, full grid, validation predictions, held-out and forward predictions, and all ratio-study outputs. This documentation makes the change auditable and allows the correction to be re-evaluated in later assessment cycles.

\subsection{Limitations}


\emph{The target and the proxy.} Sale price is simultaneously the learning target, the
benchmark against which valuations are judged, and the denominator of the assessment ratio itself,
so the correction is defined on a quantity that appears on both sides of the diagnostic. Sale price
is also an imperfect proxy for latent market value. Equation~\eqref{eq:bayes_residual_covariance}
shows that residual--sale-price covariance carries a mechanically negative component from
conditionally unpredictable transaction variation, so forcing it toward zero may spread predictions
toward realized sale noise rather than recover latent value. For the same reasons $\beta_{\log}$ is
used here as a first-order mechanism diagnostic and not as a complete definition of vertical
equity.

\emph{What the objective does not control.} The objective controls a global first-order
dependence between log residuals and log price. Nonlinear conditional-mean structure and broader
statistical dependence can remain, and neither Direct nor Surrogate optimizes PRD, PRB, VEI,
valuation level, $\Delta_{\mathrm{NL}}$, or local equity, so a countywide first-order improvement
can coexist with nonlinear, geographic, class-specific, temporal, or level inequity. The Surrogate's
$1+\rho c_i^2$ weighting is symmetric in log-price distance from the centre of the log-price
distribution and is not a protection of any particular part of that distribution. Price-scale
comparisons carry one further dependence: predictions are returned to the price scale by direct
exponentiation, which is an operational retransformation rather than a smearing or conditional-mean
correction, and Appendix~\ref{app:smearing} reports a frozen sensitivity under which some
price-scale orderings reverse while every ratio, uniformity, and residual-structure quantity is
algebraically unchanged. Finally, the closest post-processing comparator --- the centered spread map
implied by the Direct fixed-space result --- is included rather than omitted, and it reproduces much
of the Direct first-order movement over the range Direct attains, so this paper claims neither
necessity nor general superiority for training-time regularization.

\emph{Temporal evidence and its limits.} The seven rolling-origin folds are overlapping
expanding-window constructions whose sixth and seventh validation blocks share observations. They
are not independent experimental replications, and fold-level agreement among them is not a
confidence interval. Across those folds, the held-out block, and the 2025 forward block the
correction-strength relationship changes enough that the evidence identifies no portable universal
$\rho$ and no single operating region, so raw-$\rho$ landmarks may not be treated as portable.
Three alternative temporal designs were evaluated against the primary design rather than left
outstanding, and their outcomes are reported as recorded in
Section~\ref{subsec:path_stability_results}. Apart from selected group-profile intervals, the paper
reports descriptive point estimates and fold variation rather than inferential uncertainty for path
differences.

\emph{Scope, and the research/production boundary.} The evidence is centered on one
application. The primary analysis is a Python research translation of CCAO's published residential
workflow over its verified sale sample, not an official production run or assessment decision, and
it characterizes fixed-994-tree regularization paths rather than each objective family's fully
retuned frontier, because the base learner is held constant to isolate objective behavior.
Outcomes on verified transacting properties need not transport to the full roll when sale
propensity, verification, property condition, or unobserved quality differs between sold and unsold
parcels, so the paper estimates model behavior on the observed sale distribution rather than
jurisdiction-wide tax incidence or inequity relative to latent market value. External validity must
be established elsewhere, and the active paper makes no external-validity claim from the superseded
multi-county analysis. Recent Cook County evidence also documents that the degree of assessment
regressivity differs substantially across property classes \cite{CaiWiley2026}, so a countywide
result pooled over the residential classes in this sample should not be read as describing each of
them. Research evaluation is also not deployment: no operational penalty strength
is selected here, and the decisions a deployment would require --- an explicit selection rule, its
acceptance criteria, and out-of-time validation of that rule --- lie outside what this study
establishes.

\subsection{Future Work}

The next evaluation step is a clean ablation that compares the baseline model, better spatial and temporal information, the covariance correction, and their combination under the same temporal splits. Comparable-property and neighborhood features must use only information available before the target sale.
For any future deployment-oriented calibration, the selection objective and its acceptance criteria should be frozen before evaluating later assessment cycles, and temporal stability should be reported across market periods rather than inferred from a single development-period landmark.


One natural methodological extension is structured covariance control by geography, property class, and time, with safeguards for small or unstable groups and separate regularization strengths for the dimensions being targeted. The corresponding open question is no longer whether a \emph{global} first-order correction can be obtained without retraining: Sections~\ref{subsec:posthoc_comparator} and~\ref{subsec:matched_beta_results} already evaluate a global one-parameter centered-spread post-hoc correction and show that it reproduces much of Direct's first-order movement over the range Direct attains. What remains open is whether \emph{structured} corrections along those same dimensions are better performed post hoc, during training, or jointly. A follow-up study should compare structured post-hoc corrections with structured training-time penalties under the same temporal design and the same predictive, assessor-facing, nonlinear-shape, and local or subgroup diagnostics. The present paper draws no compatibility or superiority conclusion between the two families. Further work should also add block-bootstrap uncertainty, worst-period validation, and evaluation on the full assessment inventory. Desk review by assessors is needed to determine whether the changed predictions remain understandable and operationally defensible.


\printbibliography


\appendix

\section{Metric Details and Guidance Status}
\label{app:metrics}

\subsection{Assessor-facing guidance status}

The 2013 IAAO Standard on Ratio Studies is the established standard cited for PRD
and COD guidance in this paper \cite{IAAO2013Ratio}. The May 2026 revision of the
ratio-studies document is an exposure draft. Measures introduced or revised there,
including MKI and VEI, are reported as descriptive assessor-facing diagnostics rather than as
final adopted compliance standards \cite{IAAO2026ExposureRatio}. Reference bands
are used as descriptive guidance in the present
regularization-path analysis; the paper does not
infer legal compliance from a point estimate.

The weighted mean ratio satisfies
\[
\bar r_W
=
\frac{\sum_i r_iP_i}{\sum_iP_i}
=
\frac{\sum_i\widehat P_i}{\sum_iP_i},
\]
and PRD is $\bar r/\bar r_W$. PRB follows the IAAO median-normalized regression
definition. The implementation executed for this application is the metric code tracked in the
repository, and the transformations, trimming rules and proxy definition it applies are fixed in
the frozen conventions catalogued in Appendix~\ref{app:ccao_results}.

\subsection{Mechanism and residual-structure diagnostics}

The three diagnostics introduced after the assessor-facing measures play distinct
roles. The sample slope $\beta_{\log}(S)$ in Eq.~\eqref{eq:log_ratio_slope} is the
direct out-of-sample analogue of the covariance mechanism: on a fixed evaluation
sample it differs from empirical covariance only by the fixed factor
$\widehat{\Var}_S(y)^{-1}$.

For the nonlinear diagnostic, let
\[
Z=\frac{y-\E[y]}{\sqrt{\Var(y)}},
\qquad
m(z)=\E[e\mid Z=z],
\]
and let $\Pi_{\mathrm{aff}}m$ denote the $L^2(P_Z)$ projection of $m$ onto the
affine span of $\{1,Z\}$. Pearson's correlation ratio, also called the
nonparametric coefficient of determination, is
\[
\eta^2(e\mid Z)
=
\frac{\Var\!\left(m(Z)\right)}{\Var(e)}
=
\frac{\Var\!\left(\E[e\mid Z]\right)}{\Var(e)}
\]
\cite{DoksumSamarov1995}. Because $Z$ is standardized, the affine component explains
the fraction $\Corr(e,Z)^2$ of the variance of $e$. Orthogonal projection therefore
gives
\begin{equation}
\label{eq:nonlinearity_gap_projection}
\Delta_{\mathrm{NL}}
=
\eta^2(e\mid Z)-\Corr(e,Z)^2
=
\frac{
\E\!\left[
\left\{
m(Z)-(\Pi_{\mathrm{aff}}m)(Z)
\right\}^2
\right]
}{
\Var(e)
}
\geq 0.
\end{equation}
Thus $\Delta_{\mathrm{NL}}$ isolates the conditional-mean explanatory power that is
not captured by the best affine residual--price relationship. We use
\emph{correlation-ratio nonlinearity gap} as a descriptive name and
$\Delta_{\mathrm{NL}}$ as paper-specific shorthand; neither the name nor the symbol
is claimed as canonical notation from \textcite{DoksumSamarov1995}. That paper
provides the correlation-ratio/nonparametric-$R^2$ foundation and explicitly notes
its use in constructing measures of regression nonlinearity. The integrated
squared-difference goodness-of-fit framework of \textcite{HardleMammen1993} is
closely related to the projection representation above, but it is not the source of
the exact normalized statistic in Eq.~\eqref{eq:nonlinearity_gap_projection}.

For the empirical CCAO analysis, the finite-sample implementation is fixed before
the penalized paths are compared. Each held-out or forward evaluation sample $S$ is
assigned deterministically to five folds using sale identifiers. For each fold, we
fit on the complementary observations (i) an affine regression of $e$ on $Z$ and
(ii) a cubic $B$-spline regression of $e$ on $Z$ with eight quantile knots, and we
score both fits on the omitted fold. Let
$\operatorname{MSE}_{\mathrm{aff}}(S)$ and
$\operatorname{MSE}_{\mathrm{spl}}(S)$ denote the resulting pooled out-of-fold mean
squared errors. The reported sample statistic is
\begin{equation}
\label{eq:nonlinearity_gap_estimator}
\widehat{\Delta}_{\mathrm{NL}}(S)
=
\max\left\{
0,\,
\frac{
\operatorname{MSE}_{\mathrm{aff}}(S)
-
\operatorname{MSE}_{\mathrm{spl}}(S)
}{
\widehat{\Var}_S(e)
}
\right\}.
\end{equation}
The nonnegative truncation respects the population inequality in
Eq.~\eqref{eq:nonlinearity_gap_projection}; finite-sample cross-fitting can otherwise
produce a small negative difference when the flexible fit does not improve
out-of-fold squared error. For readability, the empirical tables use
$\Delta_{\mathrm{NL}}$ as shorthand for
$\widehat{\Delta}_{\mathrm{NL}}(S)$.
The estimator is reported on the primary held-out and 2025 forward paths. The same frozen cross-fitted estimator was later applied, fold by fold, to the already-stored chronological validation predictions; those CV values are a retrospective diagnostic and were not used to define the $\rho$ grid or any model-selection rule.

Finally, population distance correlation is zero if and only if independence under
the usual moment conditions \cite{SzekelyRizzoBakirov2007}. The paper uses its
sample analogue only as an omnibus diagnostic of residual--price dependence; unlike
$\beta_{\log}$ it is unsigned, and unlike $\Delta_{\mathrm{NL}}$ it is not restricted
to conditional-mean structure. It has no assessor-standard reference range.

Every reported value comes from one implementation: the
\texttt{distance\_correlation} routine of the \texttt{dcor} Python package, called on the
full evaluation sample with the package's default unit exponent and without bias
correction, so the executed estimator is the usual biased distance correlation computed from
double-centered distance matrices, rather than the $U$-centered bias-corrected
squared-distance-correlation estimator the package also exports but that this paper never calls.
The executed call passes the
signed log residual $e=\log\widehat P-\log P$ as its first argument and $y=\log P$ as its
second.


\section{Squared-Error Prediction and Residual--Value Dependence}
\label{app:population_identity}

Let $Y=\log P$, let $X$ be the observed feature vector, and let
\[
f^\star(X)=\E[Y\mid X]
\]
be the squared-error Bayes predictor. Define $e^\star=f^\star(X)-Y$. If $\E[Y^2]<\infty$, then
\begin{proposition}
\label{prop:bayes_covariance}
\[
\Cov(e^\star,Y)
=
-\E[\Var(Y\mid X)]
\le 0.
\]
\end{proposition}
\begin{proof}
Using the law of total variance,
\[
\Cov(e^\star,Y)
=
\Cov(\E[Y\mid X]-Y,Y)
=
\Var(\E[Y\mid X])-\Var(Y)
=
-\E[\Var(Y\mid X)].
\]
\end{proof}

This identity explains a source of negative log-residual dependence when important price variation remains unresolved after conditioning on $X$. It does not imply that every finite-sample model is regressive under every assessor metric. It also does not imply that adding features cannot improve both accuracy and equity.

\section{Covariance and PRD on the Price Scale}
\label{app:cov_prd}

Let $P>0$ and $r=\widehat P/P>0$. The population PRD can be written as
\[
\mathrm{PRD}
=
\frac{\E[r]\E[P]}{\E[rP]}.
\]
Then
\begin{proposition}
\label{prop:cov_prd}
\[
\operatorname{sign}\!\left(\Cov(r,P)\right)
=
\operatorname{sign}(1-\mathrm{PRD}).
\]
\end{proposition}
\begin{proof}
Since $\E[rP]>0$,
\[
1-\mathrm{PRD}
=
\frac{\E[rP]-\E[r]\E[P]}{\E[rP]}
=
\frac{\Cov(r,P)}{\E[rP]}.
\]
\end{proof}

The penalty in the paper acts on $\Cov(\log r,\log P)$, not on $\Cov(r,P)$. Proposition~\ref{prop:cov_prd} therefore gives directional motivation rather than an equality between the training target and PRD.

\section{\texorpdfstring{Fixed-Space Regularization Paths: Direct and Surrogate}{Fixed-Space Regularization Paths: Direct and Surrogate}}
\label{app:theory}


This appendix isolates the data-fit geometry of the two penalties when the prediction space is held fixed, omitting model-specific complexity penalties in order to expose the correction mechanisms cleanly. The Direct result provides an exact rank-one benchmark for the targeted covariance correction. The Surrogate result provides the corresponding log-price-distance-weighted projection benchmark and makes precise why its path need not remain one-dimensional. These finite-sample projection results are used for mechanism interpretation only; they do not assert that a fully retrained boosting path remains in a fixed prediction space or that the fixed-space formulas describe LightGBM's finite boosting trajectory exactly.

\subsection{Direct penalty: rank-one covariance path}

Let $\Vcal\subseteq\R^n$ be a closed linear prediction space, define
\[
\langle u,v\rangle_n=\frac{1}{n}u^\top v,
\qquad
\|u\|_n^2=\langle u,u\rangle_n,
\qquad
c=y-\bar y\1,
\]
and consider
\begin{equation}
\label{eq:path_problem}
f_\rho\in\argminop_{f\in\Vcal}
\left\{
\|f-y\|_n^2
+\frac{\rho}{2}\langle f-y,c\rangle_n^2
\right\}.
\end{equation}
Let $P_{\Vcal}$ be the orthogonal projection under $\langle\cdot,\cdot\rangle_n$ and set
\[
f_0=P_{\Vcal}y,\qquad
C_0=\langle f_0-y,c\rangle_n,\qquad
g=P_{\Vcal}c,\qquad
A=\|g\|_n^2.
\]

\begin{proposition}[Covariance shrinkage path]
\label{prop:path}
If $A>0$, then
\[
f_\rho=f_0-\frac{\rho}{2}C_\rho g,
\qquad
C_\rho=\frac{C_0}{1+\rho A/2}.
\]
Consequently, the slope of residuals on $y$ is multiplied by
\[
q(\rho)=\frac{1}{1+\rho A/2}.
\]
\end{proposition}

\begin{proof}
For every $h\in\Vcal$, first-order optimality gives
\[
2\langle f_\rho-y,h\rangle_n
+\rho C_\rho\langle c,h\rangle_n=0.
\]
Since $f_0=P_{\Vcal}y$, $\langle f_0-y,h\rangle_n=0$. Writing
$d_\rho=f_\rho-f_0\in\Vcal$ and $g=P_{\Vcal}c$ gives
\[
2d_\rho+\rho C_\rho g=0.
\]
Taking the inner product with $c$ yields
\[
C_\rho=C_0-\frac{\rho}{2}AC_\rho,
\]
which proves the result.
\end{proof}

{\begin{corollary}[Centered-spread form of the fixed-space Direct path]
\label{cor:path_scaling}
If $\1\in\Vcal$, then $g=f_0-\bar y\1$ and
\[
f_\rho
=
\bar y\1
+
b_\rho\bigl(f_0-\bar y\1\bigr),
\qquad
b_\rho
=
1-\frac{\rho}{2}C_\rho.
\]
Moreover, $C_0=-\|f_0-y\|_n^2$. Hence, whenever the unpenalized fit is non-perfect, $C_0<0$, $b_\rho>1$ for $\rho>0$, and $C_\rho$ approaches zero from below without changing sign at any finite $\rho$.
\end{corollary}

\begin{proof}
Because $\1\in\Vcal$,
\[
g=P_{\Vcal}(y-\bar y\1)=f_0-\bar y\1.
\]
Substituting this identity into Proposition~\ref{prop:path} gives the centered-spread representation. Also,
\[
c=(f_0-\bar y\1)+(y-f_0),
\]
and the projection residual $y-f_0$ is orthogonal to $\Vcal$, so
\[
C_0
=
\langle f_0-y,c\rangle_n
=
-\|f_0-y\|_n^2.
\]
The remaining statements follow from $C_\rho=C_0/(1+\rho A/2)$.
\end{proof}

\begin{remark}[Mechanical component of the in-sample residual--outcome covariance]
\label{rem:mechanical_covariance}
Corollary~\ref{cor:path_scaling} implies that, in this fixed-space squared-error benchmark with an intercept, $C_0$ is nonpositive even without imposing any equity interpretation: $C_0=-\|f_0-y\|_n^2$. The negative training covariance therefore contains a mechanical residual--outcome component induced by fitting the same observed $y$ that is later used to define the log valuation ratio. This identity does not hold mechanically on a new evaluation sample, and it does not imply that observed out-of-sample regressivity is spurious. Rather, it limits the interpretation of the training penalty: $C(f)$ is a controllable first-order mechanism, not a standalone estimator of inequity relative to latent market value. This is why the empirical analysis evaluates the fitted paths chronologically with proxy-aware assessor measures, ratio profiles, and residual-structure diagnostics.
\end{remark}


This equivalence is specific to the exact fixed-prediction-space benchmark. It reveals that, within that benchmark, Direct covariance regularization is no richer than a one-parameter centered spread correction. That makes it a mechanism result with an empirical prediction attached, and the prediction is tested rather than deferred: whether an explicit post-hoc correction reproduces the accuracy--equity path of a fully retrained nonlinear learner is answered for this application in Sections~\ref{subsec:posthoc_comparator} and~\ref{subsec:matched_beta_results}, where it reproduces much of the Direct first-order frontier over the range Direct attains and does not reproduce the residual-structure behaviour of the Surrogate. It does not imply that the empirical Direct-LightGBM path is one-dimensional: retraining can change tree splits and leaf structure, and the standard-library Direct implementation uses only the diagonal of the covariance Hessian.

}

The rescaling coefficient that this benchmark implies should be stated in its
general form, because the special case is easy to mistake for the rule. On a fitting block
$T$, the coefficient that drives the first-order association of a centered rescaling of a
fixed predictor $f_0$ to zero is
\[
b^{\star}_{T}
=
\frac{\widehat{\Var}_{T}(y)}{\widehat{\Cov}_{T}(f_0,y)},
\]
which is the form the frozen post-hoc comparator convention records. The familiar
$1/R^{2}$ expression is the special case in which $f_0$ is the orthogonal projection of
$y$ onto the prediction space, so that $\widehat{\Cov}_{T}(f_0,y)=\widehat{\Var}_{T}(f_0)$
and the ratio collapses. It is a theoretical diagnostic confined to that
linear/projection benchmark, and it is not the general root for a boosted predictor.
Moreover $b^{\star}_{T}$ is itself an in-sample theory diagnostic rather than an empirical
endpoint: the boosted learner used here is substantially less regressive on its own
fitting block than out of sample, so the in-sample root understates the rescaling that an
out-of-sample first-order correction requires. The frozen convention records it as a
diagnostic and does not use it to set any empirical endpoint.

\begin{proposition}[Accuracy cost]
\label{prop:path_cost}
Under Proposition~\ref{prop:path},
\[
\|f_\rho-y\|_n^2-\|f_0-y\|_n^2
=
\frac{C_0^2}{A}\bigl(1-q(\rho)\bigr)^2.
\]
\end{proposition}

Within a fixed prediction space, the targeted covariance falls to first order near $\rho=0$, while the in-sample squared-error cost is second order. For boosted trees this is a local benchmark only. Retraining can change splits, leaves, and stopping iteration.

\subsection{Surrogate penalty: weighted-projection path}

{Let
\[
D=\diag(c_1^2,\ldots,c_n^2),
\qquad
W_\rho=I+\rho D,
\qquad
e_0=f_0-y,
\]
where $f_0=P_{\Vcal}y$ is the ordinary least-squares projection defined above. The exact fixed-space Surrogate problem corresponding to Eq.~\eqref{eq:surrogate_weighted_loss} is
\begin{equation}
\label{eq:surrogate_fixed_space_problem}
f_\rho^{\mathrm{surr}}
=
\argminop_{f\in\Vcal}
\left\{
\|f-y\|_n^2
+
\rho\langle f-y,D(f-y)\rangle_n
\right\}.
\end{equation}
Because $W_\rho$ is positive definite for every $\rho\geq0$, this problem has a unique solution. Define the weighted inner product
\[
\langle u,v\rangle_{n,W_\rho}
=
\langle u,W_\rho v\rangle_n,
\]
let $P_{\Vcal}^{W_\rho}$ denote projection onto $\Vcal$ under this inner product, and define the self-adjoint positive-semidefinite operator
\[
T=P_{\Vcal}D\big|_{\Vcal}
\]
together with the baseline tail-weighted residual direction
\[
g_{\mathrm{surr}}
=
P_{\Vcal}De_0.
\]}

{\begin{proposition}[Surrogate weighted-projection path]
\label{prop:surrogate_fixed_space}
For every $\rho\geq0$,
\begin{equation}
\label{eq:surrogate_weighted_projection}
f_\rho^{\mathrm{surr}}
=
P_{\Vcal}^{W_\rho}y,
\end{equation}
and its departure from the unpenalized fit is
\begin{equation}
\label{eq:surrogate_correction_operator}
f_\rho^{\mathrm{surr}}-f_0
=
-\rho
\left(I_{\Vcal}+\rho T\right)^{-1}
g_{\mathrm{surr}}.
\end{equation}
\end{proposition}

\begin{proof}
For every $h\in\Vcal$, first-order optimality for~\eqref{eq:surrogate_fixed_space_problem} gives
\[
\langle f_\rho^{\mathrm{surr}}-y,h\rangle_n
+
\rho\langle D(f_\rho^{\mathrm{surr}}-y),h\rangle_n
=
0,
\]
or equivalently
\[
\langle f_\rho^{\mathrm{surr}}-y,h\rangle_{n,W_\rho}=0.
\]
This is exactly the weighted-projection characterization in~\eqref{eq:surrogate_weighted_projection}. Now write
\[
\delta_\rho=f_\rho^{\mathrm{surr}}-f_0\in\Vcal.
\]
Since $f_0=P_{\Vcal}y$, the ordinary projection residual $e_0=f_0-y$ is orthogonal to $\Vcal$. Substituting
$f_\rho^{\mathrm{surr}}-y=e_0+\delta_\rho$ into the first-order condition and projecting the $D$ terms onto $\Vcal$ yields
\[
\delta_\rho
+
\rho T\delta_\rho
+
\rho g_{\mathrm{surr}}
=
0.
\]
Because $T$ is positive semidefinite, $I_{\Vcal}+\rho T$ is invertible for $\rho\geq0$, which proves~\eqref{eq:surrogate_correction_operator}.
\end{proof}
}

{\begin{corollary}[Local Surrogate direction and baseline-loss cost]
\label{cor:surrogate_local}
As $\rho\downarrow0$,
\begin{equation}
\label{eq:surrogate_local_direction}
f_\rho^{\mathrm{surr}}
=
f_0
-
\rho g_{\mathrm{surr}}
+
O(\rho^2)
=
f_0
-
\rho P_{\Vcal}
\left[
c^{\odot2}\odot e_0
\right]
+
O(\rho^2).
\end{equation}
where $\odot$ denotes elementwise multiplication and $c^{\odot2}=(c_1^2,\ldots,c_n^2)^\top$. Moreover,
\begin{equation}
\label{eq:surrogate_baseline_loss_cost}
\|f_\rho^{\mathrm{surr}}-y\|_n^2
-
\|f_0-y\|_n^2
=
\rho^2\|g_{\mathrm{surr}}\|_n^2
+
O(\rho^3).
\end{equation}
\end{corollary}

\begin{proof}
The expansion
\[
(I_{\Vcal}+\rho T)^{-1}
=
I_{\Vcal}-\rho T+O(\rho^2)
\]
in~\eqref{eq:surrogate_correction_operator} gives~\eqref{eq:surrogate_local_direction}. For the second claim, $\delta_\rho=f_\rho^{\mathrm{surr}}-f_0$ belongs to $\Vcal$ and $e_0=f_0-y$ is orthogonal to $\Vcal$, so
\[
\|e_0+\delta_\rho\|_n^2-\|e_0\|_n^2
=
\|\delta_\rho\|_n^2.
\]
Using $\delta_\rho=-\rho g_{\mathrm{surr}}+O(\rho^2)$ gives~\eqref{eq:surrogate_baseline_loss_cost}.
\end{proof}
}

\subsection{Direct--Surrogate comparison in a fixed prediction space}

{The preceding results expose a structural difference that is not visible from the upper-bound relation alone. For the Direct penalty, Proposition~\ref{prop:path} shows that the entire fixed-space path moves along the single direction $P_{\Vcal}c$; when $\1\in\Vcal$, Corollary~\ref{cor:path_scaling} identifies this direction with the centered fitted-value spread $f_0-\bar y\1$. The Direct correction is therefore rank one in this benchmark. For the Surrogate, the initial direction is instead
\[
-g_{\mathrm{surr}}
=
-P_{\Vcal}
\left[c^{\odot2}\odot e_0\right],
\]
the learnable component of the baseline residual pattern after weighting observations by squared log-price distance from the sample center. The complete Surrogate path generally changes direction as $\rho$ varies.}

{To see this explicitly, let $\{u_j\}$ be an orthonormal eigenbasis of $T$ on $\Vcal$, with
\[
Tu_j=\lambda_j u_j,
\qquad
\lambda_j\geq0,
\]
and write
\[
g_{\mathrm{surr}}=\sum_j a_j u_j.
\]
Then Proposition~\ref{prop:surrogate_fixed_space} gives
\begin{equation}
\label{eq:surrogate_spectral_path}
f_\rho^{\mathrm{surr}}-f_0
=
-\sum_j
\frac{\rho}{1+\rho\lambda_j}
a_j u_j.
\end{equation}
Thus, unless the nonzero components of $g_{\mathrm{surr}}$ lie in a single eigenspace of $T$, the relative contribution of the learnable correction modes changes with $\rho$. This contrasts with the one-direction Direct path. It also makes clear that the Surrogate does not inherit the Direct formula
$C_\rho=C_0/(1+\rho A/2)$: the Surrogate minimizes a log-price-distance-weighted prediction loss rather than the signed covariance itself, so monotone covariance shrinkage is not guaranteed by the fixed-space theory. Likewise, the spectral representation does not imply that stronger Surrogate regularization must generate any particular nonlinear ratio shape; the empirical shape and $\Delta_{\mathrm{NL}}$ results remain separate diagnostics.}


{
Taken together, the fixed-space results clarify the distinct mechanisms of the two methods. Direct regularization applies a rank-one correction aligned with the projected centered-price direction and, with an intercept in the prediction space, becomes an exact centered spread rescaling. Surrogate regularization changes the projection geometry through the diagonal weights $1+\rho c_i^2$: locally it repairs the learnable tail-weighted residual pattern, and globally it can combine multiple correction modes. In both cases the baseline squared-error cost is second order near $\rho=0$, but only the Direct objective guarantees monotone shrinkage of the targeted covariance in this benchmark. The comparison therefore supports the paper's interpretation of the Surrogate as a related but distinct objective rather than as a purely computational approximation to Direct.
}

\section{Linear and Boosting Implementations}
\label{app:implementation}


\subsection{Linear Model}

Let $X\in\R^{n\times p}$ and $e=X\theta-y$. With
\[
c=y-\bar y\1,\qquad a=X^\top c,\qquad b=y^\top c,
\]
the direct penalized estimator under loss $\|X\theta-y\|_2^2/n$ satisfies
\[
\left(
\frac{1}{n}X^\top X+\frac{\rho}{2n^2}aa^\top
\right)\theta
=
\frac{1}{n}X^\top y+\frac{\rho}{2n^2}ba.
\]
This is a rank-one modification of the normal equations.

For the surrogate, let $D=\diag(c_1^2,\ldots,c_n^2)$. Then
\[
\left(
\frac{1}{n}X^\top X+\frac{\rho}{n}X^\top DX
\right)\theta
=
\frac{1}{n}X^\top y+\frac{\rho}{n}X^\top Dy.
\]

\subsection{LightGBM Custom-Objective Initialization and Scaling}

For the canonical custom objectives, let $b=\bar y$ denote the mean log sale price
in the current training sample. We fit the custom-objective model to centered labels
$\widetilde y_i=y_i-b$ with zero initial score and add $b$ back to predictions.
Equivalently, the fitted model predicts on the original log-price scale as
$f(x)=b+\widetilde f(x)$.
This construction is designed to match the native squared-error LightGBM starting prediction while leaving $e_i=f(x_i)-y_i$ and $c_i=y_i-\bar y$ unchanged.

The corresponding native/custom zero-penalty audit is reported in the next subsection.

The mathematical objectives in the main text are written as sample averages. In
the LightGBM custom-objective implementation we multiply the complete objective by
$n/2$, a common positive factor that does not change its minimizer. This yields a
native squared-error baseline gradient $g_i=e_i$ and Hessian $h_i=1$ at $\rho=0$.
For the direct objective, the supplied gradient and diagonal Hessian are
\[
g_i^{\mathrm{cov}}
=
e_i+\frac{\rho}{2}z c_i,
\qquad
h_i^{\mathrm{cov,diag}}
=
1+\frac{\rho}{2n}c_i^2,
\qquad
z=\frac{1}{n}\sum_{i=1}^n e_i c_i.
\]
For the surrogate, the supplied derivatives are
\[
g_i^{\mathrm{surr}}=e_i(1+\rho c_i^2),
\qquad
h_i^{\mathrm{surr}}=1+\rho c_i^2.
\]
Thus the direct implementation uses the exact covariance gradient and only the
diagonal of its dense covariance curvature, whereas the surrogate derivatives are
exactly observation-additive. The $\rho=0$ custom-objective path is retained as an
implementation control.

\subsection{Zero-Penalty Implementation Audit}
\label{subsec:zero_control_results}

Before attributing within-family changes to positive regularization, we compare ordinary LightGBM with the Direct and Surrogate implementations at $\rho=0$. The two custom-objective controls coincide with each other in the reported metrics, but the currently frozen artifacts are not numerically identical to ordinary LightGBM. Because the supplied zero-penalty gradient and Hessian match squared-error derivatives algebraically, this audit identifies a native-versus-custom execution-path difference rather than a positive-penalty objective effect. The paragraphs after the table report what the completed audit then established, and the boundary of that claim: the evidence isolates the setting responsible and the boosting iteration at which the two paths separate, but it does not directly instrument LightGBM's internal random-number stream, so no claim is made about library internals beyond what the tree, split and source evidence supports.

{\begin{table}[!ht]
\centering
\caption{Implementation-control comparison at $\rho=0$. Prediction differences are log-price deviations from ordinary LightGBM, aligned on sale identifiers and computed independently on the held-out and 2025 samples.}
\label{tab:rho_zero_control}
\begin{tabular}{llrrrrr}
\toprule
Sample & Model & $R^2_P$ & $\operatorname{RMSE}_{\log P}$ & $\beta_{\log}$ & Mean $|\Delta\widehat y|$ & Max $|\Delta\widehat y|$ \\
\midrule
Held-out & Ordinary LightGBM & 0.894 & 0.289 & -0.150 & -- & -- \\
Held-out & Direct, $\rho=0$ & 0.893 & 0.289 & -0.147 & 3.24e-02 & 3.09e-01 \\
Held-out & Surrogate, $\rho=0$ & 0.893 & 0.289 & -0.147 & 3.24e-02 & 3.09e-01 \\
\addlinespace
2025 forward & Ordinary LightGBM & 0.904 & 0.278 & -0.164 & -- & -- \\
2025 forward & Direct, $\rho=0$ & 0.904 & 0.279 & -0.161 & 3.06e-02 & 2.83e-01 \\
2025 forward & Surrogate, $\rho=0$ & 0.904 & 0.279 & -0.161 & 3.06e-02 & 2.83e-01 \\
\bottomrule
\end{tabular}
\end{table}
}

\paragraph{The native-versus-custom decomposition, as measured.} The parity experiment
that earlier drafts promised has been carried out, and it resolves in the negative. It compares the
three cells at three learner capacities --- a small and an intermediate diagnostic capacity, and the
canonical capacity that produced the reported paths --- on both out-of-time samples. Ordinary
LightGBM and the centered-label native L2 control agree to numerical tolerance at the two reduced
capacities: at the small capacity their held-out mean absolute log difference is 7.69e-09 with a
maximum of 2.58e-08, and the intermediate capacity is of the same order. The centered-label native
L2 control and the custom-objective $\rho=0$ origin do \emph{not} agree at any tested capacity: the
same held-out comparison gives a mean absolute log difference of 1.76e-02 with a maximum of
1.40e-01 at the small capacity, and 3.26e-02 with a maximum of 2.97e-01 at the canonical capacity.
So the label-centering step is not what separates the two implementations. At the canonical capacity
a second and separate effect --- the single-precision representation of the centered labels,
amplified by capacity --- also enters the ordinary-versus-centered comparison, which is why the
clean label-representation decomposition is legible only at reduced capacity.

Two further results fix the interpretation. Repeating the whole ladder with LightGBM's
determinism controls engaged reproduces the unpinned ladder cell for cell, with the same
centered-label-to-custom canonical-capacity held-out maximum of 2.97e-01 under both, so pinning
does not remove the difference and
run-to-run numerical variation cannot account for it; the same-host replicate check puts that
numerical floor at exactly zero. Tracing single-tree fits then isolates one setting. With the frozen
\texttt{colsample\_bytree} of 0.541 the two learners already differ in the structure of the first
tree, and pinning the histogram path changes nothing; setting \texttt{colsample\_bytree} to $1$ makes
them bitwise identical. The accepted characterization is therefore a deterministic
built-in-versus-custom feature-subsampling execution-path divergence under the canonical
feature-subsampling setting. It is not evidence of an error in the covariance objective: given the
same feature subsets, the custom $\rho=0$ objective reproduces native squared-error learning exactly.
Because \texttt{colsample\_bytree} was tuned as part of the canonical configuration, setting it to $1$
specifies a different learner rather than reproducing the frozen experiment, so it is not an
available repair. Table~\ref{tab:rho_zero_control} is accordingly reported as measured and is not
regenerated, and the workflow-benchmark comparisons in the Results are labelled as such rather than
read as penalty effects.

\paragraph{Normalization and scaling, as audited.} The executed code has been audited
against the normalization in Eqs.~\eqref{eq:baseline_learning}, \eqref{eq:direct_penalty}
and~\eqref{eq:surrogate_learning} and against the $n/2$ LightGBM scaling stated above. The audit
covers all nine fitting blocks --- the seven chronological training blocks, the development
pool, and the production block used for the forward year --- and records the same three facts in
each. The $n/2$ factor is applied analytically, so the supplied derivatives reduce to $g_i=e_i$ and
$h_i=1$ at $\rho=0$. The supplied Direct diagonal Hessian is $1+(\rho/2n)c_i^2$, against an exact
scaled Hessian $I+(\rho/2n)cc^{\top}$ whose off-diagonal rank-one term the observation-wise
interface cannot carry. The supplied Surrogate Hessian is $1+\rho c_i^2$. The executed LightGBM
parameter vector is pinned by the configuration hash
\texttt{\small 8f0f2acd83118de782604b5ca7143acfbd2af3fd186ea9376588f9bcf560585b}, re-verified against
the frozen record. The source state that produced those artifacts is bounded rather than
reconstructed. All three generating commits are ancestors of the frozen provenance head, and the
objective, model-factory, split and metric files are byte-identical at every one of them; the two
executed-path files that did drift are archived as per-file patches and adjudicated
behaviour-neutral for prediction reproduction. All three generating trees were nonetheless dirty,
and only a diff hash survives for two of them and none for the third, so the uncommitted state
cannot be reconstructed from what was recorded. The frozen provenance discloses that residual as an
irreducible limitation rather than resolving it. No factor-of-two or sample-size convention is absorbed into
a relabelled $\rho$: the $\rho$ values indexing the reported paths are the values the executed
objective received.

\subsection{Direct Penalty in Boosting}

For current predictions $f$, define $e=f-y$ and
\[
z=\frac{1}{n}e^\top c.
\]
Under squared loss, the direct objective has
\[
g_i=\frac{2e_i}{n}+\frac{\rho}{n}zc_i,
\]
and exact Hessian
\[
\nabla_f^2
=
\frac{2}{n}I+\frac{\rho}{n^2}cc^\top.
\]
The diagonal approximation uses
\[
h_i^{\mathrm{diag}}
=
\frac{2}{n}+\frac{\rho}{n^2}c_i^2.
\]
For a fixed tree with leaves $I_1,\ldots,I_J$, let
\[
b_j=\sum_{i\in I_j}e_i,
\qquad
s_j=\sum_{i\in I_j}c_i,
\qquad
n_j=|I_j|.
\]
With leaf-weight vector $w$, ridge parameter $\lambda$, and $s=(s_1,\ldots,s_J)^\top$, the exact local quadratic system is
\[
A_Tw^\star=-g_T,
\qquad
g_{T,j}=\frac{2b_j}{n}+\frac{\rho z s_j}{n},
\]
\[
A_T
=
\diag\!\left(\frac{2n_1}{n}+\lambda,\ldots,\frac{2n_J}{n}+\lambda\right)
+\frac{\rho}{n^2}ss^\top.
\]
Thus $w^\star=-A_T^{-1}g_T$ when $A_T$ is nonsingular. This solver retains the cross-observation curvature but requires a custom tree-construction procedure; it is not available through the usual observation-wise API.

\subsection{Surrogate in Boosting}

For
\[
\Psi^{\mathrm{surr}}(f)
=
\frac{1}{n}\sum_i e_i^2c_i^2,
\]
the derivatives are observation-wise:
\[
g_i=\frac{2e_i}{n}+\frac{2\rho}{n}e_ic_i^2,
\qquad
h_i=\frac{2}{n}+\frac{2\rho}{n}c_i^2.
\]
This is the canonical surrogate implementation evaluated in the CCAO regularization path.


\section{\texorpdfstring{
Primary CCAO Regularization-Path Details}{Primary CCAO Regularization-Path Details}}
\label{app:ccao_path_details}


This appendix contains the complete diagnostics underlying the compact main-text regularization-path presentation. It is organized around the final augmented paths and not around a selected configuration; the compact nominal decade-spaced anchors are presentation-only snapshots of those frozen paths.

\subsection{Complete Regularization-Path Results}

The complete regularization path is held in the frozen combined path table, which carries every
family--$\rho$ configuration together with the seven fold values, equal-weight CV mean and
standard deviation, held-out results, and 2025 forward results for the
Section~\ref{subsec:metrics} evaluation metrics. That table is an output of the executed run root
rather than a released file: the frozen evidence manifest identifies it by path, byte length and
content hash, and it is not redistributed with this paper. Appendix~\ref{app:ccao_results} sets out
which derived artifacts are tracked and which are hash-identified only.
The compact main-text presentation is split in two. Table~\ref{tab:path_anchor_frozen} carries prediction, PRD and MKI at the two positive anchors the frozen reproduction audit re-derives, and Table~\ref{tab:path_anchor_standards} carries PRB, from the adopted Standard on Ratio Studies, and VEI, an Exposure-Draft diagnostic and therefore proposed rather than adopted, at all five nominal decade-spaced anchors $\{0.01,0.1,1,10,100\}$ mapped to the nearest frozen-grid points. Both show the ordinary-LightGBM workflow benchmark and the custom-objective zero-penalty origin as explicit rows; the wider zero-penalty implementation audit stays in the implementation appendix. No path row is omitted from the machine-readable results because of its observed performance.
For completeness, the complementary baseline diagnostics and the complementary representative-$\rho$ diagnostics are reported here rather than in the main empirical narrative.

\begin{table}[!htbp]
\centering
\scriptsize
\setlength{\tabcolsep}{2.4pt}
\renewcommand{\arraystretch}{1.06}
\newcommand{\baselinecompmetric}[1]{\hspace*{0.6em}#1}
\caption{Complementary baseline level for \emph{Ordinary LightGBM (standard raw-label native)}, the assessor-facing workflow benchmark: valuation level, horizontal uniformity, and mechanism/residual structure. A level statement about one model class, not a comparison between model classes.}
\label{tab:ccao_baseline_complementary}
\begin{tabularx}{\textwidth}{@{} >{\raggedright\arraybackslash}p{6.4cm} >{\centering\arraybackslash}X >{\centering\arraybackslash}X @{}}
\toprule
\textbf{Measure} & \textbf{Held-out evaluation} & \textbf{2025 forward evaluation} \\
\midrule
\multicolumn{3}{@{}l}{\textbf{Valuation level}} \\
\cmidrule(r){1-1}
\baselinecompmetric{Median ratio $m_S$} & 0.929 & 0.950 \\
\baselinecompmetric{Mean ratio $\bar r_S$} & 0.989 & 1.015 \\
\baselinecompmetric{Weighted mean $\bar r_{W,S}$} & 0.924 & 0.941 \\
\addlinespace[5pt]
\multicolumn{3}{@{}l}{\textbf{Horizontal uniformity}} \\
\cmidrule(r){1-1}
\baselinecompmetric{COD} & 21.6\% & 21.3\% \\
\baselinecompmetric{COV} & 39.7\% & 37.2\% \\
\addlinespace[5pt]
\multicolumn{3}{@{}l}{\textbf{Mechanism and residual structure}} \\
\cmidrule(r){1-1}
\baselinecompmetric{$\beta_{\log}$} & -0.150 & -0.164 \\
\baselinecompmetric{$\Delta_{\mathrm{NL}}$} & 0.119 & 0.121 \\
\baselinecompmetric{$\operatorname{dCor}(e,y)$} & 0.387 & 0.422 \\
\bottomrule
\end{tabularx}
\vspace{1mm}
\begin{minipage}{\textwidth}
\scriptsize
\emph{Notes.}
Every entry is cell A of the frozen full zero-penalty control, which carries all eight complementary measures on both evaluations. The linear-regression columns this table previously carried have been removed rather than requalified, because no frozen artifact reproduces any of those values. There is no second column to compare against, so no entry is boldfaced.
Median, mean, weighted-mean ratio, COD, and COV are assessor-facing valuation-level or horizontal-uniformity diagnostics defined in Section~\ref{subsec:metrics}. The first-order mechanism slope $\beta_{\log}$, correlation-ratio nonlinearity gap $\Delta_{\mathrm{NL}}$, and distance correlation are defined in Eqs.~\eqref{eq:log_ratio_slope}, \eqref{eq:nonlinearity_gap}, and~\eqref{eq:dcor_diagnostic}; the last two carry no assessor-standard reference range and are residual-structure diagnostics rather than equity measures.
One reading matters for the motivation and is stated rather than left out: the workflow benchmark's COD is 21.6\% held-out and 21.3\% in 2025, both above the $[5,15]$ range the adopted Standard on Ratio Studies gives for single-family residential property (Table~\ref{tab:assessment_metrics_summary}), and COD rises further along the penalty paths (Appendix Table~\ref{tab:path_anchor_complementary}). No compliance claim is implied anywhere in this paper. This range comparison is interpretation offered here on values that do resolve from the frozen control; it is not itself a finding of the frozen inferential stage.
COD and COV are reported in percent; COV is stored as a fraction in the frozen artifact and converted for display.
\end{minipage}
\end{table}


\begin{table}[!htbp]
\centering
\scriptsize
\setlength{\tabcolsep}{3.4pt}
\renewcommand{\arraystretch}{1.08}
\caption{Complementary diagnostics at the two frozen positive-$\rho$ anchors: valuation
level, horizontal uniformity, and the first-order mechanism slope, with both reference cells shown as
explicit rows. Boldface marks movement toward the measure's own stated direction relative to
\emph{Ordinary LightGBM (standard raw-label native)}, the assessor-facing workflow benchmark; a dagger
(\dag) marks movement in that direction relative to the \emph{Custom-objective $\rho=0$ origin}, the
primary within-path penalty-isolating reference. The two reference rows carry no markers. Neither
device denotes statistical significance, standards compliance, model selection, or a preferred penalty
strength: no penalty strength is selected anywhere in this paper.}
\label{tab:path_anchor_complementary}
\resizebox{\textwidth}{!}{%
\begin{tabular}{@{}llrrrrrr@{}}
\toprule
Method & $\rho$
& Median $r$
& Mean $r$
& W. mean $r$
& COD
& COV
& $\beta_{\log}$ \\
\midrule
\multicolumn{8}{l}{\textit{Panel A: Held-out evaluation}} \\
\addlinespace[2pt]
\multicolumn{8}{l}{\textbf{Reference cells}} \\
\cmidrule(lr){1-8}
Ordinary LightGBM (standard raw-label native) & -- & 0.929 & 0.989 & 0.924 & 21.6\% & 39.7\% & -0.150 \\
Custom-objective $\rho=0$ origin & $0$ & 0.928 & 0.986 & 0.923 & 21.6\% & 39.6\% & -0.147 \\
\addlinespace[3pt]
\multicolumn{8}{l}{\textbf{Direct penalty}} \\
\cmidrule(lr){1-8}
Direct & $\approx 0.954095$ & \textbf{0.936}\textsuperscript{\dag} & \textbf{0.993}\textsuperscript{\dag} & \textbf{0.936}\textsuperscript{\dag} & 21.6\% & 39.9\% & \textbf{-0.134}\textsuperscript{\dag} \\
Direct & $100$ & \textbf{0.968}\textsuperscript{\dag} & 1.022 & \textbf{0.993}\textsuperscript{\dag} & 24.5\% & 41.5\% & \textbf{-0.079}\textsuperscript{\dag} \\
\addlinespace[3pt]
\multicolumn{8}{l}{\textbf{Surrogate penalty}} \\
\cmidrule(lr){1-8}
Surrogate & $\approx 0.954095$ & \textbf{0.929}\textsuperscript{\dag} & 0.975 & \textbf{0.930}\textsuperscript{\dag} & \textbf{21.4\%}\textsuperscript{\dag} & \textbf{38.2\%}\textsuperscript{\dag} & \textbf{-0.102}\textsuperscript{\dag} \\
Surrogate & $100$ & 0.923 & 0.937 & \textbf{0.928}\textsuperscript{\dag} & 23.5\% & 40.3\% & \textbf{-0.001}\textsuperscript{\dag} \\
\midrule
\addlinespace[2pt]
\multicolumn{8}{l}{\textit{Panel B: 2025 forward evaluation}} \\
\addlinespace[2pt]
\multicolumn{8}{l}{\textbf{Reference cells}} \\
\cmidrule(lr){1-8}
Ordinary LightGBM (standard raw-label native) & -- & 0.950 & 1.015 & 0.941 & 21.3\% & 37.2\% & -0.164 \\
Custom-objective $\rho=0$ origin & $0$ & 0.946 & 1.009 & 0.936 & 21.2\% & 37.1\% & -0.161 \\
\addlinespace[3pt]
\multicolumn{8}{l}{\textbf{Direct penalty}} \\
\cmidrule(lr){1-8}
Direct & $\approx 0.954095$ & \textbf{0.957}\textsuperscript{\dag} & 1.019 & \textbf{0.954}\textsuperscript{\dag} & \textbf{21.2\%}\textsuperscript{\dag} & 37.2\% & \textbf{-0.148}\textsuperscript{\dag} \\
Direct & $100$ & \textbf{0.992}\textsuperscript{\dag} & 1.049 & \textbf{1.012}\textsuperscript{\dag} & 23.3\% & 40.3\% & \textbf{-0.093}\textsuperscript{\dag} \\
\addlinespace[3pt]
\multicolumn{8}{l}{\textbf{Surrogate penalty}} \\
\cmidrule(lr){1-8}
Surrogate & $\approx 0.954095$ & 0.950\textsuperscript{\dag} & \textbf{1.003}\textsuperscript{\dag} & \textbf{0.946}\textsuperscript{\dag} & \textbf{20.9\%}\textsuperscript{\dag} & \textbf{36.1\%}\textsuperscript{\dag} & \textbf{-0.120}\textsuperscript{\dag} \\
Surrogate & $100$ & 0.948\textsuperscript{\dag} & 0.964 & \textbf{0.945}\textsuperscript{\dag} & 22.7\% & \textbf{36.8\%}\textsuperscript{\dag} & \textbf{-0.015}\textsuperscript{\dag} \\
\bottomrule
\end{tabular}}
\vspace{1mm}
\begin{minipage}{\textwidth}
\scriptsize
\emph{Notes.}
Median, mean and weighted-mean ratio, COD and COV are the assessor-facing valuation-level
and horizontal-uniformity diagnostics defined in Section~\ref{subsec:metrics}; the first-order
mechanism slope $\beta_{\log}$ is defined in Eq.~\eqref{eq:log_ratio_slope}. Reference-cell values are
taken from the frozen full zero-penalty control and the two positive-$\rho$ rows from the frozen
reproduction audit, which is why only the two frozen anchors appear: the other three nominal display
anchors are not reproduced by any frozen artifact and are omitted rather than shown blank.
The correlation-ratio nonlinearity gap $\Delta_{\mathrm{NL}}$ and the distance correlation
$\operatorname{dCor}(e,y)$ are deliberately \emph{not} columns of this table. The frozen reproduction
artifact carries neither measure at any positive $\rho$, so no positive-$\rho$ value for either could
be re-derived; the ordinary-LightGBM values are reported in Appendix
Table~\ref{tab:ccao_baseline_complementary} instead.
The three valuation-level ratios are compared by distance from one. COD and COV have no universal
ideal --- the frozen guidance notes that unusually low dispersion can itself warrant diagnostic review
--- so a marker on either means only that dispersion moved lower, not that the configuration is
better; the adopted 2013 COD range for single-family residential is $[5,15]$, and every COD entry in
this table lies above it. A marker on $\beta_{\log}$ means the first-order slope moved toward zero,
which is a mechanism direction and not an assessor-standard equity target. Comparisons are evaluated
on the unrounded canonical values, so two cells that appear equal after display rounding can still be
marked differently. COD and COV are reported in percent.
\end{minipage}
\end{table}


The frozen combined path table carries every family--$\rho$ configuration for the Section~\ref{subsec:metrics} metrics, together with the fold values, the equal-weight CV mean and standard deviation, and the held-out and 2025 forward results; as Appendix~\ref{app:ccao_results} records, it is hash-identified in the frozen evidence manifest rather than redistributed. Chronological-CV $\Delta_{\mathrm{NL}}$ was reconstructed fold by fold from the already-frozen validation predictions using the same fixed cross-fitted estimator specified in Appendix~\ref{app:metrics}; those CV values did not define the $\rho$ grid or any model-selection rule. Split-specific $\rho=0$ prediction-parity audits belong to the same result root.


\subsection{Prediction, Valuation-Level, and Horizontal-Uniformity Paths}


The $\rho$-evolution figures below trace the frozen tested grid descriptively. They mark no penalty interval, and this paper does not select, recommend, or certify a penalty strength.

\begin{figure}[!htbp]
\centering
\safeincludegraphics[width=0.66\textwidth,height=0.34\textheight,keepaspectratio]{img/generated_v12_994/predictive_metric_paths.pdf}\\[1mm]
\safeincludegraphics[width=0.66\textwidth,height=0.34\textheight,keepaspectratio]{img/generated_v12_994/level_uniformity_paths.pdf}
\caption{Predictive-metric and valuation-level/uniformity paths along the Direct and Surrogate grids, with held-out and 2025 evaluations overlaid. Valuation-level panels include the ratio $=1$ reference. No penalty strength is selected in this paper.}
\label{fig:other_metric_paths_placeholder}
\end{figure}\begin{figure}[!htbp]
\centering
\safeincludegraphics[width=0.8\textwidth]{img/generated_v12_994/vertical_equity_metric_paths.pdf}
\caption{Vertical-equity metric paths versus $\rho$ for Direct and Surrogate. Reference lines mark PRD $=1$, PRB $=0$, MKI $=1$, and VEI $=0$. No penalty strength is selected in this paper.}
\label{fig:vertical_equity_metric_paths}
\end{figure}

\subsection{Exact Rolling-Origin Fold Structure}

Table~\ref{tab:fold_structure} records the exact cumulative windows and sample sizes underlying the seven-fold chronological design described in Section~\ref{subsec:path_design}. In every fold, the chronologically newest 10\% of the cumulative window is used for validation; because that rule is identical across folds, it is stated once here rather than repeated as a table column.

{\begin{table}[!ht]
\centering
\caption{Rolling-origin validation structure for the final CCAO development sample. In every fold, the newest 10\% of observations in the cumulative window form the validation block. The fold-6 and fold-7 validation blocks overlap; accordingly, fold-level statistics computed across these rows are descriptive chronological-window summaries rather than independent replications, and they support no IID standard errors. See Section~\ref{subsec:path_design} for the development coordinates and the one-sale-one-vote D3 sensitivity.}
\label{tab:fold_structure}
\begin{tabular}{clrr}
\toprule
Fold & Cumulative window & $n_{\mathrm{train}}$ & $n_{\mathrm{val}}$ \\
\midrule
1 & 2016-01-01--2017-04-01 & 46,888 & 5,209 \\
2 & 2016-01-01--2018-07-01 & 100,776 & 11,197 \\
3 & 2016-01-01--2019-10-01 & 151,187 & 16,798 \\
4 & 2016-01-01--2021-01-01 & 200,908 & 22,323 \\
5 & 2016-01-01--2022-04-01 & 252,486 & 28,053 \\
6 & 2016-01-01--2023-07-01 & 298,022 & 33,113 \\
7 & 2016-01-01--2023-11-09 & 310,147 & 34,460 \\
\bottomrule
\end{tabular}
\end{table}
}

\subsection{Chronological Fold Stability}


That summary figure is removed from the active presentation because the comprehensive fold figures immediately below now report prediction, valuation level and horizontal uniformity, vertical equity, and mechanism paths separately for all available CV metrics.\begin{figure}[!htbp]
\centering
\safeincludegraphics[width=0.8\textwidth]{img/generated_v12_994/cv_predictive_metric_paths.pdf}
\caption{Chronological-fold predictive-metric paths (thin gray) and equal-weight CV means (thick). The seven folds are overlapping cumulative-window constructions rather than independent replications, so their spread is descriptive and supports no standard error. No penalty strength is selected in this paper.}
\label{fig:cv_predictive_metric_paths}
\end{figure}\begin{figure}[!htbp]
\centering
\safeincludegraphics[width=0.8\textwidth]{img/generated_v12_994/cv_level_uniformity_paths.pdf}
\caption{Chronological-fold valuation-level and uniformity paths (thin gray) and equal-weight CV means (thick). Valuation-level panels include the ratio $=1$ reference. The seven folds are overlapping cumulative-window constructions rather than independent replications, so their spread is descriptive and supports no standard error. No penalty strength is selected in this paper.}
\label{fig:cv_level_uniformity_paths}
\end{figure}\begin{figure}[!htbp]
\centering
\safeincludegraphics[width=0.8\textwidth]{img/generated_v12_994/cv_vertical_equity_metric_paths.pdf}
\caption{Chronological-fold vertical-equity paths (thin gray) and equal-weight CV means (thick). Reference lines mark PRD $=1$, PRB $=0$, MKI $=1$, and VEI $=0$. The seven folds are overlapping cumulative-window constructions rather than independent replications, so their spread is descriptive and supports no standard error. No penalty strength is selected in this paper.}
\label{fig:cv_vertical_equity_metric_paths}
\end{figure}\begin{figure}[!htbp]
\centering
\safeincludegraphics[width=0.8\textwidth]{img/generated_v12_994/cv_mechanism_metric_paths.pdf}
\caption{Chronological-fold mechanism paths for $\beta_{\log}$, $\Delta_{\mathrm{NL}}$, and distance correlation (thin gray) and equal-weight CV means (thick). The $\beta_{\log}=0$ line is a first-order neutrality reference; $\Delta_{\mathrm{NL}}$ and dCor are not given forced zero lines. CV $\Delta_{\mathrm{NL}}$ was reconstructed from frozen validation predictions with the same fixed estimator used out of time; it is a retrospective diagnostic and does not select $\rho$. The seven folds are overlapping cumulative-window constructions rather than independent replications, so their spread is descriptive and supports no standard error. No penalty strength is selected in this paper.}
\label{fig:cv_mechanism_metric_paths}
\end{figure}

The reconstructed CV $\Delta_{\mathrm{NL}}$ paths in the figure above are a retrospective diagnostic. No minimum value and no grid location is quoted from them, because the reconstructed CV path values are not part of the frozen numeric evidence, and no location on those paths is a selected operating point.

\begin{figure}[!htbp]
\centering
\safeincludegraphics[width=0.95\textwidth]{img/generated_v12_994/tradeoff_equity_vs_accuracy_heldout.pdf}
\caption{Held-out assessor-facing diagnostics (horizontal) versus predictive metrics (vertical) along the Direct and Surrogate paths. Dotted vertical lines mark metric neutrality where applicable, and shaded regions are assessor-facing guidance ranges rather than compliance claims.}
\label{fig:tradeoff_equity_vs_accuracy_heldout}
\end{figure}\begin{figure}[!htbp]
\centering
\safeincludegraphics[width=0.95\textwidth]{img/generated_v12_994/tradeoff_equity_vs_accuracy_2025.pdf}
\caption{2025 assessor-facing diagnostics (horizontal) versus predictive metrics (vertical) along the Direct and Surrogate paths. Dotted vertical lines mark metric neutrality where applicable, and shaded regions are assessor-facing guidance ranges rather than compliance claims.}
\label{fig:tradeoff_equity_vs_accuracy_2025}
\end{figure}\begin{figure}[!htbp]
\centering
\safeincludegraphics[width=0.95\textwidth]{img/generated_v12_994/tradeoff_mechanism_vs_accuracy_heldout.pdf}
\caption{Held-out mechanism/residual diagnostics (horizontal) versus predictive metrics (vertical). The dotted vertical line at $\beta_{\log}=0$ is a first-order neutrality reference where applicable; $\Delta_{\mathrm{NL}}$ and dCor are not given assessor-style zero targets.}
\label{fig:tradeoff_mechanism_vs_accuracy_heldout}
\end{figure}\begin{figure}[!htbp]
\centering
\safeincludegraphics[width=0.95\textwidth]{img/generated_v12_994/tradeoff_mechanism_vs_accuracy_2025.pdf}
\caption{2025 mechanism/residual diagnostics (horizontal) versus predictive metrics (vertical). The dotted vertical line at $\beta_{\log}=0$ is a first-order neutrality reference where applicable; $\Delta_{\mathrm{NL}}$ and dCor are not given assessor-style zero targets.}
\label{fig:tradeoff_mechanism_vs_accuracy_2025}
\end{figure}

\subsection{Centered-Spread Post-Hoc Comparator Detail}
\label{app:centered_spread_detail}

Table~\ref{tab:centered_spread} reports the comparator of
Section~\ref{subsec:posthoc_comparator} at the frozen base-value anchors, for the primary base cell
on both out-of-time blocks and for the secondary practical base cell on the held-out block. The
complete path over all frozen base values, the per-fold summaries with their descriptive
chronological spread, the base-value-one bitwise check, the linearity check against the closed form,
the root verification, and the anchor-level centering sensitivity are tracked in the repository as
part of the frozen post-hoc comparator record. Two properties of the construction are worth restating here because they bound what the
table can be read to show. First, $\beta_{\log}$ is exactly linear in the base value on a fixed
evaluation sample, so the post-hoc first-order coordinate is solved rather than searched, and the
frozen record verifies the achieved development value against the closed form. Second, the
comparator refits nothing: it is an affine transformation of a stored prediction array, which is why
it can match a first-order coordinate exactly and cannot alter non-affine conditional-mean
structure except incidentally.

{\begin{table}[!htbp]
\centering
\scriptsize
\setlength{\tabcolsep}{2.6pt}
\renewcommand{\arraystretch}{1.06}
\caption{Centered-spread post-hoc comparator at the frozen base-value anchors. The map
is $f_b(x)=\bar y_T+b(f_0(x)-\bar y_T)$ applied to cached prediction arrays; no model was refit.
D1 is the primary development coordinate and D2 the pooled out-of-fold sensitivity, so ``D1 root''
is the base value that drives the primary development coordinate to first-order neutrality and
``D2 root'' the same for the sensitivity coordinate. VEI is a May-2026 Exposure Draft diagnostic and
is proposed rather than adopted guidance, so no VEI entry is a compliance determination. No base
value is selected or recommended.}
\label{tab:centered_spread}
\resizebox{\textwidth}{!}{%
\begin{tabular}{@{}lrrrrrrrrrrr@{}}
\toprule
Anchor & $b$ & $R^2_P$ & MAE & $\operatorname{RMSE}_{\log P}$ & COD & PRD & MKI
 & VEI & $\beta_{\log}$ & $\Delta_{\mathrm{NL}}$ & dCor \\
\midrule
\addlinespace[2pt]
\multicolumn{12}{l}{\textit{Panel A: \emph{Custom-objective $\rho=0$ origin} rescaled, held-out (primary)}} \\
\addlinespace[2pt]
\cmidrule(lr){1-12}
$b=1$ (no rescaling) & 1.000000 & 0.8928 & \$75,976 & 0.2894 & 21.60 & 1.0690 & 0.9230 & -26.17 & -0.1474 & 0.1164 & 0.3822 \\
D1 root & 1.160458 & 0.8063 & \$87,182 & 0.3071 & 22.59 & 0.9871 & 1.0726 & 14.77 & -0.0106 & 0.1343 & 0.2518 \\
D2 root & 1.167027 & 0.7941 & \$88,689 & 0.3087 & 22.69 & 0.9839 & 1.0787 & 16.44 & -0.0050 & 0.1344 & 0.2526 \\
$b_{\max}$ & 1.209082 & 0.6926 & \$99,978 & 0.3200 & 23.50 & 0.9630 & 1.1175 & 26.82 & 0.0309 & 0.1338 & 0.2681 \\
\addlinespace[2pt]
\multicolumn{12}{l}{\textit{Panel B: \emph{Custom-objective $\rho=0$ origin} rescaled, 2025 forward (primary)}} \\
\addlinespace[2pt]
\cmidrule(lr){1-12}
$b=1$ (no rescaling) & 1.000000 & 0.9039 & \$79,139 & 0.2791 & 21.23 & 1.0775 & 0.9093 & -28.41 & -0.1612 & 0.1212 & 0.4168 \\
D1 root & 1.160458 & 0.7991 & \$96,243 & 0.2939 & 21.91 & 0.9934 & 1.0581 & 11.37 & -0.0266 & 0.1454 & 0.2542 \\
D2 root & 1.167027 & 0.7850 & \$98,248 & 0.2955 & 22.01 & 0.9900 & 1.0642 & 12.97 & -0.0211 & 0.1457 & 0.2536 \\
$b_{\max}$ & 1.209082 & 0.6682 & \$112,802 & 0.3069 & 22.75 & 0.9685 & 1.1029 & 23.74 & 0.0142 & 0.1459 & 0.2619 \\
\addlinespace[2pt]
\multicolumn{12}{l}{\textit{Panel C: \emph{Ordinary LightGBM (standard raw-label native)} rescaled, held-out (secondary, practical)}} \\
\addlinespace[2pt]
\cmidrule(lr){1-12}
$b=1$ (no rescaling) & 1.000000 & 0.8942 & \$75,655 & 0.2887 & 21.63 & 1.0695 & 0.9230 & -26.46 & -0.1496 & 0.1187 & 0.3871 \\
D1 root & 1.160367 & 0.8024 & \$87,658 & 0.3060 & 22.61 & 0.9875 & 1.0728 & 14.34 & -0.0132 & 0.1375 & 0.2531 \\
D2 root & 1.167266 & 0.7892 & \$89,277 & 0.3076 & 22.73 & 0.9840 & 1.0792 & 16.18 & -0.0074 & 0.1376 & 0.2537 \\
$b_{\max}$ & 1.209082 & 0.6852 & \$100,604 & 0.3188 & 23.53 & 0.9633 & 1.1180 & 27.33 & 0.0282 & 0.1371 & 0.2678 \\
\bottomrule
\end{tabular}}
\vspace{1mm}
\begin{minipage}{\textwidth}
\scriptsize
\emph{Notes.}
Panels A and B rescale the \emph{Custom-objective $\rho=0$ origin}, the primary post-hoc
base because it is the common origin of both penalized families, so rescaling it holds the execution
path fixed. Panel~C rescales \emph{Ordinary LightGBM (standard raw-label native)} and is the
secondary, practical variant; under the frozen convention it determines neither the comparison range
nor the choice of comparison targets nor the primary mechanistic test. \emph{Centered-label native
L2 (initialization-aligned)} has no post-hoc path. At $b=1$ the map returns the cached prediction
array unchanged, so the first row of each panel is the reference cell itself. VEI is reported in
percent and COD in percent; COD lies outside the adopted $[5,15]$ range throughout. The alternative
centering about the fitting block's mean prediction is a named sensitivity evaluated at these same
four base values; its analytic difference from the map used here is constant in $x$, so it cannot
move $\beta_{\log}$, COD, COV or dCor at all. $\Delta_{\mathrm{NL}}$ and dCor carry no
assessor-facing reference range.
\end{minipage}
\end{table}
}


\subsection{Matched First-Order Targets and Attainment}
\label{app:matched_beta_detail}

Tables~\ref{tab:matched_beta} and~\ref{tab:matched_beta_ext} carry the design and the
outcomes behind Section~\ref{subsec:matched_beta_results}: the six comparison targets inside the
three-way development range with the configuration each family used, the held-out outcome at those
targets, and the three extended targets outside the range with an explicit attainment state for
every family. The D1 coordinate is primary throughout and the one-sale-one-vote D3 coordinate is the
sensitivity; the D3 reselection counts and sign-agreement classification reported in
Section~\ref{subsec:matched_beta_results} come from the frozen D3 branch, which re-runs the
identical deterministic construction on the row-balanced coordinate and reproduces its targets from
the pre-outcome record. The equal-weight and forward-year outcome columns, the pairwise difference
tables, the crossing locations, the origin-identity check and the D3 twins of all four tables are
tracked in the repository as part of the frozen matched-target record.

{\begin{table}[!htbp]
\centering
\scriptsize
\setlength{\tabcolsep}{4pt}
\renewcommand{\arraystretch}{1.06}
\caption{Matched first-order comparison on the primary development coordinate D1.
Panel~A gives the six comparison targets and the configuration each family used; Panel~B gives the
held-out outcome at those targets for the three primary families. Every configuration is an actually
fitted model or an exactly solved rescaling, and no metric is interpolated between penalty values.
The match is on the \emph{development} coordinate and does not equalise held-out first-order
correction. No configuration is selected or recommended.}
\label{tab:matched_beta}
\resizebox{\textwidth}{!}{%
\begin{tabular}{@{}llrrrrrr@{}}
\toprule
\multicolumn{8}{l}{\textit{Panel A: comparison targets and the matched configurations}} \\
\cmidrule(lr){1-8}
$j$ & Target $\beta_{\log}$ (D1) & Direct $\rho$ & Surrogate $\rho$ & C-posthoc $b$
 & A-posthoc $b$ & & \\
\midrule
0 & -0.138271 & 0 & 0 & 1.000000 & 0.999922 \\
1 & -0.126490 & 0.7197 & 0.1326 & 1.013671 & 1.013592 \\
2 & -0.114733 & 1.6768 & 0.4095 & 1.027316 & 1.027235 \\
3 & -0.102033 & 3.9069 & 0.8286 & 1.042053 & 1.041972 \\
4 & -0.090902 & 9.1030 & 1.4563 & 1.054971 & 1.054888 \\
5 & -0.078651 & 86.8511 & 2.5595 & 1.069186 & 1.069103 \\
\midrule
\multicolumn{8}{l}{\textit{Panel B: held-out outcome at those targets}} \\
\cmidrule(lr){1-8}
$j$ & Family & $R^2_P$ & MAE & $\operatorname{RMSE}_{\log P}$ & COD
 & $\Delta_{\mathrm{NL}}$ & dCor \\
\midrule
\addlinespace[2pt]
0 & Direct & 0.8928 & \$75,976 & 0.2894 & 21.60 & 0.1164 & 0.3822 \\
 & Surrogate & 0.8928 & \$75,976 & 0.2894 & 21.60 & 0.1164 & 0.3822 \\
 & C-posthoc & 0.8928 & \$75,976 & 0.2894 & 21.60 & 0.1164 & 0.3822 \\
\addlinespace[2pt]
1 & Direct & 0.8984 & \$74,516 & 0.2888 & 21.60 & 0.1216 & 0.3637 \\
 & Surrogate & 0.8954 & \$75,416 & 0.2909 & 21.50 & 0.1132 & 0.3660 \\
 & C-posthoc & 0.8972 & \$74,897 & 0.2893 & 21.55 & 0.1192 & 0.3625 \\
\addlinespace[2pt]
2 & Direct & 0.9006 & \$74,167 & 0.2900 & 21.57 & 0.1231 & 0.3424 \\
 & Surrogate & 0.8948 & \$75,831 & 0.2948 & 21.36 & 0.1035 & 0.3358 \\
 & C-posthoc & 0.9000 & \$74,151 & 0.2895 & 21.52 & 0.1218 & 0.3438 \\
\addlinespace[2pt]
3 & Direct & 0.8993 & \$74,164 & 0.2910 & 21.70 & 0.1292 & 0.3227 \\
 & Surrogate & 0.8963 & \$75,505 & 0.2976 & 21.40 & 0.0983 & 0.3134 \\
 & C-posthoc & 0.9009 & \$73,725 & 0.2900 & 21.52 & 0.1243 & 0.3248 \\
\addlinespace[2pt]
4 & Direct & 0.8958 & \$75,296 & 0.2937 & 21.93 & 0.1298 & 0.3017 \\
 & Surrogate & 0.8955 & \$76,083 & 0.3030 & 21.39 & 0.0938 & 0.2938 \\
 & C-posthoc & 0.8999 & \$73,755 & 0.2908 & 21.55 & 0.1263 & 0.3096 \\
\addlinespace[2pt]
5 & Direct & 0.8760 & \$82,764 & 0.3170 & 23.97 & 0.1198 & 0.2730 \\
 & Surrogate & 0.8963 & \$76,128 & 0.3094 & 21.40 & 0.0914 & 0.2714 \\
 & C-posthoc & 0.8966 & \$74,236 & 0.2920 & 21.61 & 0.1282 & 0.2944 \\
\bottomrule
\end{tabular}}
\vspace{1mm}
\begin{minipage}{\textwidth}
\scriptsize
\emph{Notes.}
Targets are placed by a deterministic rule: probe equally spaced points in the
three-way development range, then take as the target the achieved development $\beta_{\log}$ of the
fitted Direct configuration nearest the probe. The target is therefore an achieved value rather than
the probe, Direct matches exactly by construction, and no selection discretion enters. Surrogate
rows use the nearest actually fitted penalty value and the post-hoc rows an exact solve, all within
the frozen tolerance $\tau=$ 0.002. The $j=0$ row is an origin identity check: Direct at
$\rho=0$, Surrogate at $\rho=0$ and the post-hoc map at $b=1$ are the same object. A-posthoc is the
rescaling of \emph{Ordinary LightGBM (standard raw-label native)} and is carried as a secondary,
practical variant; per the frozen convention it determines neither the comparison range, nor the
target selection, nor whether a target exists, and its outcome columns are therefore not reported
here. Direct and Surrogate penalty magnitudes are not comparable across families. COD is in percent
and lies outside the adopted $[5,15]$ range throughout. $\Delta_{\mathrm{NL}}$ and dCor are
residual-structure diagnostics with no assessor-facing reference range.
\end{minipage}
\end{table}
}

{\begin{table}[!htbp]
\centering
\scriptsize
\setlength{\tabcolsep}{3.4pt}
\renewcommand{\arraystretch}{1.06}
\caption{Extended development targets, which lie outside the three-way common range,
with an explicit attainment state for every family. Held-out evaluation. Rows recorded as not
attained carry blank outcome columns: they are neither dropped nor filled by interpolation. The
states describe the penalty grid this experiment evaluated, not the objectives and not penalty
values the experiment never examined.}
\label{tab:matched_beta_ext}
\begin{tabular}{@{}llllrrrrrrr@{}}
\toprule
Target $\beta_{\log}$ (D1) & Family & Configuration & State & Max attained $\beta_{\log}$
 & $R^2_P$ & MAE & $\operatorname{RMSE}_{\log P}$ & COD & $\Delta_{\mathrm{NL}}$ & dCor \\
\midrule
\addlinespace[2pt]
-0.06 & Direct & --- & \verb|NOT_ATTAINED| & -0.078651 & --- & --- & --- & --- & --- & --- \\
 & Surrogate & 5.9636 & attained & --- & 0.8934 & \$77,813 & 0.3230 & 21.81 & 0.0968 & 0.2518 \\
 & C-posthoc & 1.090830 & attained & --- & 0.8868 & \$75,847 & 0.2944 & 21.75 & 0.1306 & 0.2752 \\
\addlinespace[2pt]
-0.03 & Direct & --- & \verb|NOT_ATTAINED| & -0.078651 & --- & --- & --- & --- & --- & --- \\
 & Surrogate & 28.1177 & attained & --- & 0.8917 & \$79,939 & 0.3443 & 22.84 & 0.1126 & 0.2553 \\
 & C-posthoc & 1.125644 & attained & --- & 0.8569 & \$80,444 & 0.2999 & 22.10 & 0.1332 & 0.2559 \\
\addlinespace[2pt]
0.00 & Direct & --- & \verb|NOT_ATTAINED| & -0.078651 & --- & --- & --- & --- & --- & --- \\
 & Surrogate & --- & \verb|NOT_ATTAINED| & -0.018609 & --- & --- & --- & --- & --- & --- \\
 & C-posthoc & 1.160458 & attained & --- & 0.8063 & \$87,182 & 0.3071 & 22.59 & 0.1343 & 0.2518 \\
\bottomrule
\end{tabular}
\vspace{1mm}
\begin{minipage}{\textwidth}
\scriptsize
\emph{Notes.}
Configuration is a fitted penalty value for the Surrogate rows and an exactly solved
base value for the post-hoc rows. The ``Max attained'' column is filled only for rows recorded as
not attained, and reports the largest development first-order correction that family reached over
the evaluated frozen path. Over that path the Direct family did not attain any of the three extended
targets and the Surrogate family did not attain zero; both post-hoc families attain all three
exactly, because a one-parameter rescaling can be solved for any first-order target. Reaching a
target is not by itself an advantage: driving the post-hoc coordinate to neutrality costs the
price-scale accuracy recorded in Table~\ref{tab:centered_spread}. COD is in percent;
$\Delta_{\mathrm{NL}}$ and dCor carry no assessor-facing reference range.
\end{minipage}
\end{table}
}


\subsection{Retransformation Sensitivity}
\label{app:smearing}

Predictions are returned to the price scale by direct exponentiation,
Equation~\eqref{eq:price_prediction}. That is an operational retransformation and not a
conditional-mean correction, so it is a convention rather than a result. A frozen sensitivity
therefore re-evaluates every display configuration under a Duan smearing factor
\[
s=\frac{\sum_i w_i\exp(u_i)}{\sum_i w_i},
\qquad u_i=\log P_i-f(x_i),
\]
applied as $f(x_i)\mapsto f(x_i)+\log s$. The residual convention here is the negative of the one
used by the residual-dependence diagnostics, where $e_i=f(x_i)-\log P_i$; the two are not
interchangeable. The factor is estimated from development out-of-fold residuals under the
one-sale-one-vote weighting and is never re-estimated on an evaluation block.

The consequences separate by construction. Valuation-level ratios---median, mean and
weighted-mean---are multiplied exactly by $s$. COD, COV, PRD, PRB, MKI, VEI, $\beta_{\log}$,
$\Delta_{\mathrm{NL}}$ and $\operatorname{dCor}(e,y)$ are invariant to a common multiplicative
factor as a matter of algebra, so no vertical-equity, uniformity, or residual-structure result in
this paper depends on the retransformation choice. The frozen sensitivity corroborates that
invariance numerically on the configurations it checks: exhaustively for the level, uniformity,
vertical-equity and distance-correlation measures, and on a smaller checked subset for
$\Delta_{\mathrm{NL}}$, whose guarantee rests on the algebra rather than on that subset.
$R^2_P$, MAE and MAPE are not invariant and move nontrivially.
$\operatorname{RMSE}_{\log P}$ is defined on the log scale and is not a price-scale rescaling
quantity, so the frozen design does not recompute it under smearing.

Table~\ref{tab:smearing_matched_beta} reports the movement at the three matched
first-order configurations of Table~\ref{tab:matched_beta_headline}. The factor is not common
across the arms---it is largest for the Surrogate and smallest for the post-hoc comparator---so
direct exponentiation does not affect them equally. Under the sensitivity the price-scale
ordering of the Surrogate and the post-hoc comparator reverses on both $R^2_P$ and MAE. The
price-scale half of that matched comparison is therefore convention-dependent and is reported as
such rather than as a property of the mechanisms; the residual-structure half is not, since the
$\Delta_{\mathrm{NL}}$ and dCor entries of Table~\ref{tab:matched_beta_headline} are unchanged. The
frozen sensitivity was evaluated over the whole frozen configuration set rather than this table
alone; what is reported here is the matched-target illustration, because the purpose is to show that
the price-scale half of that comparison is convention-dependent, not to introduce a second ranking.

\begin{table}[!ht]
\centering
\caption{Retransformation sensitivity at the matched first-order targets, development
coordinate D1 at the strongest common target, held-out evaluation. Baseline columns reproduce
Table~\ref{tab:matched_beta_headline}; smeared columns apply the frozen Duan factor $s$ to the
same stored predictions. Nothing is refitted, no configuration is selected or recommended, and
the table is a convention check rather than a result.}
\label{tab:smearing_matched_beta}
\begin{tabular}{@{}lrrrrr@{}}
\toprule
Family & $s$ & $R^2_P$ base & $R^2_P$ smeared & MAE base & MAE smeared \\
\midrule
Direct & 1.0636 & 0.8760 & 0.8498 & \$82,764 & \$87,979 \\
Surrogate & 1.0879 & 0.8963 & 0.9010 & \$76,128 & \$73,908 \\
C-posthoc & 1.0524 & 0.8966 & 0.8860 & \$74,236 & \$75,548 \\
\bottomrule
\end{tabular}
\end{table}

One asymmetry is recorded rather than resolved. The factor is estimated on development
out-of-fold residuals, whereas the 2025 forward block is evaluated from a production refit that
has no out-of-fold analogue, so the development-estimated factor is applied there unchanged.
That is an explicit assumption of the sensitivity, not a validated property of the forward
period.


\subsection{VEI Percentile-Group Profiles}

\begin{figure}[!htbp]
\centering
\safeincludegraphics[width=0.8\textwidth]{img/generated_v12_994/vei_percentile_group_profile.pdf}
\caption{VEI percentile-group median valuation ratios for ordinary LightGBM, drawn beside the historical township-level linear-regression trace, on the held-out and 2025 samples. The dashed line marks a valuation ratio of $1$; the dotted line marks the overall sample median ratio. The shaded ribbons are descriptive, deterministic percentile-bootstrap intervals for the plotted group medians at a nominal 90\% level, drawn only to convey sampling variability in this visualization. They are \emph{not} the rank-based order-statistic median interval that the May-2026 IAAO exposure draft specifies for its own VEI procedure in Appendix~D.2, and they are \emph{not} a standards-facing bootstrap-sensitivity calculation. No exposure-draft VEI verdict is attached to this figure. The historical township-level linear-regression trace is drawn for visual context only: no linear-regression quantitative result is reported in this paper, and none should be read off this figure.}
\label{fig:vei_group_profile_placeholder}
\end{figure}

\section{Reproducibility and Archived Artifacts}
\label{app:ccao_results}

This appendix states what the repository archives, what a reader can therefore check, and what
cannot be reproduced from public material. A paper-specific guide, \path{paper/REPRODUCIBILITY.md},
carries the path-level detail, the machine-readable inventory and the exact commands; only the
substance is summarized here. Three levels are distinguished, and they are deliberately not
collapsed into a single claim of reproducibility.

\paragraph{Auditability of the reported results.} Every number printed in this paper can be
re-derived without refitting anything. The manuscript source, the figures it includes, the frozen
specifications and reference conventions, the provenance records, and a numeric ledger are tracked
in the repository. The ledger binds each printed value to a tracked artifact by content hash,
together with the selector that locates it, the transform applied to it and the rounding rule used
to print it; a validator reopens each artifact, verifies its hash against both the file and the
frozen manifest, re-executes the selector and compares the result with what the source prints. The
executed metric code, including the ratio-study measures, the nonlinearity-gap estimator and the
distance-correlation call site, is tracked and is byte-identical at every provenance commit.

\paragraph{Reconstruction from frozen derived artifacts.} The derived evidence behind the reported
comparisons is tracked and can be re-examined without the raw extract. This covers the matched
first-order design and its outcomes, with the extended targets, pairwise differences, crossing
locations, origin-identity check and the one-sale-one-vote sensitivity twins; the post-hoc
comparator path over all frozen base values, with its per-fold summaries, closed-form linearity
check, root verification and centering sensitivity; the zero-penalty control and the parity
ladders; the retransformation-sensitivity, PRB and Exposure-Draft VEI inference tables; and the
distance-correlation implementation record. What this level does not include is the
per-configuration fold, held-out and forward prediction arrays, which are record-level outputs
derived from restricted assessor data and are not released.

\paragraph{End-to-end rerun.} A full rerun of the empirical pipeline is not publicly reproducible.
It requires the Cook County extract identified in Section~\ref{subsec:ccao_design}, which is
hash-pinned but not redistributed and for which no public availability is claimed, together with
the executed run roots holding the stored predictions and the combined path table. The frozen
evidence manifest records each of those inputs by path, byte length and content hash, so the reader
can confirm which file was used; it does not put the file itself in the reader's hands. A reader
without authorized access to that exact input can audit and reconstruct the reported results, but
cannot refit the models.

\paragraph{Two limitations belong with this account.} The recorded environment is the one the
executed pipeline observed and not a lockfile: the interpreter and the seven libraries that matter
for the results are pinned to exact versions in two independent frozen records, while the
repository requirement file states only lower bounds and does not reconstruct that environment. And
the source state that generated the frozen artifacts is bounded rather than recovered, for the
reason given in Appendix~\ref{app:implementation}: the generating trees were dirty and the
uncommitted portion was never archived.

Nothing in this appendix is a statement about operational readiness. Before operational use, the
same analysis should report block-bootstrap intervals and ratio curves by township, property class,
value range, and time, so that model-selection uncertainty, implementation effects and local
variation are separated. This paper reports none of those.


\section*{Data, Code, and Institutional Statement}

The primary application uses Cook County Assessor's Office data together with a research
translation of the office's published residential modeling workflow. The analysis code, the frozen
specifications and conventions, the derived result tables behind every reported comparison, the
provenance records and the numeric ledger are tracked in the public repository, and
\path{paper/REPRODUCIBILITY.md} is the entry point for auditing this paper specifically; the
repository's general development README describes a broader legacy workflow and is not the protocol
behind these results. The raw CCAO extract is not redistributed. It is identified by content hash
and byte length rather than by a version string, and no public availability is claimed for it, so an
external reader can verify which input was used but cannot obtain it here. Consequently the reported
results can be audited and reconstructed from tracked artifacts, while refitting the models requires
authorized access to that exact input. The per-record prediction artifacts are not released. The
results reported here are research findings from a research translation of the CCAO workflow, not
official CCAO assessment decisions, and the paper makes no claim that CCAO has adopted the
correction or placed it into production.


\end{document}
